% Leçon 11 — Expression orale : importance de la pratique du sport — Chinois 1ère A — Cours signé OnBuch+
% Programme officiel MINESEC. Compiler avec : tectonic ZHA11-oral-importance-de-la-pratique-du-sport.tex
\newcommand{\DOCMATIERE}{Chinois}
\newcommand{\DOCNIVEAU}{1ère A}
\newcommand{\DOCMODULE}{Module 2 : Environnement et santé}
\newcommand{\DOCLECON}{ZHA11}
\newcommand{\DOCTITRE}{Leçon 11 — Expression orale : importance de la pratique du sport}
% OnBuch+ — préambule « cours signés » ENRICHI (compilé avec tectonic / XeLaTeX)
% Système visuel « Pop » d'OnBuch+ : fond crème (jamais blanc), bordures encre
% épaisses, ombres « dures » décalées, titres Archivo Black en capitales.
% Version MATHÉMATIQUES : graphes (pgfplots/tikz), boîtes pédagogiques
% (définition, propriété, méthode, attention, le savais-tu, exercice résolu,
% exercice, TP, bilan), page de garde et signature OnBuch+ ancrée en bas.
%
% Macros attendues dans main.tex AVANT \input{preamble} :
%   \DOCMATIERE  \DOCNIVEAU  \DOCMODULE  \DOCLECON (numéro)  \DOCTITRE
\documentclass[11pt]{article}

\usepackage{fontspec}
\usepackage[french]{babel} % français = langue principale ; le chinois via \zh{..} / textebox (xeCJK)
\usepackage{xeCJK}
\usepackage{pifont} % \ding{51} \ding{55} utilisés par les modèles
\usepackage[a4paper,margin=2cm,top=3.4cm,bottom=2.8cm,headheight=1.4cm,footskip=1.3cm]{geometry}
\usepackage{amsmath,amssymb,amsfonts}
\usepackage{mathtools} % \xmapsto, \coloneqq... souvent utilisés par les modèles
\usepackage{cancel} % simplifications barrées (bilans de forces)
% \abs{z} (module, valeur absolue) et \norm{u} (norme), utilisés par les modèles
\providecommand{\abs}{}\let\abs\relax\DeclarePairedDelimiter{\abs}{\lvert}{\rvert}
\providecommand{\norm}{}\let\norm\relax\DeclarePairedDelimiter{\norm}{\lVert}{\rVert}
\usepackage[table]{xcolor} % [table] : fournit aussi \rowcolors (alternance de lignes)
\usepackage{pagecolor}
\usepackage{tikz}
\usetikzlibrary{arrows.meta,positioning,calc,shapes.geometric,shapes.misc,shapes.arrows,decorations.pathmorphing,decorations.pathreplacing,decorations.markings,patterns,shadows,mindmap,trees,fit,backgrounds,matrix,intersections,angles,quotes}
\usepackage{pgfplots}
\usepackage[version=4]{mhchem} % équations nucléaires \ce{^{235}_{92}U}
\usepackage[european,siunitx]{circuitikz} % schémas électriques
\pgfplotsset{compat=1.18}
\usepgfplotslibrary{fillbetween,groupplots} % aires entre courbes (calcul intégral)
\usepackage{enumitem}
\usepackage{fancyhdr}
% [most] charge toutes les bibliothèques tcolorbox (listings, minted, raster,
% documentation...) et gonfle inutilement la mémoire TeX sur un document long
% (~300 boîtes) : on ne charge que ce qui est réellement utilisé.
\usepackage[skins,breakable]{tcolorbox}
\usepackage{graphicx}
\usepackage{colortbl}
\usepackage{booktabs}
\usepackage{tabularx}
\usepackage{diagbox} % case d'en-tête diagonale des tableaux à double entrée
% Colonnes extensibles centrée / gauche / droite (C, L, R), souvent utilisées par les modèles
\newcolumntype{C}{>{\centering\arraybackslash}X}
\newcolumntype{L}{>{\raggedright\arraybackslash}X}
\newcolumntype{R}{>{\raggedleft\arraybackslash}X}
\usepackage{array}
\usepackage{multicol}
\usepackage{multirow}
\usepackage{siunitx}
% parse-numbers=false : accepte aussi \SI{2,5 \times 10^{-3}}{...}
\sisetup{locale=FR,output-decimal-marker={,},group-minimum-digits=4,parse-numbers=false}
\DeclareSIUnit\ohm{\ensuremath{\symup{\Omega}}} % U+2126 absent de la police
\DeclareSIUnit\division{div} % réglages d'oscilloscope (ms/div, V/div)
\DeclareSIUnit\tr{tr} % tours (vitesse de rotation, ex. \SI{1500}{\tr\per\minute})
\DeclareSIUnit\molar{M} % concentration molaire (ex. \SI{3}{\molar}, \SI{1}{\micro\molar})
\DeclareSIUnit\an{an} % année (le modèle confond parfois avec \year, un compteur TeX) — ex. \SI{5}{\kilo\meter\per\an}
\DeclareSIUnit\FCFA{FCFA} % franc CFA (ex. \SI{500}{\FCFA\per\kilo\gram})
\DeclareSIUnit\degreeBrix{\textdegree Brix} % teneur en sucre d'un sirop/jus (réfractométrie)
\DeclareSIUnit\month{mois} % le modèle utilise parfois \month, un compteur TeX, comme unité de durée
\usepackage{qrcode}
% \degree : commande fréquemment utilisée par les modèles pour un angle ou une
% température en dehors de \SI{}{} (non fournie par les paquets chargés).
\providecommand{\degree}{^{\circ}}
% \qed (fin de démonstration), utilisé par les modèles sans amsthm
\providecommand{\qed}{\hfill$\square$}
% \barre : double barre des valeurs interdites dans un tableau de variations (tkz-tab)
\providecommand{\barre}{\ensuremath{\|}}
% environnement proof (amsthm non chargé) utilisé par les modèles
\ifdefined\proof\else\newenvironment{proof}[1][Démonstration]{\par\noindent\textit{#1.}\ }{\qed\par}\fi
\usepackage{lastpage}
\usepackage{hyperref}
\hypersetup{hidelinks}

% ============================================================
% Palette « Pop » OnBuch+ — mêmes hex que la classe P (plus_ui.dart)
% ============================================================
\definecolor{poporange}{HTML}{F59321}
\definecolor{popdark}{HTML}{E07A0C}
\definecolor{popcream}{HTML}{FFF6E8}
\definecolor{popink}{HTML}{1C1714}
\definecolor{poppurple}{HTML}{7A5AE0}
\definecolor{popgreen}{HTML}{2E9E6A}
\definecolor{poppink}{HTML}{E0457B}
\definecolor{popblue}{HTML}{2F7DE1}
\definecolor{popgold}{HTML}{C98A12}
\definecolor{popmuted}{HTML}{8A7F76}
\definecolor{popline}{HTML}{EADFD2}
\definecolor{popgray}{HTML}{8A7F76}
\colorlet{popgrey}{popgray}
% Alias tolérants (le modèle nomme parfois les couleurs différemment)
\colorlet{popwhite}{white}
\colorlet{popblack}{popink}
\colorlet{popred}{poppink}
\colorlet{popyellow}{popgold}
% Teintes claires (fonds de tableaux/encadrés) — le modèle nomme parfois
% les couleurs avec un suffixe L (ex. popblueL) sans les avoir définies.
\colorlet{poporangeL}{poporange!20}
\colorlet{popdarkL}{popdark!20}
\colorlet{poppurpleL}{poppurple!20}
\colorlet{popgreenL}{popgreen!20}
\colorlet{poppinkL}{poppink!20}
\colorlet{popblueL}{popblue!20}
\colorlet{popgoldL}{popgold!20}
\colorlet{popredL}{popred!20}
\colorlet{popgrayL}{popgray!20}
% Teintes claires pour fonds de tableaux / zones
\colorlet{popblueL}{popblue!10!white}
\colorlet{poporangeL}{poporange!14!white}
\colorlet{popgreenL}{popgreen!12!white}
\colorlet{poppurpleL}{poppurple!12!white}
\colorlet{poppinkL}{poppink!12!white}
\colorlet{popgoldL}{popgold!14!white}

\pagecolor{popcream}

% ============================================================
% Typographie
% ============================================================
\defaultfontfeatures{Ligatures=TeX}
\setmainfont{PlusJakartaSans-Medium.ttf}[Path=../../../fonts/,BoldFont=PlusJakartaSans-Bold.ttf]
% Symboles, API et emojis absents de la police principale : repli sur DejaVu Sans (caractères actifs)
\newfontface\symfont{DejaVuSans.ttf}[Path=../../../fonts/]
\makeatletter
\newcommand\symactive[1]{\catcode#1=13 \begingroup\lccode`\~=#1 \lowercase{\endgroup\def~}{{\symfont\char#1}}}
\newcommand\symblank[1]{\catcode#1=13 \begingroup\lccode`\~=#1 \lowercase{\endgroup\def~}{}}
\newcommand\symhyph[1]{\catcode#1=13 \begingroup\lccode`\~=#1 \lowercase{\endgroup\def~}{\mbox{-}}}
\newcount\symc
\newcommand\symrange[2]{\symc=#1 \@whilenum\symc<#2 \do{\expandafter\symactive\expandafter{\the\symc}\advance\symc by 1 }}
\newcommand\symblankrange[2]{\symc=#1 \@whilenum\symc<#2 \do{\expandafter\symblank\expandafter{\the\symc}\advance\symc by 1 }}
\makeatother
\symrange{"0250}{"02FF}\symactive{"2713}\symactive{"2714}\symactive{"2717}\symactive{"2718}\symactive{"2610}\symactive{"2611}\symactive{"2612}
\symrange{"2600}{"27BF}
\catcode"2705=13 \begingroup\lccode`\~="2705 \lowercase{\endgroup\def~}{{\symfont\char"2713}}\symblank{"26BD}\symblank{"26A1}
\symrange{"2460}{"257F}\symactive{"223C}\symactive{"2248}\symactive{"2260}
\symhyph{"2011}\symhyph{"2E1A}\symblankrange{"1F300}{"1FAFF}\symblank{"FE0F}
% Caractères chinois : WenQuanYi Zen Hei (sous-ensemble GB2312 + pinyin), gras/italique simulés
\setCJKmainfont{WenQuanYiZenHei-subset.ttf}[Path=../../../fonts/,AutoFakeBold=3,AutoFakeSlant=0.2]
\xeCJKsetup{CJKspace=false}
% Pinyin : voyelles à caron (ǎ ǐ ǒ ǔ ǖ ǘ ǚ ǜ) absentes de la police principale -> caractères actifs
\newfontface\pyfont{WenQuanYiZenHei-subset.ttf}[Path=../../../fonts/]
\newcommand\pyactive[1]{\catcode#1=13 \begingroup\lccode`\~=#1 \lowercase{\endgroup\def~}{{\pyfont\char#1}}}
\pyactive{"01CD}\pyactive{"01CE}\pyactive{"01CF}\pyactive{"01D0}\pyactive{"01D1}\pyactive{"01D2}\pyactive{"01D3}\pyactive{"01D4}\pyactive{"01D5}\pyactive{"01D6}\pyactive{"01D7}\pyactive{"01D8}\pyactive{"01D9}\pyactive{"01DA}\pyactive{"01DB}\pyactive{"01DC}
% Maths en sans-serif (Fira Math) avec chiffres et lettres droites de Plus
% Jakarta, pour que les formules se fondent dans le texte.
\usepackage[math-style=ISO,bold-style=ISO]{unicode-math}
\setmathfont{FiraMath-Regular.otf}[Path=../../../fonts/]
\setmathfont{PlusJakartaSans-Medium.ttf}[Path=../../../fonts/,range={up/num,up/{Latin,latin},\mathup}]
\usepackage{icomma} % virgule décimale sans espace en mode math
% Caractères mathématiques Unicode absents des polices de texte (Plus Jakarta,
% Archivo Black) : rendus par leur équivalent mathématique, en texte comme en
% formule — sinon ils s'affichent en carré vide (ex. « Arithmétique dans ℤ »).
\usepackage{newunicodechar}
\newunicodechar{ℝ}{\ensuremath{\mathbb{R}}}
\newunicodechar{ℤ}{\ensuremath{\mathbb{Z}}}
\newunicodechar{ℂ}{\ensuremath{\mathbb{C}}}
\newunicodechar{ℕ}{\ensuremath{\mathbb{N}}}
\newunicodechar{ℚ}{\ensuremath{\mathbb{Q}}}
\newunicodechar{𝔻}{\ensuremath{\mathbb{D}}}
\newunicodechar{α}{\ensuremath{\alpha}}
\newunicodechar{β}{\ensuremath{\beta}}
\newunicodechar{γ}{\ensuremath{\gamma}}
\newunicodechar{δ}{\ensuremath{\delta}}
\newunicodechar{ε}{\ensuremath{\varepsilon}}
\newunicodechar{θ}{\ensuremath{\theta}}
\newunicodechar{λ}{\ensuremath{\lambda}}
\newunicodechar{μ}{\ensuremath{\mu}}
\newunicodechar{σ}{\ensuremath{\sigma}}
\newunicodechar{τ}{\ensuremath{\tau}}
\newunicodechar{φ}{\ensuremath{\varphi}}
\newunicodechar{ψ}{\ensuremath{\psi}}
\newunicodechar{ω}{\ensuremath{\omega}}
\newunicodechar{Σ}{\ensuremath{\Sigma}}
% majuscules produites par \MakeUppercase dans les titres (α-aminés -> Α)
\newunicodechar{Α}{\ensuremath{\alpha}}
\newunicodechar{Β}{\ensuremath{\beta}}
\newunicodechar{Γ}{\ensuremath{\Gamma}}
\newunicodechar{Δ}{\ensuremath{\Delta}}
\newunicodechar{Ε}{\ensuremath{\varepsilon}}
\newunicodechar{Θ}{\ensuremath{\Theta}}
\newunicodechar{Λ}{\ensuremath{\Lambda}}
\newunicodechar{Μ}{\ensuremath{\mu}}
\newunicodechar{Τ}{\ensuremath{\tau}}
\newunicodechar{Φ}{\ensuremath{\Phi}}
\newunicodechar{Ψ}{\ensuremath{\Psi}}
\newunicodechar{Ω}{\ensuremath{\Omega}}
\newunicodechar{↦}{\ensuremath{\mapsto}}
\newunicodechar{∈}{\ensuremath{\in}}
\newunicodechar{∉}{\ensuremath{\notin}}
\newunicodechar{∧}{\ensuremath{\wedge}}
\newunicodechar{⇔}{\ensuremath{\Leftrightarrow}}
\newunicodechar{⟺}{\ensuremath{\Longleftrightarrow}}
\newunicodechar{⇒}{\ensuremath{\Rightarrow}}
\newunicodechar{⟹}{\ensuremath{\Longrightarrow}}
\newunicodechar{∘}{\ensuremath{\circ}}
\newunicodechar{∪}{\ensuremath{\cup}}
\newunicodechar{∩}{\ensuremath{\cap}}
\newunicodechar{≡}{\ensuremath{\equiv}}
\newunicodechar{⊂}{\ensuremath{\subset}}
\newunicodechar{⊥}{\ensuremath{\perp}}
\newunicodechar{∃}{\ensuremath{\exists}}
\newunicodechar{∀}{\ensuremath{\forall}}
\newunicodechar{∛}{\ensuremath{\sqrt[3]{\,}}}
\newunicodechar{∜}{\ensuremath{\sqrt[4]{\,}}}
\newunicodechar{⃗}{\ensuremath{{}^{\to}}}
\newunicodechar{ⁿ}{\ifmmode^{n}\else\textsuperscript{\ensuremath{n}}\fi}
\newunicodechar{ˣ}{\ifmmode^{x}\else\textsuperscript{\ensuremath{x}}\fi}
\newunicodechar{ᵇ}{\ifmmode^{b}\else\textsuperscript{\ensuremath{b}}\fi}
\newunicodechar{ʳ}{\ifmmode^{r}\else\textsuperscript{\ensuremath{r}}\fi}
\newunicodechar{ᵘ}{\ifmmode^{u}\else\textsuperscript{\ensuremath{u}}\fi}
\newunicodechar{ʸ}{\ifmmode^{y}\else\textsuperscript{\ensuremath{y}}\fi}
\newunicodechar{ᵃ}{\ifmmode^{a}\else\textsuperscript{\ensuremath{a}}\fi}
\newunicodechar{ᵉ}{\ifmmode^{e}\else\textsuperscript{\ensuremath{e}}\fi}
\newunicodechar{ᵏ}{\ifmmode^{k}\else\textsuperscript{\ensuremath{k}}\fi}
\newunicodechar{ᵖ}{\ifmmode^{p}\else\textsuperscript{\ensuremath{p}}\fi}
\newunicodechar{ᴺ}{\ifmmode^{N}\else\textsuperscript{\ensuremath{N}}\fi}
\newunicodechar{ᶜ}{\ifmmode^{c}\else\textsuperscript{\ensuremath{c}}\fi}
\newunicodechar{ʰ}{\ifmmode^{h}\else\textsuperscript{\ensuremath{h}}\fi}
\newunicodechar{ᵠ}{\ifmmode^{q}\else\textsuperscript{\ensuremath{q}}\fi}
\newunicodechar{ₙ}{\ifmmode_{n}\else\textsubscript{\ensuremath{n}}\fi}
\newunicodechar{ₐ}{\ifmmode_{a}\else\textsubscript{\ensuremath{a}}\fi}
\newunicodechar{ᵢ}{\ifmmode_{i}\else\textsubscript{\ensuremath{i}}\fi}
\newunicodechar{ᵣ}{\ifmmode_{r}\else\textsubscript{\ensuremath{r}}\fi}
\newunicodechar{ₖ}{\ifmmode_{k}\else\textsubscript{\ensuremath{k}}\fi}
\newunicodechar{ᵦ}{\ifmmode_{\beta}\else\textsubscript{\ensuremath{\beta}}\fi}
\newunicodechar{ₓ}{\ifmmode_{x}\else\textsubscript{\ensuremath{x}}\fi}
\newfontfamily\popdisplay{ArchivoBlack-Regular.ttf}[Path=../../../fonts/]
\newfontfamily\poplogo{SpaceGrotesk-Bold.ttf}[Path=../../../fonts/]
\newfontfamily\popsemi{PlusJakartaSans-SemiBold.ttf}[Path=../../../fonts/]
\newfontfamily\popxbold{PlusJakartaSans-ExtraBold.ttf}[Path=../../../fonts/]
\newfontfamily\popxboldls{PlusJakartaSans-ExtraBold.ttf}[Path=../../../fonts/,LetterSpace=3.0]

\setlist[enumerate]{leftmargin=1.7em,itemsep=5pt,topsep=4pt}
\setlist[itemize]{leftmargin=1.7em,itemsep=4pt,topsep=4pt,label={\color{poporange}\textbullet}}
\setlength{\parindent}{0pt}
\setlength{\parskip}{5pt}
\renewcommand{\arraystretch}{1.25}

% pgfplots : style Pop par défaut
\pgfplotsset{
  every axis/.append style={axis line style={popink,line width=1pt},tick style={popink},
    grid style={popline,line width=0.5pt},label style={font=\small},tick label style={font=\footnotesize}},
  popcurve/.style={poporange,line width=1.8pt,no markers,smooth},
}
% Même style disponible dans un \draw TikZ simple (hors pgfplots)
\tikzset{popcurve/.style={poporange,line width=1.8pt}}
\tikzset{popgrid/.style={popline,very thin}}
\colorlet{popcurve}{poporange} % le nom du style est parfois utilisé comme couleur

% ============================================================
% Pill / squiggle
% ============================================================
\newcommand{\poppill}[3]{%
  \tikz[baseline=-0.55ex]\node[draw=popink,line width=1.2pt,rounded corners=2.6pt,%
    fill=#2,inner xsep=6.5pt,inner ysep=3pt]{%
    \popxboldls\fontsize{7}{7}\selectfont\color{#3}\MakeUppercase{#1}};%
}
\newcommand{\popsquiggle}[1]{%
  \tikz[baseline=-0.4ex]{
    \draw[#1,line width=2.6pt,line cap=round]
      plot [smooth] coordinates {(0,0.09) (0.34,0.01) (0.68,0.21) (1.02,0.01) (1.36,0.10)};
  }%
}
% Mot-clé surligné (marqueur orange sous le texte)
\newcommand{\cle}[1]{{\popxbold\color{popdark}#1}}

% ============================================================
% En-tête / pied
% ============================================================
\pagestyle{fancy}
\fancyhf{}
\renewcommand{\headrulewidth}{0pt}
\renewcommand{\footrulewidth}{0pt}
\lhead{\raisebox{-0.15ex}{\poplogo\fontsize{13}{13}\selectfont\color{popink}OnBuch\color{poporange}+}}
\rhead{\poppill{\DOCMATIERE\ \textbullet\ \DOCNIVEAU\ \textbullet\ Leçon \DOCLECON}{white}{popink}}
\lfoot{\popxbold\fontsize{7.6}{7.6}\selectfont\color{popmuted}\MakeUppercase{Cours signé OnBuch+}\ \textbullet\ {\popsemi\DOCTITRE}}
\rfoot{\tikz[baseline=-0.6ex]\node[rounded corners=8.5pt,fill=popink,minimum height=17pt,minimum width=17pt,inner xsep=5pt,inner sep=0]{\popxbold\fontsize{8}{8}\selectfont\color{popcream}\thepage\,/\,\pageref*{LastPage}};}

% ============================================================
% Titres
% ============================================================
\newcommand{\chaptitre}[2]{%
  \begin{tcolorbox}[enhanced,colback=poporange,colframe=popink,boxrule=2.6pt,arc=11pt,
      drop shadow=popink,left=15pt,right=15pt,top=12pt,bottom=11pt,before skip=3pt,after skip=18pt]
    \poppill{#2}{popink}{popcream}\par\vspace{9pt}
    {\raggedright\hyphenpenalty=10000\exhyphenpenalty=10000\popdisplay\fontsize{22}{24}\selectfont\color{popcream}\MakeUppercase{#1}\par}
  \end{tcolorbox}
}
% Section numérotée : pastille orange avec le numéro + titre Archivo
\newcounter{popsec}
\newcommand{\coursec}[1]{%
  \stepcounter{popsec}\setcounter{popsub}{0}%
  \par\vspace{16pt}\noindent
  \begin{minipage}{\linewidth}
  \tikz[baseline=-0.7ex]\node[draw=popink,line width=1.6pt,fill=poporange,rounded corners=4pt,minimum size=22pt,inner sep=0]{\popdisplay\fontsize{12}{12}\selectfont\color{popcream}\thepopsec};%
  \hspace{8pt}{\popdisplay\fontsize{15}{17}\selectfont\color{popink}\MakeUppercase{#1}}\par
  \vspace{4pt}\noindent{\color{poporange}\rule{36pt}{3.6pt}}
  \end{minipage}\par\nopagebreak\vspace{8pt}
}
\newcounter{popsub}
\newcommand{\courssub}[1]{%
  \stepcounter{popsub}%
  \par\vspace{9pt}\noindent{\popxbold\fontsize{11.5}{13}\selectfont\color{popink}{\color{poporange}\thepopsec.\thepopsub}\hspace{6pt}#1}\par\nopagebreak\vspace{4pt}
}

% ============================================================
% Boîtes pédagogiques — même recette : fond blanc, bordure épaisse colorée,
% ombre dure encre, badge plein. Titre optionnel [#1].
% ============================================================
\tcbset{popbase/.style={breakable,enhanced,colback=white,boxrule=2.3pt,arc=9pt,
  shadow={3pt}{-3pt}{0pt}{popink},left=14pt,right=14pt,top=11pt,bottom=11pt,before skip=11pt,after skip=15pt,
  fontupper=\popsemi\fontsize{10.6}{14.5}\selectfont\color{popink},
  fontlower=\popsemi\fontsize{10.6}{14.5}\selectfont\color{popink}}}
% Coche dessinée (le glyphe ✓ n'existe pas dans les polices du document)
\newcommand{\popcheck}[1]{\tikz[baseline=-0.5ex]\draw[#1,line width=1.6pt,line cap=round,line join=round](-0.13,0.0)--(-0.04,-0.09)--(0.13,0.1);}
\providecommand{\coche}{\popcheck{popgreen}}
\providecommand{\resultat}[1]{\par\smallskip\noindent\fcolorbox{popgreen}{popgreenL}{\parbox{\dimexpr\linewidth-2\fboxsep-2\fboxrule\relax}{\textbf{Résultat.} #1}}\par\smallskip}
\newcommand{\popboxhead}[3]{\poppill{#1}{#2}{white}\ifx\relax#3\relax\else\hspace{6pt}{\popxbold\fontsize{10.6}{12}\selectfont\color{popink}#3}\fi\par\vspace{7pt}}

\newtcolorbox{aretenir}[1][]{popbase,colframe=popgreen,before upper={\popboxhead{À retenir}{popgreen}{#1}}}
% Texte en chinois (document de lecture, dialogue, extrait) : cadre
% d'étude. Un titre optionnel (source, auteur) s'affiche dans l'en-tête.
\newtcolorbox{textebox}[1][]{popbase,colframe=popblue,before upper={\popboxhead{文本 · Texte}{popblue}{#1}},after upper={}}
% Marques ✓ / ✗ (forme correcte / fautive) souvent utilisées par les modèles
\providecommand{\cross}{\textcolor{poppink}{\ensuremath{\times}}}
\providecommand{\xmark}{\cross}\providecommand{\crossmark}{\cross}\providecommand{\cmark}{\ensuremath{\checkmark}}
% Ligne de réponse / justification à compléter (inventée par les modèles)
\providecommand{\justification}[1]{\par\noindent\textit{Justification~:} \dotfill\par}
% \textipa (package tipa non chargé) : la police ne couvre pas l'API -> simple passage
\providecommand{\textipa}[1]{#1}
% Soulignement (ulem non chargé) : \uline, \ul
\providecommand{\uline}[1]{\underline{#1}}\providecommand{\ul}[1]{\underline{#1}}
% \glossaire{\item ...} inventée par les modèles : liste de glossaire
\providecommand{\glossaire}[1]{\begin{itemize}#1\end{itemize}}
% \alert (beamer) inventée par les modèles : texte mis en évidence
\providecommand{\alert}[1]{\textbf{\textcolor{poppink}{#1}}}
% variantes ASCII d'ordinaux écrites en macro
\providecommand{\iere}{\textsuperscript{re}}\providecommand{\ieme}{\textsuperscript{e}}\providecommand{\ere}{\textsuperscript{re}}\providecommand{\eme}{\textsuperscript{e}}\providecommand{\er}{\textsuperscript{er}}\providecommand{\ie}{\textsuperscript{e}}
% Ordinaux français écrits en macro par les modèles : 1\ière, 3\ième
\newcommand{\ière}{\textsuperscript{re}}\newcommand{\ième}{\textsuperscript{e}}
% \exemple (inventée par les modèles) : étiquette « Exemple : »
\providecommand{\exemple}{\textbf{Exemple~:}\ }
% \square utilisable aussi hors mode math (cases à cocher)
\AtBeginDocument{\let\squareMath\square\renewcommand{\square}{\ensuremath{\squareMath}}}
% Mot ou phrase en chinois à l'intérieur d'un paragraphe français
\newcommand{\zh}[1]{\ifmmode\text{#1}\else{#1}\fi}
\let\eng\zh \let\esp\zh \let\all\zh % alias tolérés
% flèches d'intonation inventées par les modèles
\providecommand{\textsearrow}{}\renewcommand{\textsearrow}{\ensuremath{\searrow}}\providecommand{\textnearrow}{}\renewcommand{\textnearrow}{\ensuremath{\nearrow}}
% \fr{..} : mot français (inventé par les modèles)
\providecommand{\fr}[1]{\textit{#1}}
% zhbox : encadré de texte chinois inventé par les modèles
\newenvironment{zhbox}{\begin{quote}}{\end{quote}}
% \vs : « versus » inventé par les modèles
\providecommand{\vs}{\textit{vs}\ }
% \blank : case à compléter inventée par les modèles
\providecommand{\blank}[1][]{\underline{\hspace{1.6em}}}
\newtcolorbox{exemplebox}[1][]{popbase,colframe=poporange,before upper={\popboxhead{Exemple}{poporange}{#1}}}
\newtcolorbox{definition}[1][]{popbase,colframe=popblue,before upper={\popboxhead{Définition}{popblue}{#1}}}
\newtcolorbox{propriete}[1][]{popbase,colframe=poppurple,before upper={\popboxhead{Propriété}{poppurple}{#1}}}
\newtcolorbox{methode}[1][]{popbase,colframe=popgold,before upper={\popboxhead{Méthode}{popgold}{#1}}}
\newtcolorbox{attention}[1][]{popbase,colframe=poppink,before upper={\popboxhead{Attention}{poppink}{#1}}}
\newtcolorbox{experience}[1][]{popbase,colframe=popblue,colback=popblueL,before upper={\popboxhead{Expérience}{popblue}{#1}}}
% « Le savais-tu ? » : carte violette pleine (contexte, culture, Cameroun)
\newtcolorbox{savaistu}[1][]{popbase,colback=poppurple,colframe=popink,
  fontupper=\popsemi\fontsize{10.4}{14.2}\selectfont\color{white},
  before upper={\poppill{Le savais-tu\,?}{white}{poppurple}\ifx\relax#1\relax\else\hspace{6pt}{\popxbold\color{white}#1}\fi\par\vspace{7pt}}}
% Exercice résolu : énoncé en haut, \tcblower puis la solution sur fond crème
\newcounter{popexo}\newcounter{popexr}
\newtcolorbox{exoresolu}[1][]{popbase,colframe=poporange,colbacklower=popcream,
  segmentation style={popink,line width=1.2pt,dashed},
  before upper={\stepcounter{popexr}\popboxhead{Exercice résolu \thepopexr}{poporange}{#1}},
  before lower={\poppill{Solution}{popink}{popcream}\par\vspace{6pt}}}
% Exercice (non résolu en place) — la correction est dans la partie Corrigés
\newtcolorbox{exercice}[1][]{popbase,colframe=popink,
  before upper={\stepcounter{popexo}\popboxhead{Exercice \thepopexo}{popink}{#1}}}
% Corrigé
\newtcolorbox{corrige}[1][]{popbase,colframe=popgreen,colback=popgreenL,
  before upper={\popboxhead{Corrigé}{popgreen}{#1}}}
% Objectifs / prérequis / bilan
\newtcolorbox{objectifs}{popbase,colframe=popink,colback=poporangeL,
  before upper={\popboxhead{Objectifs de la leçon}{popink}{}}}
\newtcolorbox{prerequis}{popbase,colframe=popink,colback=popblueL,
  before upper={\popboxhead{Prérequis}{popblue}{}}}
\newtcolorbox{bilan}{popbase,colframe=popink,colback=white,boxrule=2.8pt,
  before upper={\popboxhead{Fiche bilan}{poporange}{L'essentiel à connaître pour l'examen}}}
% Figure encadrée avec légende
\newtcolorbox{popfigure}[1][]{enhanced,breakable=false,colback=white,colframe=popline,boxrule=1.4pt,arc=8pt,
  left=10pt,right=10pt,top=10pt,bottom=8pt,before skip=10pt,after skip=12pt,halign=center,
  lower separated=false,halign lower=center,
  fontlower=\popsemi\fontsize{9}{11.5}\selectfont\color{popmuted},
  #1}
\newcommand{\legende}[1]{\par\vspace{6pt}{\popsemi\fontsize{9}{11.5}\selectfont\color{popmuted}#1}}

% ============================================================
% Signature OnBuch+
% ============================================================
\newcommand{\poplogoinline}[1]{{\poplogo\fontsize{#1}{#1}\selectfont\color{popink}OnBuch\color{poporange}+}}
\newcommand{\signatureonbuch}{%
  \begin{tcolorbox}[enhanced,colback=popink,colframe=popink,boxrule=2.2pt,arc=10pt,
      drop shadow=poporange,left=14pt,right=14pt,top=10pt,bottom=10pt,before skip=0pt,after skip=0pt]
    \tikz[baseline=-0.7ex]\node[circle,fill=popgreen,draw=popcream,line width=1.3pt,minimum size=22pt,inner sep=0]{\popcheck{white}};%
    \hspace{9pt}\begin{minipage}[c]{0.62\linewidth}
      {\popxbold\fontsize{12}{13}\selectfont\color{popcream}Signé\ \textbullet\ {\poplogo OnBuch\color{poporange}+}}\par\vspace{2pt}
      {\raggedright\popsemi\fontsize{8}{10.5}\selectfont\color{popline}\MakeUppercase{Équipe pédagogique OnBuch+}\\\MakeUppercase{Conforme au programme officiel MINESEC}\par}
    \end{minipage}\hfill
    {\poplogo\fontsize{22}{22}\selectfont\color{popcream}OnBuch\color{poporange}+}
  \end{tcolorbox}%
}

% ============================================================
% Page de garde
% #1 titre · #2 matière · #3 niveau · #4 module · #5 n° leçon · #6 durée ·
% #7 description / objectifs (texte ou itemize)
% ============================================================
\newcommand{\pagegarde}[7]{%
  \thispagestyle{empty}
  \newgeometry{a4paper,margin=1.7cm,top=1.6cm,bottom=1.4cm}%
  \begin{tikzpicture}[remember picture,overlay]
    \fill[poporange!18] ([xshift=-3.2cm,yshift=-3.6cm]current page.north east) circle (4.6cm);
    \fill[poppurple!14] ([xshift=2.4cm,yshift=7.6cm]current page.south west) circle (3.8cm);
  \end{tikzpicture}%
  {\poplogo\fontsize{30}{30}\selectfont\color{popink}OnBuch\color{poporange}+}\hfill
  \raisebox{0.6ex}{\poppill{Cours signé \textbullet\ Premium}{popink}{popcream}}\par
  \vspace{1pt}\popsquiggle{poppurple}\par
  \vspace{8pt}
  \begin{tcolorbox}[enhanced,colback=poporange,colframe=popink,boxrule=2.8pt,arc=13pt,
      drop shadow=popink,left=18pt,right=18pt,top=12pt,bottom=13pt,after skip=10pt]
    \poppill{#2 \textbullet\ #3}{popink}{popcream}\hspace{5pt}\poppill{#4}{white}{popink}\par\vspace{8pt}
    {\popdisplay\fontsize{14}{14}\selectfont\color{popink}LEÇON #5}\par\vspace{4pt}
    {\raggedright\hyphenpenalty=10000\exhyphenpenalty=10000\popdisplay\fontsize{25}{27}\selectfont\color{popcream}\MakeUppercase{#1}\par}
    \vspace{8pt}
    {\popxbold\fontsize{9.5}{9.5}\selectfont\color{popink}\MakeUppercase{Durée conseillée : #6}}
  \end{tcolorbox}
  \begin{tcolorbox}[enhanced,colback=white,colframe=popink,boxrule=2.2pt,arc=10pt,
      drop shadow=popink,left=16pt,right=16pt,top=10pt,bottom=10pt,after skip=8pt]
    \poppill{Dans cette leçon}{poppurple}{white}\par\vspace{5pt}
    \popsemi\fontsize{9.3}{12.6}\selectfont\color{popink}#7
  \end{tcolorbox}
  \noindent\begin{minipage}[t]{0.487\linewidth}
    \begin{tcolorbox}[enhanced,colback=popgreen,colframe=popink,boxrule=2.2pt,arc=10pt,
        drop shadow=popink,left=11pt,right=11pt,top=10pt,bottom=10pt,height=3.1cm,valign=center]
      \tcbox[colback=white,colframe=popink,boxrule=1.3pt,arc=5pt,boxsep=0pt,nobeforeafter,
        left=3pt,right=3pt,top=3pt,bottom=3pt,valign=center]{\qrcode[height=1.9cm]{https://play.google.com/store/apps/details?id=com.luvvix.onbuch&hl=fr}}%
      \hspace{8pt}\begin{minipage}[c]{0.5\linewidth}
        {\popdisplay\fontsize{11}{12.5}\selectfont\color{white}\MakeUppercase{Télécharge OnBuch}}\par\vspace{4pt}
        {\raggedright\popsemi\fontsize{8.4}{11}\selectfont\color{white}Gratuit \textbullet\ Google Play\\Tuteur IA, annales, concours, résultats\par}
      \end{minipage}
    \end{tcolorbox}
  \end{minipage}\hfill
  \begin{minipage}[t]{0.487\linewidth}
    \begin{tcolorbox}[enhanced,colback=poppurple,colframe=popink,boxrule=2.2pt,arc=10pt,
        drop shadow=popink,left=11pt,right=11pt,top=10pt,bottom=10pt,height=3.1cm,valign=center]
      \tikz[baseline=-0.7ex]\node[circle,fill=white,draw=popink,line width=1.3pt,minimum size=1.6cm,inner sep=0]{\popdisplay\fontsize{24}{24}\selectfont\color{poppurple}+};%
      \hspace{8pt}\begin{minipage}[c]{0.6\linewidth}
        {\popdisplay\fontsize{11}{12.5}\selectfont\color{white}\MakeUppercase{Passe à OnBuch+}}\par\vspace{4pt}
        {\raggedright\popsemi\fontsize{8.4}{11}\selectfont\color{white}Tous les cours signés, Tuteur IA illimité, fiches PDF — onglet OnBuch+ de l'app\par}
      \end{minipage}
    \end{tcolorbox}
  \end{minipage}
  \vspace{8pt}
  \popcontenu
  \vfill
  \signatureonbuch
  \restoregeometry
  \newpage
}

% Rangée « Ce cours contient » (page de garde)
\newcommand{\poptile}[3]{% #1 couleur #2 gros texte #3 légende
  \begin{tcolorbox}[enhanced,colback=#1,colframe=popink,boxrule=2pt,arc=9pt,shadow={2.5pt}{-2.5pt}{0pt}{popink},
      left=6pt,right=6pt,top=9pt,bottom=9pt,halign=center,valign=center,height=1.75cm,width=\linewidth,nobeforeafter]
    {\popdisplay\fontsize{12.5}{13}\selectfont\color{white}\MakeUppercase{#2}}\par\vspace{3pt}
    {\centering\popsemi\fontsize{7.8}{9.5}\selectfont\color{white}#3\par}
  \end{tcolorbox}}
\newcommand{\popcontenu}{%
  \par\noindent{\popxboldls\fontsize{7.5}{7.5}\selectfont\color{popmuted}CE COURS SIGNÉ CONTIENT}\par\vspace{7pt}
  \noindent\begin{minipage}[t]{0.235\linewidth}\poptile{popblue}{Cours}{complet, illustré, pas à pas}\end{minipage}\hfill
  \begin{minipage}[t]{0.235\linewidth}\poptile{poporange}{Méthodes}{et exercices résolus}\end{minipage}\hfill
  \begin{minipage}[t]{0.235\linewidth}\poptile{poppink}{Exercices}{type examen + corrigés}\end{minipage}\hfill
  \begin{minipage}[t]{0.235\linewidth}\poptile{popgreen}{Fiche bilan}{l'essentiel à retenir}\end{minipage}\par}

% Sommaire
\newtcolorbox{sommaire}{enhanced,colback=white,colframe=popink,boxrule=2.2pt,arc=10pt,
  drop shadow=popink,left=18pt,right=18pt,top=13pt,bottom=13pt,before skip=6pt,after skip=20pt,
  before upper={\poppill{Sommaire}{poppurple}{white}\par\vspace{9pt}\popsemi\fontsize{10.6}{14.5}\selectfont\color{popink}}}

% Signature de fin de cours — poussée tout en bas de la dernière page
\newcommand{\findecours}{%
  \par\vfill\nopagebreak
  \begin{center}\popsquiggle{poporange}\end{center}\vspace{4pt}
  \signatureonbuch
}
\AtBeginDocument{\def\llbracket{\mathopen{[\mkern-3.5mu[}}\def\rrbracket{\mathclose{]\mkern-3.5mu]}}} % intervalles d'entiers (glyphe ⟦ absent de la police)
% Glyphes absents de Fira Math : redéfinis à partir de glyphes disponibles
\AtBeginDocument{%
  \def\setminus{\mathbin{\backslash}}\let\smallsetminus\setminus
  \def\vdots{\mathord{\vbox{\baselineskip3.2pt\lineskiplimit0pt\kern4pt\hbox{.}\hbox{.}\hbox{.}}}}%
  \def\ddots{\mathinner{\mkern1mu\raise6.5pt\hbox{.}\mkern2mu\raise3.6pt\hbox{.}\mkern2mu\raise0.7pt\hbox{.}\mkern1mu}}%
  \def\triangle{\mathord{\scalebox{0.9}{$\symup{\Delta}$}}}%
  \def\nabla{\mathord{\raisebox{\depth}{\scalebox{1}[-1]{$\symup{\Delta}$}}}}%
  \def\checkmark{\popcheck{popdark}}%
  \def\star{\mathbin{\ast}}\def\bigstar{\mathord{\ast}}%
  \def\diamond{\mathbin{\scalebox{0.6}{$\Diamond$}}}%
  \def\bigcup{\mathop{\mathchoice{\raisebox{-0.3em}{\scalebox{1.7}{$\cup$}}}{\scalebox{1.25}{$\cup$}}{\cup}{\cup}}\displaylimits}%
  \def\bigcap{\mathop{\mathchoice{\raisebox{-0.3em}{\scalebox{1.7}{$\cap$}}}{\scalebox{1.25}{$\cap$}}{\cap}{\cap}}\displaylimits}%
  \def\lesseqqgtr{\mathrel{\lessgtr}}\def\bigoplus{\mathop{\oplus}\displaylimits}%
}
\providecommand{\ssi}{si et seulement si} % abréviation fréquente
\AtBeginDocument{\providecommand{\wideparen}{\overparen}} % arc de cercle

%% FIGFIX-BEGIN v4 — correctifs globaux des figures (docs/onbuch-generation/tools/figfix)
\usepackage{adjustbox}\usepackage{etoolbox}
% toute figure TikZ/pgfplots plus large que la ligne est réduite à la largeur disponible
\BeforeBeginEnvironment{tikzpicture}{\begin{adjustbox}{max width=\linewidth}}
\AfterEndEnvironment{tikzpicture}{\end{adjustbox}}
% légendes pgfplots lisibles par-dessus les courbes
\pgfplotsset{every axis/.append style={legend style={fill=white,fill opacity=0.92,text opacity=1,draw=black!15,inner sep=3pt}}}
% étiquettes d'arêtes (syntaxe edge["..."]) avec fond blanc
\tikzset{every edge quotes/.append style={fill=white,inner sep=1.2pt}}
%% FIGFIX-END
\begin{document}

\pagegarde{Leçon 11 — Expression orale : importance de la pratique du sport}{Chinois}{1ère A}{Module 2 : Environnement et santé}{ZHA11}{2 h}{Cette leçon d'expression orale apprend à l'élève de 1ère A à présenter un exposé simple et à participer à un débat guidé en chinois sur le thème « l'importance de la pratique du sport ». Elle fournit le vocabulaire sportif essentiel, les formules fonctionnelles pour présenter, donner son avis et relancer, ainsi qu'un entraînement ciblé sur les tons.
\vspace{4pt}
\begin{multicols}{2}\small
\begin{itemize}
\item À la fin de cette leçon, je sais utiliser le vocabulaire de base du sport et de la santé en chinois (运动, 健康, 锻炼, 比赛...).
\item À la fin de cette leçon, je sais présenter oralement un exposé simple de 5 à 8 phrases sur l'importance du sport.
\item À la fin de cette leçon, je sais donner mon avis avec les structures 我觉得..., 我认为..., 对我来说...
\item À la fin de cette leçon, je sais exprimer l'accord et le désaccord poliment (我同意 / 我不同意, 你说得对...).
\item À la fin de cette leçon, je sais relancer un échange avec des questions simples (你呢？你觉得怎么样？为什么？).
\item À la fin de cette leçon, je sais prononcer correctement les tons dans les mots clés de la leçon et éviter les confusions tonales.
\end{itemize}\end{multicols}}

\begin{sommaire}
\begin{multicols}{2}
\begin{enumerate}
\item Vocabulaire du sport et de la santé
\item Dialogue modèle 1 : parler de ses sports
\item Formules fonctionnelles : présenter, opiner, relancer
\item Travail sur les tons et la prononciation
\item Dialogue modèle 2 : mini-débat « faut-il faire du sport tous les jours ? »
\item Jeux de rôle et débat guidé en classe
\item Activité d'intégration
\item Méthodes et exercices résolus
\item Exercices
\item Corrigés des exercices
\item Fiche bilan
\end{enumerate}
\end{multicols}
\end{sommaire}

\begin{objectifs}
\begin{itemize}
\item À la fin de cette leçon, je sais utiliser le vocabulaire de base du sport et de la santé en chinois (运动, 健康, 锻炼, 比赛...).
\item À la fin de cette leçon, je sais présenter oralement un exposé simple de 5 à 8 phrases sur l'importance du sport.
\item À la fin de cette leçon, je sais donner mon avis avec les structures 我觉得..., 我认为..., 对我来说...
\item À la fin de cette leçon, je sais exprimer l'accord et le désaccord poliment (我同意 / 我不同意, 你说得对...).
\item À la fin de cette leçon, je sais relancer un échange avec des questions simples (你呢？你觉得怎么样？为什么？).
\item À la fin de cette leçon, je sais prononcer correctement les tons dans les mots clés de la leçon et éviter les confusions tonales.
\item À la fin de cette leçon, je sais participer à un jeu de rôle et à un débat guidé en respectant mon tour de parole.
\item À la fin de cette leçon, je sais comparer brièvement les pratiques sportives au Cameroun et en Chine.
\end{itemize}
\end{objectifs}

\begin{prerequis}
\begin{itemize}
\item Salutations, présentations et structures de base de Seconde (我叫..., 我是..., 我喜欢...)
\item Le verbe 喜欢 et les verbes modaux 会, 能, 应该
\item Les tons du mandarin (les 4 tons + ton neutre) et la lecture du pinyin
\item Le vocabulaire des activités quotidiennes et des loisirs (Module 1)
\item La structure 因为... 所以... vue en début d'année
\end{itemize}
\end{prerequis}

\newpage

\coursec{Leçon 11 — Expression orale : importance de la pratique du sport}

\courssub{Vocabulaire du sport et de la santé}

Au Cameroun, le football est une passion nationale : chaque victoire des Lions Indomptables rassemble tout le pays, de Yaoundé à Bamenda. En Chine aussi, le sport occupe une place centrale : le tai-chi au parc le matin, le tennis de table à l'école, et les Jeux olympiques de Pékin 2008 et 2022 ont marqué le monde. Imagine que ton école organise une journée sportive avec une délégation d'élèves chinois en visite : comment présenterais-tu, en chinois, l'importance du sport pour la santé et la vie des jeunes ?

\textbf{Problématique :} quels mots faut-il connaître pour parler du sport et de la santé en chinois, et comment les employer correctement dans une phrase ? Cette première section construit ton lexique : c'est la base indispensable de l'exposé oral et du débat que tu mèneras dans les sections suivantes.

\begin{textebox}[Li Ming se présente : le sport et moi]
大家好！我叫李明，我是中国学生。我非常喜欢运动。我每天早上跑步，因为运动对身体很好。\\
Dàjiā hǎo! Wǒ jiào Lǐ Míng, wǒ shì Zhōngguó xuésheng. Wǒ fēicháng xǐhuan yùndòng. Wǒ měitiān zǎoshang pǎobù, yīnwèi yùndòng duì shēntǐ hěn hǎo.\\
(Bonjour à tous ! Je m'appelle Li Ming, je suis un élève chinois. J'aime beaucoup le sport. Je cours chaque matin, car le sport est très bon pour le corps.)

在学校，我常常和同学们打篮球，有时候我们也打乒乓球。乒乓球很有意思！\\
Zài xuéxiào, wǒ chángcháng hé tóngxuémen dǎ lánqiú, yǒu shíhou wǒmen yě dǎ pīngpāngqiú. Pīngpāngqiú hěn yǒu yìsi!\\
(À l'école, je joue souvent au basket avec mes camarades, parfois nous jouons aussi au tennis de table. Le tennis de table est très intéressant !)

我知道喀麦隆人很喜欢足球。足球运动员身体很健康。运动很重要，也很有用！\\
Wǒ zhīdào Kāmàilóngrén hěn xǐhuan zúqiú. Zúqiú yùndòngyuán shēntǐ hěn jiànkāng. Yùndòng hěn zhòngyào, yě hěn yǒuyòng!\\
(Je sais que les Camerounais aiment beaucoup le football. Les footballeurs ont le corps en bonne santé. Le sport est très important, et aussi très utile !)
\end{textebox}

\textbf{Questions d'observation :}
\begin{enumerate}
\item Quels sports Li Ming pratique-t-il ? Souligne-les dans le texte.
\item Pourquoi court-il chaque matin ? Retrouve la structure \zh{因为……} (parce que…).
\item Quels mots expriment la fréquence de ses activités ?
\end{enumerate}

\begin{corrige}[Questions d'observation]
\begin{enumerate}
\item Li Ming court (\zh{跑步}), joue au basket (\zh{打篮球}) et au tennis de table (\zh{打乒乓球}).
\item Il court parce que le sport est bon pour le corps : \zh{因为运动对身体很好} (\emph{yīnwèi yùndòng duì shēntǐ hěn hǎo}). La structure \zh{因为} introduit la cause, comme « parce que » en français.
\item Les mots de fréquence sont : \zh{每天} (chaque jour), \zh{常常} (souvent), \zh{有时候} (parfois).
\end{enumerate}
\end{corrige}

\courssub{Les sports courants}

\begin{definition}[Le mot-clé : 运动]
\zh{运动} (\emph{yùndòng}, nom) signifie « sport, activité physique ». C'est le mot générique : on l'utilise pour parler du sport en général, jamais pour nommer un sport précis. Attention à la prononciation : deux quatrièmes tons descendants \cle{yùn-dòng}.
\end{definition}

\begin{center}
\begin{tabularx}{\linewidth}{|X|X|X|X|}
\hline
\rowcolor{popblueL} Caractères & Pinyin & Nature & Traduction \\
\hline
运动 & yùndòng & nom & sport, activité physique \\
\hline
足球 & zúqiú & nom & football \\
\hline
篮球 & lánqiú & nom & basketball \\
\hline
乒乓球 & pīngpāngqiú & nom & tennis de table \\
\hline
跑步 & pǎobù & verbe/nom & courir, course à pied \\
\hline
游泳 & yóuyǒng & verbe/nom & nager, natation \\
\hline
太极拳 & tàijíquán & nom & tai-chi \\
\hline
\end{tabularx}
\end{center}

Observe la formation des mots : \zh{足球} se compose de \zh{足} (pied) + \zh{球} (balle) : littéralement « balle de pied ». \zh{篮球} = \zh{篮} (panier) + \zh{球} (balle) : « balle de panier ». \zh{乒乓球} imite le son de la balle sur la table : \emph{pīng-pāng}. Le chinois construit souvent ses mots de façon logique et transparente : apprends à deviner le sens à partir des caractères.

\courssub{La santé et le corps}

\begin{center}
\begin{tabularx}{\linewidth}{|X|X|X|X|}
\hline
\rowcolor{popblueL} Caractères & Pinyin & Nature & Traduction \\
\hline
健康 & jiànkāng & nom/adjectif & santé ; en bonne santé \\
\hline
身体 & shēntǐ & nom & corps, santé physique \\
\hline
锻炼 & duànliàn & verbe & s'entraîner, faire de l'exercice \\
\hline
比赛 & bǐsài & nom/verbe & match, compétition ; disputer \\
\hline
队 & duì & nom & équipe \\
\hline
运动员 & yùndòngyuán & nom & sportif, athlète \\
\hline
\end{tabularx}
\end{center}

\begin{exemplebox}[La santé dans la phrase]
\zh{运动员的身体很健康。} \emph{Yùndòngyuán de shēntǐ hěn jiànkāng.} « Les sportifs ont le corps en bonne santé. »

\zh{我们每天锻炼身体。} \emph{Wǒmen měitiān duànliàn shēntǐ.} « Nous entraînons notre corps chaque jour. »

\zh{明天有足球比赛。} \emph{Míngtiān yǒu zúqiú bǐsài.} « Demain, il y a un match de football. »

\zh{我们队很好！} \emph{Wǒmen duì hěn hǎo!} « Notre équipe est très bonne ! »
\end{exemplebox}

Remarque la logique de \zh{运动员} : \zh{运动} (sport) + \zh{员} (membre, personne) = « personne du sport », c'est-à-dire sportif. Le suffixe \zh{员} forme des noms de métiers ou de membres, comme dans \zh{学员} (élève stagiaire).

\courssub{Carte mentale du lexique}

\begin{popfigure}
\begin{tikzpicture}[mindmap, grow cyclic, every node/.style={concept, align=center, font=\small},
root concept/.append style={concept color=popblue, font=\large},
level 1/.append style={level distance=3.2cm, sibling angle=120},
level 2/.append style={level distance=2.2cm, sibling angle=45}]
\node[root concept, align=center]{运动\\ \emph{yùndòng}\\ sport}
  child[concept color=poporange] { node {sports de balle}
    child { node[concept color=poporangeL, align=center]{足球\\ zúqiú} }
    child { node[concept color=poporangeL, align=center]{篮球\\ lánqiú} }
    child { node[concept color=poporangeL, align=center]{乒乓球\\ pīngpāngqiú} }
  }
  child[concept color=popgreen] { node {activités}
    child { node[concept color=popgreenL, align=center]{跑步\\ pǎobù} }
    child { node[concept color=popgreenL, align=center]{游泳\\ yóuyǒng} }
    child { node[concept color=popgreenL, align=center]{太极拳\\ tàijíquán} }
  }
  child[concept color=poppurple] { node {santé}
    child { node[concept color=poppurpleL, align=center]{健康\\ jiànkāng} }
    child { node[concept color=poppurpleL, align=center]{身体\\ shēntǐ} }
    child { node[concept color=poppurpleL, align=center]{锻炼\\ duànliàn} }
  };
\end{tikzpicture}
\legende{Carte mentale du vocabulaire : le sport et la santé autour du mot-clé 运动.}
\end{popfigure}

\courssub{Les verbes de jeu : 打, 踢, 练, 参加}

En français, on « joue » à tout : jouer au foot, jouer au basket, jouer au tennis de table. Le chinois, lui, choisit le verbe selon la \cle{partie du corps} utilisée ou la nature de l'activité. C'est une règle essentielle de cette leçon.

\begin{propriete}[Le verbe change selon la partie du corps]
En chinois, on dit \zh{打篮球} (jouer au basket, avec les mains) mais \zh{踢足球} (jouer au foot, avec les pieds) : le verbe change selon la partie du corps utilisée.
\begin{itemize}
\item \zh{打} \emph{dǎ} (frapper avec la main) + sports de balle joués avec les mains : \zh{打篮球}, \zh{打乒乓球}.
\item \zh{踢} \emph{tī} (donner un coup de pied) + sports joués avec les pieds : \zh{踢足球}.
\item \zh{练} \emph{liàn} (s'entraîner) + disciplines pratiquées seul : \zh{练太极拳}.
\item \zh{参加} \emph{cānjiā} (participer à) + événements : \zh{参加比赛} (participer à une compétition).
\item Pour la course et la natation, le verbe suffit seul : \zh{跑步}, \zh{游泳}.
\end{itemize}
\end{propriete}

\begin{center}
\begin{tabularx}{\linewidth}{|X|X|X|}
\hline
\rowcolor{popblueL} Verbe & Emploi & Exemple + traduction \\
\hline
打 dǎ (mains) & sports de balle avec les mains & 打篮球 dǎ lánqiú : jouer au basket \\
\hline
打 dǎ (mains) & idem & 打乒乓球 dǎ pīngpāngqiú : jouer au ping-pong \\
\hline
踢 tī (pieds) & sports de balle avec les pieds & 踢足球 tī zúqiú : jouer au foot \\
\hline
练 liàn & disciplines, entraînement & 练太极拳 liàn tàijíquán : pratiquer le tai-chi \\
\hline
参加 cānjiā & événements, compétitions & 参加比赛 cānjiā bǐsài : participer à un match \\
\hline
\end{tabularx}
\end{center}

\begin{attention}[Le piège du calque français]
Erreur fréquente des francophones : dire \zh{打足球} par calque du français « jouer au football ». C'est faux : le football se joue avec les pieds, il faut donc dire \zh{踢足球} \emph{tī zúqiú}. De même, on ne dit pas \zh{踢篮球} : le basket se joue avec les mains, donc \zh{打篮球}. Retiens : \cle{mains → 打, pieds → 踢}.
\end{attention}

\begin{exemplebox}[Exemples traduits]
\zh{他常常游泳。} \emph{Tā chángcháng yóuyǒng.} « Il nage souvent. »

\zh{我们参加篮球比赛。} \emph{Wǒmen cānjiā lánqiú bǐsài.} « Nous participons à un match de basket. »

\zh{同学们下午踢足球。} \emph{Tóngxuémen xiàwǔ tī zúqiú.} « Les camarades jouent au foot l'après-midi. »

\zh{我爷爷每天早上练太极拳。} \emph{Wǒ yéye měitiān zǎoshang liàn tàijíquán.} « Mon grand-père pratique le tai-chi chaque matin. »
\end{exemplebox}

\courssub{Les adverbes de fréquence}

Pour dire à quel rythme on pratique un sport, le chinois utilise des adverbes de fréquence placés \cle{avant le verbe}, entre le sujet et le verbe : \textbf{Sujet + adverbe de fréquence + verbe}.

\begin{center}
\begin{tabularx}{\linewidth}{|X|X|X|X|}
\hline
\rowcolor{popblueL} Caractères & Pinyin & Nature & Traduction \\
\hline
每天 & měitiān & adverbe & chaque jour \\
\hline
常常 & chángcháng & adverbe & souvent \\
\hline
有时候 & yǒu shíhou & adverbe & parfois \\
\hline
很少 & hěn shǎo & adverbe & rarement \\
\hline
从不 & cóng bù & adverbe & jamais (ne… jamais) \\
\hline
\end{tabularx}
\end{center}

\begin{exemplebox}[Du plus fréquent au plus rare]
\zh{我每天跑步。} \emph{Wǒ měitiān pǎobù.} « Je cours chaque jour. »

\zh{他常常打篮球。} \emph{Tā chángcháng dǎ lánqiú.} « Il joue souvent au basket. »

\zh{我们有时候游泳。} \emph{Wǒmen yǒu shíhou yóuyǒng.} « Nous nageons parfois. »

\zh{她很少踢足球。} \emph{Tā hěn shǎo tī zúqiú.} « Elle joue rarement au foot. »

\zh{我从不锻炼。} \emph{Wǒ cóng bù duànliàn.} « Je ne m'entraîne jamais. »
\end{exemplebox}

Comparaison avec le français : en français, « souvent » et « parfois » se placent souvent après le verbe (« il nage souvent ») ; en chinois, l'adverbe de fréquence se place \textbf{toujours avant le verbe}. Dire \zh{他游泳常常} est une faute d'ordre des mots.

\courssub{Les adjectifs d'opinion}

Pour donner ton avis sur un sport (compétence visée par l'exposé oral), tu as besoin d'adjectifs d'opinion. En chinois, l'adjectif se comporte comme un verbe : on dit directement \zh{运动很重要} « le sport est important », sans verbe « être ».

\begin{center}
\begin{tabularx}{\linewidth}{|X|X|X|X|}
\hline
\rowcolor{popblueL} Caractères & Pinyin & Nature & Traduction \\
\hline
重要 & zhòngyào & adjectif & important \\
\hline
有意思 & yǒu yìsi & adjectif & intéressant, amusant \\
\hline
累 & lèi & adjectif & fatigant, fatigué \\
\hline
好 & hǎo & adjectif & bon \\
\hline
有用 & yǒuyòng & adjectif & utile \\
\hline
\end{tabularx}
\end{center}

Structure : \textbf{Sujet + 很 + adjectif}. Le mot \zh{很} (\emph{hěn}, « très ») sert surtout de liaison naturelle entre le sujet et l'adjectif ; il n'insiste pas forcément.

\begin{exemplebox}[Donner son avis]
\zh{运动很重要。} \emph{Yùndòng hěn zhòngyào.} « Le sport est important. »

\zh{乒乓球很有意思。} \emph{Pīngpāngqiú hěn yǒu yìsi.} « Le ping-pong est très amusant. »

\zh{跑步很累，但是很有用。} \emph{Pǎobù hěn lèi, dànshì hěn yǒuyòng.} « Courir est fatigant, mais très utile. »

\zh{游泳对身体好。} \emph{Yóuyǒng duì shēntǐ hǎo.} « La natation est bonne pour le corps. »
\end{exemplebox}

Note la structure utile \zh{对……好} (\emph{duì… hǎo}) : « être bon pour… ». Exemple : \zh{运动对健康很好} « le sport est très bon pour la santé ». C'est une phrase-clé pour ton exposé.

\courssub{Prononciation ciblée : les tons du lexique}

Trois mots de cette leçon demandent un entraînement particulier des tons :

\begin{itemize}
\item \zh{运动} \emph{yùndòng} : 4e ton + 4e ton. Deux tons descendants, francs et secs, comme deux ordres brefs : \cle{yùn ! dòng !}. Ne les laisse pas « monter » à la française.
\item \zh{健康} \emph{jiànkāng} : 4e ton + 1er ton. On descend fort sur \emph{jiàn}, puis on tient un son haut et plat sur \emph{kāng}, comme une note de musique tenue.
\item \zh{锻炼} \emph{duànliàn} : 4e ton + 4e ton. Même profil que \zh{运动} : deux chutes nettes.
\end{itemize}

\begin{methode}[Bien prononcer les doubles quatrièmes tons]
\begin{enumerate}
\item Prononce chaque syllabe séparément, en descendant fortement la voix : \emph{yùn}… \emph{dòng}…
\item Enchaîne sans hésitation, en gardant les deux chutes : \emph{yùndòng}.
\item Place le mot dans une phrase courte et répète : \zh{我喜欢运动。} \emph{Wǒ xǐhuan yùndòng.}
\item Compare avec un mot à tons montants (\zh{足球} \emph{zúqiú}, 2e + 2e ton) pour sentir le contraste.
\end{enumerate}
\end{methode}

\begin{exoresolu}[Associer le bon verbe au bon sport]
Énoncé : complète chaque phrase avec le verbe correct : 打, 踢, 练 ou 参加.
\begin{enumerate}
\item 我们下午 \_\_\_\_ 足球。(Nous jouons au foot l'après-midi.)
\item 他常常 \_\_\_\_ 篮球。(Il joue souvent au basket.)
\item 爷爷早上 \_\_\_\_ 太极拳。(Grand-père pratique le tai-chi le matin.)
\item 我们 \_\_\_\_ 游泳比赛。(Nous participons à une compétition de natation.)
\end{enumerate}
\tcblower
Solution détaillée :
\begin{enumerate}
\item \zh{我们下午踢足球。} \emph{Wǒmen xiàwǔ tī zúqiú.} Le football se joue avec les pieds → \zh{踢}.
\item \zh{他常常打篮球。} \emph{Tā chángcháng dǎ lánqiú.} Le basket se joue avec les mains → \zh{打}.
\item \zh{爷爷早上练太极拳。} \emph{Yéye zǎoshang liàn tàijíquán.} Le tai-chi est une discipline d'entraînement → \zh{练}.
\item \zh{我们参加游泳比赛。} \emph{Wǒmen cānjiā yóuyǒng bǐsài.} Une compétition est un événement → \zh{参加}.
\end{enumerate}
Méthode : demande-toi toujours « avec quelle partie du corps joue-t-on ? » ou « est-ce un événement ? » avant de choisir le verbe.
\end{exoresolu}

\begin{exercice}[Entraînement écrit : vocabulaire et structures]
Traduis en chinois (caractères + pinyin) :
\begin{enumerate}
\item Je cours chaque matin.
\item Elle joue rarement au tennis de table.
\item Le sport est bon pour la santé.
\item Nous participons à un match de football.
\item Mon père ne fait jamais de sport.
\end{enumerate}
\end{exercice}

\begin{corrige}[Entraînement écrit : vocabulaire et structures]
\begin{enumerate}
\item \zh{我每天早上跑步。} \emph{Wǒ měitiān zǎoshang pǎobù.}
\item \zh{她很少打乒乓球。} \emph{Tā hěn shǎo dǎ pīngpāngqiú.} (ping-pong = mains → \zh{打} ; \zh{很少} avant le verbe.)
\item \zh{运动对健康很好。} \emph{Yùndòng duì jiànkāng hěn hǎo.} (structure \zh{对……好}.)
\item \zh{我们参加足球比赛。} \emph{Wǒmen cānjiā zúqiú bǐsài.}
\item \zh{我爸爸从不运动。} \emph{Wǒ bàba cóng bù yùndòng.} (\zh{从不} avant le verbe.)
\end{enumerate}
\end{corrige}

\begin{savaistu}[Le ping-pong, sport national chinois]
Le tennis de table \zh{乒乓球} est le sport national chinois : on y joue à l'école, dans les parcs, au bureau, et la Chine domine ce sport aux Jeux olympiques depuis des décennies. Au Cameroun, c'est le football \zh{足球} qui rassemble le plus de supporters : les Lions Indomptables sont une fierté nationale, et les enfants jouent au foot dans chaque quartier de Douala ou de Yaoundé. Quant au tai-chi \zh{太极拳}, pratiqué le matin dans les parcs chinois, il montre que le sport peut aussi être doux et bon pour la santé à tout âge : \zh{运动对身体很好} !
\end{savaistu}

\courssub{Dialogue modèle 1 : parler de ses sports}

Ce dialogue met en scène Li Ming (chinois) et Paul (camerounais) qui font connaissance lors de la journée sportive. Observe les formules de présentation et les verbes de sport.

\begin{textebox}[Dialogue : Paul et Li Ming au stade]
李明：保罗，你喜欢运动吗？\\
Lǐ Míng: Bǎoluó, nǐ xǐhuan yùndòng ma?\\
(Li Ming : Paul, aimes-tu le sport ?)

保罗：喜欢。我最喜欢足球。喀麦隆人都爱踢足球！你呢？\\
Bǎoluó: Xǐhuan. Wǒ zuì xǐhuan zúqiú. Kāmàilóngrén dōu ài tī zúqiú! Nǐ ne?\\
(Paul : Oui. J'aime le football par-dessus tout. Les Camerounais aiment tous jouer au foot ! Et toi ?)

李明：我喜欢打篮球和乒乓球。我也每天早上跑步，锻炼身体。\\
Lǐ Míng: Wǒ xǐhuan dǎ lánqiú hé pīngpāngqiú. Wǒ yě měitiān zǎoshang pǎobù, duànliàn shēntǐ.\\
(Li Ming : J'aime jouer au basket et au ping-pong. Je cours aussi chaque matin pour m'entraîner.)

保罗：跑步很累吧？\\
Bǎoluó: Pǎobù hěn lèi ba?\\
(Paul : Courir est fatigant, non ?)

李明：有点儿累，但是运动对健康很有用。我们一起去打球吧！\\
Lǐ Míng: Yǒudiǎnr lèi, dànshì yùndòng duì jiànkāng hěn yǒuyòng. Wǒmen yīqǐ qù dǎ qiú ba!\\
(Li Ming : Un peu fatiguant, mais le sport est très utile pour la santé. Allons jouer au ballon ensemble !)

保罗：好啊！我们去打篮球。\\
Bǎoluó: Hǎo a! Wǒmen qù dǎ lánqiú.\\
(Paul : D'accord ! Allons jouer au basket.)
\end{textebox}

\begin{aretenir}[Points clés du dialogue]
\begin{itemize}
\item \zh{你呢？} (\emph{Nǐ ne?}) : « Et toi ? » (relance simple).
\item \zh{最喜欢} (\emph{zuì xǐhuan}) : « aimer le plus / préférer ».
\item \zh{都爱} (\emph{dōu ài}) : « tous aiment » (\zh{都} placé avant le verbe).
\item \zh{有点儿} (\emph{yǒudiǎnr}) : « un peu » (souvent pour un aspect négatif : fatigant, difficile).
\item \zh{一起} (\emph{yīqǐ}) : « ensemble » (placé avant le verbe : \zh{一起去}).
\item \zh{好啊！} (\emph{Hǎo a!}) : « D'accord ! / Avec plaisir ! » (particule \zh{啊} pour l'enthousiasme).
\end{itemize}
\end{aretenir}

\courssub{Formules fonctionnelles : présenter, opiner, relancer}

Pour réussir ton exposé oral et participer au débat, tu dois maîtriser des « blocs de langue » prêts à l'emploi. Voici les formules essentielles classées par fonction.

\begin{propriete}[Formules pour l'exposé et le débat]
\begin{center}
\begin{tabularx}{\linewidth}{|X|X|X|}
\hline
\rowcolor{popblueL} Fonction & Chinois (Pinyin) & Français \\
\hline
\textbf{Présenter son sujet} & \zh{今天我要谈谈……} (\emph{Jīntiān wǒ yào tántan…}) & Aujourd'hui, je vais parler de… \\
& \zh{我的题目是……} (\emph{Wǒ de tímù shì…}) & Mon sujet est… \\
\hline
\textbf{Structurer} & \zh{首先……} (\emph{Shǒuxiān…}) & D'abord… \\
& \zh{其次……} (\emph{Qícì…}) & Ensuite… / En second lieu… \\
& \zh{最后……} (\emph{Zuìhòu…}) & Enfin… \\
\hline
\textbf{Donner son avis} & \zh{我认为……} (\emph{Wǒ rènwéi…}) & Je pense que… / J'estime que… \\
& \zh{我觉得……} (\emph{Wǒ juéde…}) & Je trouve que… / J'ai l'impression que… \\
& \zh{在我看来……} (\emph{Zài wǒ kànlái…}) & À mon avis… / Selon moi… \\
\hline
\textbf{Argumenter} & \zh{因为……所以……} (\emph{Yīnwèi… suǒyǐ…}) & Parce que… donc… \\
& \zh{一方面……另一方面……} (\emph{Yī fāngmiàn… lìng yī fāngmiàn…}) & D'un côté… de l'autre… \\
\hline
\textbf{Exprimer accord / désaccord} & \zh{我同意。/ 我赞成。} (\emph{Wǒ tóngyì. / Wǒ zànchéng.}) & Je suis d'accord. / J'approuve. \\
& \zh{我不太同意。} (\emph{Wǒ bù tài tóngyì.}) & Je ne suis pas tout à fait d'accord. \\
& \zh{恐怕不太对。} (\emph{Kǒngpà bù tài duì.}) & Je crains que ce ne soit pas exact. \\
\hline
\textbf{Relancer / Interagir} & \zh{你怎么看？} (\emph{Nǐ zěnme kàn?}) & Qu'en penses-tu ? / Quel est ton avis ? \\
& \zh{对吗？} (\emph{Duì ma?}) & N'est-ce pas ? / D'accord ? \\
& \zh{比如说呢？} (\emph{Bǐrú shuō ne?}) & Par exemple ? \\
\hline
\textbf{Conclure} & \zh{总之……} (\emph{Zǒngzhī…}) & En résumé… / Bref… \\
& \zh{谢谢大家。} (\emph{Xièxie dàjiā.}) & Merci à tous. \\
\hline
\end{tabularx}
\end{center}
\end{propriete}

\begin{exemplebox}[Mini-exposé modèle (1 minute)]
\textbf{Sujet :} \zh{运动的重要性} (\emph{Yùndòng de zhòngyàoxìng}) — L'importance du sport.

\zh{大家好。今天我要谈谈运动的重要性。} (\emph{Dàjiā hǎo. Jīntiān wǒ yào tántan yùndòng de zhòngyàoxìng.})
« Bonjour à tous. Aujourd'hui je vais parler de l'importance du sport. »

\zh{首先，运动对身体很好。} (\emph{Shǒuxiān, yùndòng duì shēntǐ hěn hǎo.})
« D'abord, le sport est très bon pour le corps. »

\zh{其次，运动能让我们心情好。} (\emph{Qícì, yùndòng néng ràng wǒmen xīnqíng hǎo.})
« Ensuite, le sport nous met de bonne humeur. » (\zh{让} = faire en sorte que)

\zh{最后，我认为每天锻炼很有必要。} (\emph{Zuìhòu, wǒ rènwéi měitiān duànliàn hěn yǒu bìyào.})
« Enfin, je pense qu'il est nécessaire de s'entraîner chaque jour. »

\zh{总之，请大家多运动！谢谢大家。} (\emph{Zǒngzhī, qǐng dàjiā duō yùndòng! Xièxie dàjiā.})
« En résumé, faites tous plus de sport ! Merci à tous. »
\end{exemplebox}

\begin{attention}[Pièges à l'oral]
\begin{itemize}
\item Ne confonds pas \zh{觉得} (sentiment, impression subjective) et \zh{认为} (opinion raisonnée, jugement). Pour un exposé formel, \zh{我认为} est plus approprié.
\item \zh{因为……所以……} : en chinois, les deux termes sont souvent présents (contrairement au français où « donc » est parfois omis). Ne dis pas seulement \zh{因为……} sans conclusion.
\item \zh{心情} (\emph{xīnqíng}) = humeur, moral. Ne pas confondre avec \zh{身体} (corps physique).
\end{itemize}
\end{attention}

\courssub{Travail sur les tons et la prononciation}

\begin{methode}[Entraînement prosodique : le « ton sandwich » (4-4 et 4-1)]
Les mots \zh{运动} (4-4), \zh{锻炼} (4-4) et \zh{健康} (4-1) sont des « pièges à tons » pour les francophones.
\begin{enumerate}
\item \textbf{Isolation :} Répète chaque syllabe en exagérant la chute (4e) ou la tenue haute (1er). \emph{Yùn! Dòng!} \quad \emph{Jiàn! Kāng—}
\item \textbf{Enchaînement :} Lie les syllabes sans pause. Le 4e ton reste une chute brève, pas un son long.
\item \textbf{En phrase (contexte) :} Le 4e ton final d'un mot peut s'adoucir légèrement s'il est suivi d'un autre 4e ton (règle du « demi-4e ton »), mais en articulation soignée (exposé), on garde la chute nette.
\item \textbf{Exercices de contraste :} Lis à haute voix : \zh{足球} (2-2, montants) vs \zh{运动} (4-4, descendants) ; \zh{游泳} (2-3) vs \zh{锻炼} (4-4).
\end{enumerate}
\end{methode}

\begin{experience}[Atelier prononciation : « Le championnat des tons »]
\begin{enumerate}
\item Par binômes. L'élève A lit une phrase du vocabulaire (ex: \zh{我每天跑步}). L'élève B note les tons entendus sur un mini-tableau (ex: 3-1-2-3-4).
\item Inversion des rôles.
\item L'enseignant dicte 5 phrases « pièges » mêlant \zh{运动}, \zh{健康}, \zh{锻炼}, \zh{足球}, \zh{篮球}. Les élèves écrivent en pinyin avec tons.
\item Correction collective au tableau : focus sur \zh{yùndòng} vs \zh{zúqiú}.
\end{enumerate}
\end{experience}

\courssub{Dialogue modèle 2 : mini-débat « Faut-il faire du sport tous les jours ? »}

Ce dialogue illustre un échange d'opinions structuré. Note l'utilisation des formules fonctionnelles (accord, désaccord nuancé, relance).

\begin{textebox}[Débat : 每天运动，好不好？]
陈老师：同学们，我们来辩论一下：“每天运动，好不好？”小王，你先说。\\
Chén lǎoshī: Tóngxuémen, wǒmen lái biànlùn yīxià: “Měitiān yùndòng, hǎo bu hǎo?” Xiǎo Wáng, nǐ xiān shuō.\\
(Mme Chen : Camarades, débattons : « Faire du sport tous les jours, bien ou pas bien ? » Xiao Wang, à toi.)

小王：我认为每天运动很好。因为运动对健康有用，身体会更强壮。\\
Xiǎo Wáng: Wǒ rènwéi měitiān yùndòng hěn hǎo. Yīnwèi yùndòng duì jiànkāng yǒuyòng, shēntǐ huì gèng qiángzhuàng.\\
(Xiao Wang : Je pense que faire du sport chaque jour est très bien. Parce que le sport est utile pour la santé, le corps devient plus fort.)

小李：我不同意。我觉得每天运动太累了。学生作业多，没有时间。\\
Xiǎo Lǐ: Wǒ bù tóngyì. Wǒ juéde měitiān yùndòng tài lèi le. Xuésheng zuòyè duō, méiyǒu shíjiān.\\
(Xiao Li : Je ne suis pas d'accord. Je trouve que faire du sport tous les jours est trop fatigant. Les élèves ont beaucoup de devoirs, pas de temps.)

小王：可是，不运动身体会不好。你怎么看？\\
Xiǎo Wáng: Kěshì, bù yùndòng shēntǐ huì bù hǎo. Nǐ zěnme kàn?\\
(Xiao Wang : Mais si on ne fait pas de sport, le corps ira mal. Qu'en penses-tu ?)

小李：我也赞成运动，但是一周三次就够了。休息也重要。\\
Xiǎo Lǐ: Wǒ yě zànchéng yùndòng, dànshì yī zhōu sān cì jiù gòu le. Xiūxi yě zhòngyào.\\
(Xiao Li : Moi aussi j'approuve le sport, mais trois fois par semaine suffit. Le repos est aussi important.)

陈老师：很好。双方都有道理。一方面要锻炼，另一方面要休息。平衡最重要。\\
Chén lǎoshī: Hěn hǎo. Shuāngfāng dōu yǒu dàolǐ. Yī fāngmiàn yào duànliàn, lìng yī fāngmiàn yào xiūxi. Pínghéng zuì zhòngyào.\\
(Mme Chen : Très bien. Les deux côtés ont raison. D'un côté il faut s'entraîner, de l'autre se reposer. L'équilibre est le plus important.)
\end{textebox}

\begin{aretenir}[Structures du débat]
\begin{itemize}
\item \zh{辩论} (\emph{biànlùn}) : débattre.
\item \zh{太……了} (\emph{tài… le}) : « trop… » (insistance). Ex: \zh{太累了}.
\item \zh{可是} (\emph{kěshì}) / \zh{但是} (\emph{dànshì}) : « mais / cependant ».
\item \zh{也} (\emph{yě}) : « aussi / moi aussi » (placé avant le verbe : \zh{我也赞成}).
\item \zh{就够了} (\emph{jiù gòu le}) : « cela suffit / c'est assez ».
\item \zh{有道理} (\emph{yǒu dàolǐ}) : « avoir raison / être fondé ».
\item \zh{平衡} (\emph{pínghéng}) : équilibre (nom/verbe).
\end{itemize}
\end{aretenir}

\courssub{Jeux de rôle et débat guidé en classe}

\begin{experience}[Activité 1 : Jeu de rôle — « Journée sportive sino-camerounaise »]
\textbf{Consigne :} Par groupes de 3 (1 élève chinois en échange, 2 élèves camerounais). Durée : 5 min préparation + 2 min jeu par groupe.
\textbf{Rôles :}
\begin{itemize}
\item \textbf{A (Chinois) :} Présente ton sport favori (tai-chi, ping-pong, badminton) et ses bienfaits. Utilise \zh{我认为}, \zh{因为……所以……}.
\item \textbf{B (Camerounais) :} Présente le football (ou un autre sport local : lutte traditionnelle, handball). Utilise \zh{我觉得}, \zh{对……好}.
\item \textbf{C (Journaliste) :} Pose des questions pour relancer : \zh{你怎么看？}, \zh{比如说呢？}, \zh{对吗？}.
\end{itemize}
\textbf{Contraintes linguistiques :} Minimum 5 phrases par élève. Utiliser au moins 3 verbes de sport (\zh{打}, \zh{踢}, \zh{练}, \zh{跑}, \zh{游}) et 2 adverbes de fréquence. Soigner les tons 4-4 (\zh{运动}, \zh{锻炼}).
\end{experience}

\begin{experience}[Activité 2 : Débat guidé — « Le sport tous les jours : pour ou contre ? »]
\textbf{Organisation :} Deux camps (Pour / Contre) + Jury (3 élèves).
\textbf{Étapes :}
\begin{enumerate}
\item \textbf{Préparation (10 min) :} Chaque camp rédige 3 arguments sur fiche (mots-clés en chinois + pinyin). Ex: \zh{Pour: 健康, 强壮, 心情好. Contre: 太累, 没时间, 受伤}. 
\item \textbf{Exposé d'ouverture (1 min/camp) :} Un porte-parole lit l'exposé structuré (\zh{首先… 其次… 最后…}).
\item \textbf{Échanges libres (5 min) :} Le jury donne la parole alternativement. Obligation d'utiliser \zh{我同意/不同意}, \zh{可是}, \zh{你怎么看？}.
\item \textbf{Synthèse du Jury (1 min) :} \zh{总之… 平衡最重要}. Vote de la meilleure argumentation (prononciation + structures + vocabulaire).
\end{enumerate}
\textbf{Évaluation (grille simplifiée) :}
\begin{center}
\begin{tabularx}{\linewidth}{|X|X|X|}
\hline
\rowcolor{popblueL} Critère & Indicateur & Points \\
\hline
Prononciation (Tons) & Tons 4-4 (\zh{运动}, \zh{锻炼}) corrects, 4-1 (\zh{健康}) stable & /5 \\
\hline
Vocabulaire & Verbes sport (\zh{打}/\zh{踢}/\zh{练}), adverbes fréquence, adjectifs opinion & /5 \\
\hline
Structures & \zh{因为…所以…}, \zh{我认为}, \zh{一方面…}, \zh{对…好} & /5 \\
\hline
Interaction & Relance (\zh{你怎么看?}), réactivité, fluidité & /5 \\
\hline
\end{tabularx}
\end{center}
\end{experience}

\courssub{Exercices de consolidation (Écrit / Préparation orale)}

\begin{exercice}[Préparation de ton exposé : « Mon sport et moi »]
Rédige un court exposé écrit (80-100 caractères) que tu présenteras oralement à la prochaine séance. Structure imposée :
\begin{enumerate}
\item Introduction : \zh{大家好，今天我要谈谈……}
\item Argument 1 (Santé) : \zh{首先，运动对……很好。}
\item Argument 2 (Moral/Plaisir) : \zh{其次，……很有意思 / 能让……}
\item Fréquence personnelle : \zh{我每天/常常/有时候……} (verbe + sport).
\item Conclusion : \zh{总之，运动很重要。谢谢大家。}
\end{enumerate}
\textbf{Aide :} Utilise le vocabulaire des tableaux (sports, santé, verbes, fréquence, adjectifs). Vérifie l'ordre des mots (Sujet + Adv + Verbe / Sujet + 很 + Adj).
\end{exercice}

\begin{exercice}[Traduction thème : du français vers le chinois]
Traduis les phrases suivantes en chinois (caractères + pinyin).
\begin{enumerate}
\item Je pense que le basketball est très intéressant.
\item Mon frère joue au football tous les après-midis.
\item Faire du sport est bon pour la santé, mais il ne faut pas être trop fatigué.
\item À mon avis, trois fois par semaine, c'est suffisant.
\item Est-ce que tu es d'accord ? Qu'en penses-tu ?
\end{enumerate}
\end{exercice}

\begin{corrige}[Traduction thème]
\begin{enumerate}
\item \zh{我认为篮球很有意思。} \emph{Wǒ rènwéi lánqiú hěn yǒu yìsi.}
\item \zh{我哥哥每天下午踢足球。} \emph{Wǒ gēge měitiān xiàwǔ tī zúqiú.} (Préciser \zh{哥哥} grand frère ou \zh{弟弟} petit frère ; \zh{每天下午} avant verbe).
\item \zh{运动对健康很好，但是不要太累。} \emph{Yùndòng duì jiànkāng hěn hǎo, dànshì bùyào tài lèi.} (\zh{不要} impératif négatif).
\item \zh{在我看来，一周三次就够了。} \emph{Zài wǒ kànlái, yī zhōu sān cì jiù gòu le.}
\item \zh{你同意吗？你怎么看？} \emph{Nǐ tóngyì ma? Nǐ zěnme kàn?}
\end{enumerate}
\end{corrige}

\begin{aretenir}[Synthèse de la leçon 11 — Check-list pour l'examen oral]
\begin{itemize}
\item \textbf{Lexique :} 10+ mots sport/santé (\zh{运动, 足球, 篮球, 乒乓球, 跑步, 游泳, 太极拳, 健康, 身体, 锻炼, 运动员, 比赛, 队}).
\item \textbf{Grammaire clé :} Verbes sport (\zh{打}/\zh{踢}/\zh{练}/\zh{参加}), Adverbes fréquence (place avant V), Structure \zh{因为…所以…}, \zh{对…好}, \zh{太…了}, \zh{就够了}.
\item \textbf{Fonctionnel :} Présenter (\zh{今天我要谈谈…}), Opiner (\zh{我认为/我觉得}), Relancer (\zh{你怎么看?}), Conclure (\zh{总之}).
\item \textbf{Prononciation :} Maîtrise des tons 4-4 (\zh{运动, 锻炼}) et 4-1 (\zh{健康}). Intonation interrogative (\zh{吗}, \zh{呢}).
\item \textbf{Culture :} Ping-pong (Chine) vs Football (Cameroun) ; Tai-chi (santé douce) ; Équilibre effort/repos.
\end{itemize}
\end{aretenir}

\coursec{Leçon 11 — Expression orale : importance de la pratique du sport}

\courssub{Vocabulaire du sport et de la santé : mots-outils et verbes d'action}

Avant d'entrer dans le dialogue, fixons le lexique de base. En chinois, \textbf{aimer un sport} ne se dit pas \zh{喜欢 + nom du sport} (sauf pour la natation, la course), mais \zh{喜欢 + verbe d'action + nom du sport}. Le choix du verbe dépend du geste technique : c'est la première grande différence avec le français.

\begin{center}
\begin{tabularx}{\linewidth}{|X|X|X|X|}
\hline
\rowcolor{popblueL} Caractères & Pinyin & Nature & Traduction \\
\hline
运动 & yùndòng & nom / verbe & sport, exercice physique ; faire du sport \\
\hline
锻炼 & duànliàn & verbe & s'entraîner, faire de l'exercice (général) \\
\hline
健康 & jiànkāng & nom / adj. & santé ; sain, en bonne santé \\
\hline
身体 & shēntǐ & nom & corps, condition physique \\
\hline
踢足球 & tī zúqiú & v. + n. & jouer au football (lit. donner un coup de pied au ballon) \\
\hline
打篮球 & dǎ lánqiú & v. + n. & jouer au basket-ball (lit. frapper le ballon) \\
\hline
打乒乓球 & dǎ pīngpāngqiú & v. + n. & jouer au tennis de table (ping-pong) \\
\hline
打排球 & dǎ páiqiú & v. + n. & jouer au volley-ball \\
\hline
游泳 & yóuyǒng & verbe & nager, faire de la natation \\
\hline
跑步 & pǎobù & verbe & courir, faire de la course à pied \\
\hline
喜欢 & xǐhuan & verbe & aimer, apprécier \\
\hline
觉得 & juéde & verbe & trouver, penser (avis subjectif) \\
\hline
有意思 & yǒuyìsi & adjectif & intéressant, amusant \\
\hline
常常 & chángcháng & adverbe & souvent, fréquemment \\
\hline
每天 & měitiān & compl. de temps & tous les jours \\
\hline
都 & dōu & adverbe & tous, tout (marque la totalité) \\
\hline
因为 & yīnwèi & conjonction & parce que (cause) \\
\hline
所以 & suǒyǐ & conjonction & donc, par conséquent (conséquence) \\
\hline
对 & duì & préposition & pour, envers (marque le bénéficiaire) \\
\hline
重要 & zhòngyào & adjectif & important \\
\hline
应该 & yīnggāi & verbe modal & devoir, devrait (obligation morale/conseil) \\
\hline
最好 & zuìhǎo & adv. + adj. & il vaudrait mieux, le mieux est de \\
\hline
但是 & dànshì & conjonction & mais, cependant \\
\hline
也 & yě & adverbe & aussi, également \\
\hline
\end{tabularx}
\end{center}

\begin{aretenir}[Règle d'or : le verbe d'action obligatoire]
Pour les sports de balle ou d'opposition, le chinois \textbf{exige} un verbe d'action devant le nom du sport après \zh{喜欢}, \zh{会}, \zh{想} :
\begin{itemize}
\item \zh{踢} (tī, pied) $\rightarrow$ \zh{踢足球} (football), \zh{踢毽子} (jeu de volant).
\item \zh{打} (dǎ, main/raquette) $\rightarrow$ \zh{打篮球}, \zh{打乒乓球}, \zh{打排球}, \zh{打羽毛球} (badminton), \zh{打网球} (tennis).
\item \zh{游} (yóu, nager) $\rightarrow$ \zh{游泳} (natation) — \zh{游} est déjà le verbe, pas de nom ajouté.
\item \zh{跑} (pǎo, courir) $\rightarrow$ \zh{跑步} (course) — idem.
\end{itemize}
\emph{Erreur type :} \zh{我喜欢足球} (je suis fan de foot / j'aime le ballon) $\neq$ \zh{我喜欢踢足球} (j'aime jouer au foot).
\end{aretenir}

\courssub{Dialogue modèle 1 : parler de ses sports préférés}

Ce dialogue est la \cle{colonne vertébrale} de la leçon. Il met en scène deux élèves camerounais, \zh{阿明} (Àmíng) et \zh{小红} (Xiǎohóng), au lycée de Yaoundé. Il active les cinq fonctions cibles : saluer, demander, répondre, relancer, justifier.

\begin{textebox}[Dialogue 1 : 你喜欢什么运动？ — À lire à voix haute, marquer les tons au crayon]
阿明：你好！你喜欢什么运动？\\
\emph{Āmíng: Nǐ hǎo ! Nǐ xǐhuan shénme yùndòng ?}\\
« Salut ! Quel sport aimes-tu ? »\\[4pt]
小红：我喜欢踢足球，你呢？\\
\emph{Xiǎohóng: Wǒ xǐhuan tī zúqiú, nǐ ne ?}\\
« J'aime jouer au football, et toi ? »\\[4pt]
阿明：我喜欢打篮球。我觉得篮球很有意思。\\
\emph{Āmíng: Wǒ xǐhuan dǎ lánqiú. Wǒ juéde lánqiú hěn yǒuyìsi.}\\
« J'aime jouer au basket-ball. Je trouve que le basket, c'est très intéressant. »\\[4pt]
小红：你常常锻炼吗？\\
\emph{Xiǎohóng: Nǐ chángcháng duànliàn ma ?}\\
« Tu t'entraînes souvent ? »\\[4pt]
阿明：对，我每天都锻炼，因为运动对健康很重要。\\
\emph{Āmíng: Duì, wǒ měitiān dōu duànliàn, yīnwèi yùndòng duì jiànkāng hěn zhòngyào.}\\
« Oui, je m'entraîne tous les jours, parce que le sport est très important pour la santé. »
\end{textebox}

\textbf{Consigne de première lecture (3 min).}
\begin{enumerate}
\item Lis en suivant le pinyin du doigt. Deuxième lecture : marque au crayon le signe du ton (ā á ǎ à) au-dessus de \textbf{chaque syllabe}.
\item Réponds oralement (en chinois) : 
\begin{enumerate}
\item Qui pose la question initiale ? Avec quel mot interrogatif ?
\item Comment \zh{小红} relance-t-elle sans répéter le verbe ?
\item Quel mot \zh{阿明} utilise-t-il pour introduire la cause ?
\end{enumerate}
\end{enumerate}

\textbf{Réponses d'observation.}
(1) \zh{阿明} pose la question avec \zh{什么} (\emph{shénme}). (2) \zh{小红} utilise la particule \zh{呢} (\emph{ne}) : \zh{你呢？} (« et toi ? »). (3) \zh{阿明} utilise \zh{因为} (\emph{yīnwèi}, « parce que »).

\courssub{La logique du dialogue : l'escalier en cinq marches}

Un échange oral structuré ne s'improvise pas : il gravit des marches précises. Le schéma ci-dessous visualise l'enchaînement pragmatique.

\begin{popfigure}
\begin{tikzpicture}[
  stepS/.style={rectangle, draw=popblue, fill=popblueL, rounded corners=4pt, minimum width=3.2cm, minimum height=1cm, align=center, font=\small\bfseries},
  arr/.style={-{Stealth[length=3mm]}, thick, poporange},
  node distance=1.2cm and 1.5cm
]
\node[stepS, align=center] (s1){\textbf{1. Saluer}\\你好！};
\node[stepS, above right=of s1, align=center] (s2){\textbf{2. Questionner}\\你喜欢什么运动？};
\node[stepS, above right=of s2, align=center] (s3){\textbf{3. Répondre}\\我喜欢踢足球。};
\node[stepS, below right=of s3, align=center] (s4){\textbf{4. Relancer}\\你呢？};
\node[stepS, below right=of s4, align=center] (s5){\textbf{5. Argumenter}\\我觉得... 因为...};
\draw[arr] (s1) -- (s2);
\draw[arr] (s2) -- (s3);
\draw[arr] (s3) -- (s4);
\draw[arr] (s4) -- (s5);
\end{tikzpicture}
\legende{Les cinq marches de l'échange oral : saluer $\rightarrow$ questionner $\rightarrow$ répondre $\rightarrow$ relancer (\zh{呢}) $\rightarrow$ argumenter (\zh{觉得}/\zh{因为}).}
\end{popfigure}

\begin{aretenir}[Les 5 marches à automatiser]
\zh{你好！} $\rightarrow$ \zh{你喜欢什么运动？} $\rightarrow$ \zh{我喜欢 + [verbe+sport].} $\rightarrow$ \zh{你呢？} $\rightarrow$ \zh{我觉得...很...，因为...对...很重要。}
\end{aretenir}

\courssub{Formules fonctionnelles \& Grammaire en action}

Quatre structures grammaticales portent ce dialogue. Maîtrise-les pour les réutiliser en débat (section 5).

\begin{propriete}[F1 — Questionner les goûts : \zh{什么} en position d'objet]
\textbf{Structure :} Sujet + \zh{喜欢} + \zh{什么} + Nom ? \\
Le mot interrogatif \zh{什么} (\emph{shénme}) reste \textbf{à la place de la réponse} (in situ). Pas d'inversion, pas de \zh{吗}.
\zh{你喜欢什么运动？} \emph{Nǐ xǐhuan shénme yùndòng ?} \\
Réponse : \zh{我喜欢踢足球。} (\zh{什么} $\leftarrow$ \zh{踢足球}).
\end{propriete}

\begin{propriete}[F2 — Relancer l'échange : la particule \zh{呢} (ton neutre)]
\textbf{Structure :} Sujet (pronom/nom) + \zh{呢} ? \\
\zh{呢} (\emph{ne}, \textbf{ton neutre} : court, léger, sans contour de ton) évite la répétition. \zh{你呢？} = « Et toi ? » (sous-entendu : \zh{你喜欢什么运动？}).
\zh{我喜欢游泳，你呢？} \emph{Wǒ xǐhuan yóuyǒng, nǐ ne ?}
\end{propriete}

\begin{attention}[Le piège du ton neutre \zh{呢}]
\emph{Ne prononce pas} \zh{né} (2e ton montant) ni \zh{nè} (4e ton tombant). Le ton neutre n'a pas de cible de hauteur fixe : il est \textbf{court, faible, attaché à la syllabe précédente}. Exercice : \emph{nǐ ne -- nǐ ne -- nǐ ne} (le \emph{ne} dure la moitié du \emph{nǐ}). Enregistre-toi et compare à l'enseignant.
\end{attention}

\begin{propriete}[F3 — Exprimer un avis : \zh{我觉得} + \zh{很} + adjectif]
\textbf{Structure :} \zh{我觉得} + Sujet + \zh{很} + Adjectif. \\
\zh{觉得} (\emph{juéde}) = « trouver, avoir l'impression que ». \textbf{Règle cruciale :} en chinois, un adjectif predicatif (seul après le sujet) sonne incomplet à l'affirmatif. On insère \zh{很} (\emph{hěn}) comme \cle{marqueur d'adjectif predicatif}. \zh{很} a ici une valeur grammaticale (pas forcément « très »).
\zh{我觉得篮球很有意思。} \emph{Wǒ juéde lánqiú hěn yǒuyìsi.}
\end{propriete}

\begin{propriete}[F4 — Justifier : \zh{因为}... + structure \zh{对...重要}]
\textbf{Cause :} \zh{因为} + phrase (souvent en position médiane ou finale). \\
\textbf{Importance :} \textbf{A \zh{对} B \zh{很} + Adj.} = « A est très [Adj] pour B ». \zh{对} (\emph{duì}) est ici \textbf{préposition} (« pour, envers »), \emph{pas} l'adjectif « correct ».
\zh{运动对健康很重要。} \emph{Yùndòng duì jiànkāng hěn zhòngyào.} \\
\zh{运动对身体很好。} \emph{Yùndòng duì shēntǐ hěn hǎo.}
\end{propriete}

\begin{exemplebox}[Variantes pour l'oral (substitution libre)]
\zh{我觉得游泳很有意思。} \emph{Wǒ juéde yóuyǒng hěn yǒuyìsi.} « Je trouve la natation intéressante. »\\[3pt]
\zh{我觉得跑步很好。} \emph{Wǒ juéde pǎobù hěn hǎo.} « Je trouve la course bien. »\\[3pt]
\zh{我每天锻炼，因为运动对身体很重要。} \emph{Wǒ měitiān duànliàn, yīnwèi yùndòng duì shēntǐ hěn zhòngyào.} « Je m'entraîne chaque jour car le sport est important pour le corps. »\\[3pt]
\zh{你常常打乒乓球吗？} \emph{Nǐ chángcháng dǎ pīngpāngqiú ma ?} « Tu joues souvent au ping-pong ? »
\end{exemplebox}

\begin{center}
\begin{tabularx}{\linewidth}{|X|X|X|}
\hline
\rowcolor{popblueL} Fonction communicative & Structure type (caractères + pinyin) & Équivalent français \\
\hline
Saluer & 你好！Nǐ hǎo ! & Salut ! \\
\hline
Demander le sport & 你喜欢什么运动？Nǐ xǐhuan shénme yùndòng ? & Quel sport aimes-tu ? \\
\hline
Répondre (verbe d'action) & 我喜欢踢足球。Wǒ xǐhuan tī zúqiú. & J'aime jouer au foot. \\
\hline
Relancer (économie) & 你呢？Nǐ ne ? & Et toi ? \\
\hline
Donner son avis & 我觉得...很有意思。Wǒ juéde... hěn yǒuyìsi. & Je trouve que... c'est intéressant. \\
\hline
Demander la fréquence & 你常常锻炼吗？Nǐ chángcháng duànliàn ma ? & Tu t'entraînes souvent ? \\
\hline
Répondre fréquence & 我每天都锻炼。Wǒ měitiān dōu duànliàn. & Je m'entraîne tous les jours. \\
\hline
Justifier (cause) & 因为运动对健康很重要。Yīnwèi yùndòng duì jiànkāng hěn zhòngyào. & Parce que le sport est important pour la santé. \\
\hline
\end{tabularx}
\end{center}

\courssub{Contraste Chinois / Français : trois points de vigilance}

\begin{enumerate}
\item \textbf{Ordre des mots (Question) :} FR \emph{« Quel sport aimes-\textbf{tu} ? »} (inversion) vs CN \zh{你喜欢什么运动？} (SVO inchangé, \zh{什么} in situ).
\item \textbf{Articles :} FR \emph{« \textbf{le} football, \textbf{la} santé »} vs CN \zh{足球}, \zh{健康} (zéro article).
\item \textbf{Verbe support selon le sport :} FR \emph{« jouer »} universel vs CN \zh{踢} (pied/foot), \zh{打} (main/raquette/basket), \zh{游} (nage), \zh{跑} (course).
\item \textbf{Position du complément de temps :} FR \emph{« Je m'entraîne \textbf{tous les jours} »} (fin) vs CN \zh{我\textbf{每天都}锻炼} (avant le verbe, \zh{都} renforce la totalité).
\end{enumerate}

\begin{attention}[Top 4 des erreurs francophones à bannir]
\begin{itemize}
\item \textbf{✗} \zh{我喜欢运动足球} \quad \textbf{✓} \zh{我喜欢\textbf{踢}足球} (verbe d'action manquant).
\item \textbf{✗} \zh{和你？} ou \zh{你和？} \quad \textbf{✓} \zh{\textbf{你呢？}} (seule relance correcte).
\item \textbf{✗} \zh{我觉得篮球有意思} \quad \textbf{✓} \zh{我觉得篮球\textbf{很}有意思} (\zh{很} obligatoire à l'affirmatif).
\item \textbf{✗} \zh{我锻炼每天} \quad \textbf{✓} \zh{我\textbf{每天都}锻炼} (complément de temps \textbf{avant} le verbe).
\end{itemize}
\end{attention}

\courssub{Travail ciblé sur les tons et la prononciation}

La correction des tons n'est pas optionnelle : elle conditionne l'intelligibilité. Trois points critiques dans ce dialogue.

\begin{methode}[Entraînement phonétique en 4 temps (10 min)]
\begin{enumerate}
\item \textbf{Le ton neutre \zh{呢} / \zh{ma} / \zh{de}.} Professeur modèle $\rightarrow$ Classe en chœur (exagérer la brièveté) $\rightarrow$ Rangées $\rightarrow$ Individuel. Contrôle : le ton neutre ne doit pas « porter » plus que la syllabe pleine précédente.
\item \textbf{Le 3e ton \zh{很} (hěn) + 3e ton \zh{有} (yǒu) dans \zh{很有意思}.} Règle du \textbf{3e ton + 3e ton} : le premier devient 2e ton montant. \rightarrow Prononcer \emph{hényǒuyìsi} (pas \emph{hěnyǒu...}). Exercice : \zh{很好} \emph{hén hǎo}, \zh{很忙} \emph{hén máng}, \zh{很有意思} \emph{hén yǒuyìsi}.
\item \textbf{L'intonation de la question \zh{吗} / \zh{呢}.} Montée finale nette sur la dernière syllabe (\emph{ma?}, \emph{ne?}), sans crier. Flèche montante \tikz[baseline=-0.5ex]\draw[poporange, thick, ->] (0,0) -- (0.5,0.5); tracée au crayon sur le texte.
\item \textbf{L'enchaînement \zh{每天都} (měitiān dōu).} \zh{每} (3e) + \zh{天} (1e) + \zh{都} (1e). Attention au 3e ton isolé \zh{每} (bas-montant) qui devient souvent demi-3e ton (bas) en phrase rapide. Entraînement : \emph{měi-tiān-dōu} fluide, sans pause.
\end{enumerate}
\end{methode}

\begin{experience}[Atelier prononciation : « Le miroir tonal »]
En binôme. L'élève A lit une réplique du dialogue \textbf{en masquant le pinyin} (regarde seulement les caractères). L'élève B, qui a le pinyin devant lui, note sur une fiche : (1) Tons neutres respectés ? (2) Sandhi \zh{很有} correct ? (3) Intonation finale ? Inversion des rôles. L'enseignant circule et corrige les 3 points ciblés uniquement.
\end{experience}

\courssub{Dialogue modèle 2 : Mini-débat « Faut-il faire du sport tous les jours ? »}

On passe de l'échange personnel (goûts) au \cle{débat d'idées} (argumentation). Nouveau dialogue entre \zh{王老师} (Professeur Wang) et deux élèves. Nouvelles structures : \zh{应该} (devoir/conseil), \zh{最好} (il vaut mieux), \zh{但是} (concession), \zh{也} (aussi).

\begin{textebox}[Dialogue 2 : 每天锻炼，好不好？ — Débat guidé]
王老师：同学们，你们觉得每天锻炼好不好？\\
\emph{Wáng lǎoshī: Tóngxuémen, nǐmen juéde měitiān duànliàn hǎo bu hǎo ?}\\
« Élèves, selon vous, s'entraîner tous les jours, est-ce bien ou pas ? »\\[4pt]
阿明：我觉得每天锻炼很好，因为运动对健康很重要。\\
\emph{Āmíng: Wǒ juéde měitiān duànliàn hěn hǎo, yīnwèi yùndòng duì jiànkāng hěn zhòngyào.}\\
« Je trouve que c'est très bien, car le sport est important pour la santé. »\\[4pt]
小红：我同意。学生应该每天运动，最好锻炼三十分钟。\\
\emph{Xiǎohóng: Wǒ tóngyì. Xuésheng yīnggāi měitiān yùndòng, zuìhǎo duànliàn sānshí fēnzhōng.}\\
« Je suis d'accord. Les élèves devraient faire du sport chaque jour, l'idéal c'est 30 minutes. »\\[4pt]
阿明：但是，作业很多，没有时间。\\
\emph{Āmíng: Dànshì, zuòyè hěn duō, méiyǒu shíjiān.}\\
« Mais il y a beaucoup de devoirs, pas le temps. »\\[4pt]
王老师：没关系。跑步、打乒乓球也可以，不需要很多时间。\\
\emph{Wáng lǎoshī: Méiguānxi. Pǎobù, dǎ pīngpāngqiú yě kěyǐ, bù xūyào hěn duō shíjiān.}\\
« Pas de problème. Courir, jouer au ping-pong ça marche aussi, ça ne demande pas beaucoup de temps. »
\end{textebox}

\begin{propriete}[Nouvelles structures pour le débat]
\begin{enumerate}
\item \textbf{Accord / Désaccord :} \zh{我同意} (Wǒ tóngyì, Je suis d'accord) / \zh{我不同意} (Wǒ bù tóngyì, Je ne suis pas d'accord).
\item \textbf{Conseil / Obligation morale :} \zh{应该} (\emph{yīnggāi}) + Verbe. \zh{学生应该每天运动。} « Les élèves devraient faire du sport chaque jour. »
\item \textbf{Recommandation optimale :} \zh{最好} (\emph{zuìhǎo}) + Verbe. \zh{最好锻炼三十分钟。} « Il vaudrait mieux s'entraîner 30 min. »
\item \textbf{Concession / Opposition :} \zh{但是} (\emph{dànshì}) en tête de réplique. \zh{但是作业很多。} « Mais les devoirs sont nombreux. »
\item \textbf{Alternative / Addition :} \zh{也} (\emph{yě}) placé \textbf{avant} le verbe. \zh{打乒乓球也可以。} « Jouer au ping-pong est aussi possible. »
\item \textbf{Négation de nécessité :} \zh{不需要} (\emph{bù xūyào}) + Verbe/Objet. \zh{不需要很多时间。} « N'a pas besoin de beaucoup de temps. »
\end{enumerate}
\end{propriete}

\begin{aretenir}[Boîte à outils du débat oral (Niveau HSK 3)]
\begin{center}
\begin{tabularx}{\linewidth}{|X|X|X|}
\hline
\rowcolor{poppurpleL} Fonction & Formule chinoise + pinyin & Français \\
\hline
Lancer le débat & 你们觉得...好不好？Nǐmen juéde... hǎo bu hǎo ? & Qu'en pensez-vous, bien ou pas ? \\
\hline
Donner son avis & 我觉得...很... / 我认为... & Je trouve que... / Je pense que... \\
\hline
Exprimer accord & 我同意。Wǒ tóngyì. / 对，没错。Duì, méicuò. & Je suis d'accord. / Oui, exact. \\
\hline
Exprimer désaccord & 我不同意。Wǒ bù tóngyì. / 不一定吧。Bù yídìng ba. & Je ne suis pas d'accord. / Pas sûr. \\
\hline
Argumenter (cause) & 因为...所以... Yīnwèi... suǒyǐ... & Parce que... donc... \\
\hline
Argumenter (importance) & ...对...很重要。...duì... hěn zhòngyào. & ... est important pour... \\
\hline
Conseiller & 应该... Yīnggāi... / 最好... Zuìhǎo... & Il faut... / Il vaut mieux... \\
\hline
Nuancer (mais) & 但是... Dànshì... & Mais... \\
\hline
Proposer alternative & ...也可以。... yě kěyǐ. & ... marche aussi / est possible. \\
\hline
Clôturer & 总之，运动很重要。Zǒngzhī, yùndòng hěn zhòngyào. & En résumé, le sport est important. \\
\hline
\end{tabularx}
\end{center}
\end{aretenir}

\courssub{Jeux de rôle et débat guidé en classe (Production orale finale)}

C'est l'évaluation formative de la leçon. Trois activités progressives : substitution guidée $\rightarrow$ jeu de rôle à cartes $\rightarrow$ débat ouvert.

\begin{experience}[Activité 1 : Substitution guidée « Mon sport à moi » (10 min)]
\textbf{Consigne :} En binôme. Reprenez le Dialogue 1. Remplacez les sports et les adjectifs par vos vrais goûts. Utilisez \textbf{obligatoirement} le verbe d'action correct.
\textbf{Support :} Fiche « Menu Sport » au tableau :
\begin{itemize}
\item \zh{踢足球} (tī zúqiú) — \zh{打篮球} (dǎ lánqiú) — \zh{打乒乓球} (dǎ pīngpāngqiú) — \zh{打羽毛球} (dǎ yǔmáoqiú)
\item \zh{游泳} (yóuyǒng) — \zh{跑步} (pǎobù) — \zh{打太极拳} (dǎ tàijíquán, Tai-chi)
\end{itemize}
Adjectifs : \zh{有意思} (intéressant), \zh{累} (léi, fatigant), \zh{健康} (sain), \zh{快乐} (kuàilè, joyeux), \zh{难} (nán, difficile).
\textbf{Produit attendu :} Chaque binôme présente sa version personnalisée devant la classe (1 min). Critères : tons, verbe d'action, \zh{很}+adj, \zh{你呢?} ton neutre.
\end{experience}

\begin{experience}[Activité 2 : Jeu de rôle « Au club de sport » (15 min)]
\textbf{Mise en situation :} Vous êtes à \zh{雅温得体育馆} (Stade de Yaoundé) ou \zh{北京大学体育馆}. Trois rôles par groupe.
\begin{center}
\begin{tabularx}{\linewidth}{|X|X|}
\hline
\rowcolor{poporangeL} Rôle & Objectif langagier \\
\hline
\textbf{A : Nouveau membre} & Se présenter, dire son sport préféré (\zh{我喜欢...}), demander conseil (\zh{你觉得...怎么样？}). \\
\hline
\textbf{B : Coach / Ancien} & Répondre, donner avis (\zh{我觉得...}), conseiller (\zh{你应该...}, \zh{最好...}), justifier (\zh{因为...}). \\
\hline
\textbf{C : Observateur} & Grille d'évaluation : (1) Verbes d'action corrects ? (2) \zh{呢/吗} tons neutres ? (3) \zh{很}+adj ? (4) \zh{因为/所以} ? (5) Fluidité. \\
\hline
\end{tabularx}
\end{center}
\textbf{Déroulement :} 3 min préparation (notes autorisées en pinyin seulement) $\rightarrow$ 5 min jeu $\rightarrow$ 2 min feedback observateur $\rightarrow$ Rotation des rôles (3 tours). L'enseignant note les productions pour correction collective finale.
\end{experience}

\begin{experience}[Activité 3 : Débat structuré « Le sport tous les jours : pour ou contre ? » (15 min)]
\textbf{Thèse :} \zh{学生应该每天锻炼。} (Les élèves devraient s'entraîner tous les jours.)
\textbf{Organisation :} Deux camps (Pour / Contre) + Jury (3 élèves). Temps de préparation 5 min (recherche d'arguments sur brouillon en chinois).
\textbf{Arguments suggérés (affichés au tableau en pinyin/caractères) :}
\begin{itemize}
\item \textbf{Pour :} \zh{对健康重要} (important pour la santé), \zh{成绩会变好} (les notes s'améliorent), \zh{缓解压力} (huǎnjiě yālì, soulager le stress), \zh{交朋友} (jiāo péngyou, se faire des amis).
\item \textbf{Contre :} \zh{作业太多} (devoirs trop nombreux), \zh{没有时间} (pas de temps), \zh{太累} (trop fatiguant), \zh{受伤危险} (shòushāng wéixiǎn, risque de blessure), \zh{天气热} (tiānqì rè, il fait chaud — contexte Cameroun).
\end{itemize}
\textbf{Règles :} 1 argument = 1 phrase complète avec \zh{因为/所以} ou \zh{应该/最好}. Jury note : justesse grammaticale + richesse vocabulaire + prononciation. Déclaration du verdict en chinois : \zh{正方胜} / \zh{反方胜} / \zh{平局}.
\end{experience}

\begin{savaistu}[Sport : Cameroun vs Chine — Regards croisés]
\begin{itemize}
\item \textbf{Football (\zh{足球} zúqiú) :} Religion nationale au Cameroun (\zh{喀麦隆} Kāmàilóng) : \zh{不可触摸的狮子} (Bùkě chùmō de shīzi, Lions Indomptables), Coupe d'Afrique, Samuel Eto'o. En Chine : suivi énorme de la Premier League et de la Ligue 1, développement de la \zh{中超联赛} (Zhōngchāo Liánsài, Chinese Super League), objectif Coupe du Monde.
\item \textbf{Sports « nationaux » :} Chine : \zh{乒乓球} (pīngpāngqiú, ping-pong — sport n°1 historique), \zh{羽毛球} (yǔmáoqiú, badminton), \zh{太极拳} (tàijíquán, Tai-chi — santé/seniors), \zh{跳广场舞} (tiào guǎngchǎngwǔ, danse de place — seniors femmes). Cameroun : \zh{篮球} (lánqiú, basket — très populaire en milieu scolaire/urbain), \zh{手球} (shǒuqiú, handball), \zh{lutte traditionnelle} (rurale), \zh{course à pied} (montagne/cross).
\item \textbf{École :} Chine : \zh{广播体操} (guǎngbō tǐcāo, gymnastique matinale quotidienne obligatoire), 2h sport/semaine + \zh{阳光体育} (yángguāng tǐyù, « sport au soleil » 1h/jour). Cameroun : EPS 2h/semaine (souvent football/basket), inter-classes, \zh{Jeux Universitaires} (JU), \zh{FENASCO} (fédérations scolaires).
\item \textbf{Vocabulaire culturel :} \zh{奥运会} (Àoyùnhuì, JO), \zh{全运会} (Quányùnhuì, Jeux Nationaux Chine), \zh{金牌} (jīnpái, médaille d'or), \zh{国球} (guóqiú, sport national — ping-pong pour la Chine).
\end{itemize}
\end{savaistu}

\courssub{Exercices de consolidation et évaluation}

\begin{exoresolu}[Reconstitution de dialogue \& Traduction]
Énoncé. Remettez les répliques dans l'ordre logique du Dialogue 1. Traduisez la réplique (3).
\begin{enumerate}
\item[(a)] 我喜欢打篮球。我觉得篮球很有意思。
\item[(b)] 对，我每天都锻炼，因为运动对健康很重要。
\item[(c)] 你喜欢什么运动？
\item[(d)] 你常常锻炼吗？
\item[(e)] 我喜欢踢足球，你呢？
\end{enumerate}
\tcblower
\textbf{Solution.}
Ordre : \textbf{(c) $\rightarrow$ (e) $\rightarrow$ (a) $\rightarrow$ (d) $\rightarrow$ (b)}.
Logique : Question initiale (c) $\rightarrow$ Réponse + Relance (e) $\rightarrow$ Réponse + Avis (a) $\rightarrow$ Question fréquence (d) $\rightarrow$ Justification (b).
Traduction (c) : \zh{你喜欢什么运动？} \emph{Nǐ xǐhuan shénme yùndòng ?} « Quel sport aimes-tu ? » — \zh{什么} occupe la place de l'objet (le sport), exactement où viendra se loger la réponse (\zh{踢足球}, \zh{打篮球}...).
\end{exoresolu}

\begin{exercice}[Expression écrite guidée : « Mon journal de sport »]
Rédigez un court paragraphe (5-6 phrases) en caractères chinois (+ pinyin si besoin) pour vous présenter et parler de votre pratique sportive. Utilisez \textbf{au moins} : 2 sports avec verbes d'action, \zh{觉得} + \zh{很}+adj, \zh{每天/常常} + \zh{都}, \zh{因为}... \zh{对...很重要}.
\textbf{Amorce :} \zh{我叫... 我喜欢... 和... 我觉得... 我每天/常常... 因为...}
\end{exercice}

\begin{corrige}[Exercice « Mon journal de sport »]
\textbf{Modèle de production attendue (niveau HSK 2/3) :}
\zh{我叫阿明。我喜欢踢足球和打篮球。我觉得踢足球很有意思，打篮球也很快乐。我常常锻炼，每周三次。因为运动对身体很重要，对学习也很好。}
\emph{Wǒ jiào Āmíng. Wǒ xǐhuan tī zúqiú hé dǎ lánqiú. Wǒ juéde tī zúqiú hěn yǒuyìsi, dǎ lánqiú yě hěn kuàilè. Wǒ chángcháng duànliàn, měi zhōu sān cì. Yīnwèi yùndòng duì shēntǐ hěn zhòngyào, duì xuéxí yě hěn hǎo.}
« Je m'appelle Amin. J'aime jouer au foot et au basket. Je trouve le foot intéressant, le basket aussi joyeux. Je m'entraîne souvent, trois fois par semaine. Parce que le sport est important pour le corps, et bon pour les études aussi. »
\textbf{Points de vigilance corrigés :} \zh{和} (hé) pour lier deux VO ; \zh{也} (yě) avant l'adjectif pour « aussi » ; \zh{每周三次} (měi zhōu sān cì, 3 fois/semaine) structure Fréquence ; \zh{对...好} structure importance.
\end{corrige}

\begin{exercice}[Traduction Thème (Français $\rightarrow$ Chinois) — Évaluation sommative]
Traduisez en chinois (caractères + pinyin).
\begin{enumerate}
\item Quel sport aimes-tu ? — J'aime nager. Et toi ?
\item Je trouve que le ping-pong est très intéressant.
\item Tu t'entraînes souvent ? — Oui, je cours tous les jours.
\item Parce que le sport est important pour la santé.
\item Les élèves devraient faire du sport 30 minutes chaque jour.
\item Mais j'ai beaucoup de devoirs, je n'ai pas le temps.
\item Courir ne demande pas beaucoup de temps, c'est aussi possible.
\end{enumerate}
\end{exercice}

\begin{corrige}[Exercice de Thème]
\begin{enumerate}
\item \zh{你喜欢什么运动？— 我喜欢游泳，你呢？} \emph{Nǐ xǐhuan shénme yùndòng? — Wǒ xǐhuan yóuyǒng, nǐ ne?}
\item \zh{我觉得打乒乓球很有意思。} \emph{Wǒ juéde dǎ pīngpāngqiú hěn yǒuyìsi.}
\item \zh{你常常锻炼吗？— 对，我每天都跑步。} \emph{Nǐ chángcháng duànliàn ma? — Duì, wǒ měitiān dōu pǎobù.}
\item \zh{因为运动对健康很重要。} \emph{Yīnwèi yùndòng duì jiànkāng hěn zhòngyào.}
\item \zh{学生应该每天锻炼三十分钟。} \emph{Xuésheng yīnggāi měitiān duànliàn sānshí fēnzhōng.}
\item \zh{但是作业很多，没有时间。} \emph{Dànshì zuòyè hěn duō, méiyǒu shíjiān.}
\item \zh{跑步不需要很多时间，也可以。} \emph{Pǎobù bù xūyào hěn duō shíjiān, yě kěyǐ.}
\end{enumerate}
\end{corrige}

\begin{aretenir}[Bilan de la Leçon 11 — Ce que tu dois savoir faire]
\begin{enumerate}
\item \textbf{Vocabulaire :} Citer 8 sports avec le \textbf{verbe d'action correct} (\zh{踢/打/游/跑}).
\item \textbf{Grammaire :} Construire une question avec \zh{什么} (in situ) ; relancer avec \zh{呢} (ton neutre) ; donner avis avec \zh{觉得...很...} ; justifier avec \zh{因为...对...很重要} ; conseiller avec \zh{应该/最好} ; nuancer avec \zh{但是/也}.
\item \textbf{Prononciation :} Maîtriser le ton neutre (\zh{呢, 吗, 的}) ; appliquer le sandhi 3e+3e ton (\zh{很有} $\rightarrow$ \emph{hén yǒu}) ; placer l'intonation montante finale sur les questions.
\item \textbf{Oral :} Tenir un échange de 5 répliques (Dialogue 1) ; participer à un débat de 3 arguments (Dialogue 2) ; improviser un jeu de rôle « Coach / Membre ».
\item \textbf{Culture :} Citer 2 différences / similitudes sport Cameroun-Chine (football, ping-pong, EPS, gymnastique matinale).
\end{enumerate}
\end{aretenir}

\coursec{Leçon 11 — Expression orale : importance de la pratique du sport}

\courssub{Rappel \& Mise en route : vocabulaire du sport et dialogue modèle 1}
\begin{experience}[Révision active — Vocabulaire sport \& santé (HSK 2-3)]
En binôme, l'un mime un sport, l'autre devine en chinois. Utilisez la liste ci-dessous (révision Module 2 + extensions). L'enseignant corrige les tons sur les mots-clés.
\begin{center}
\begin{tabularx}{\linewidth}{|X|X|X|X|}
\hline
\rowcolor{popblueL} Caractères & Pinyin & Nature & Traduction \\ \hline
跑步 & pǎobù & verbe & courir, faire du jogging \\ \hline
游泳 & yóuyǒng & verbe & nager \\ \hline
打球 & dǎ qiú & verbe (VO) & jouer au ballon (basket, foot, etc.) \\ \hline
打篮球 & dǎ lánqiú & verbe (VO) & jouer au basket \\ \hline
踢足球 & tī zúqiú & verbe (VO) & jouer au football \\ \hline
锻炼 & duànliàn & verbe / nom & s'exercer, exercice physique \\ \hline
健身 & jiànshēn & verbe / nom & fitness, musculation \\ \hline
散步 & sànbù & verbe & se promener, marcher \\ \hline
瑜伽 & yújiā & nom & yoga \\ \hline
健康 & jiànkāng & adj & en bonne santé \\ \hline
强身健体 & qiángshēn jiàntǐ & locution & renforcer le corps (chengyu courant) \\ \hline
缓解压力 & huǐjiě yālì & verbe + obj & soulager le stress \\ \hline
肥胖 & fèipàng & adj / nom & obésité, obèse \\ \hline
坚持 & jiānchí & verbe & persévérer, persister \\ \hline
习惯 & xíguàn & nom & habitude \\ \hline
\end{tabularx}
\end{center}
\end{experience}

\begin{textebox}[Dialogue modèle 1 : Deux camarades parlent de leurs sports — Révision]
\zh{小王：} 老陈，你平时喜欢做什么运动？\\
\emph{Lǎo Chén, nǐ píngshí xǐhuan zuò shénme yùndòng?}\\
« Lao Chen, tu aimes faire quel sport d'habitude ? »\\

\zh{老陈：} 我喜欢跑步，每天早上跑三十分钟。你呢？\\
\emph{Wǒ xǐhuan pǎobù, měitiān zǎoshang pǎo sānshí fēnzhōng. Nǐ ne?}\\
« J'aime courir, tous les matins je cours 30 minutes. Et toi ? »\\

\zh{小王：} 我比较喜欢游泳，因为游泳能锻炼全身，而且不伤膝盖。\\
\emph{Wǒ bǐjiào xǐhuan yóuyǒng, yīnwèi yóuyǒng néng duànliàn quánshēn, érqiě bù shāng xīgài.}\\
« Moi je préfère la natation, car la natation muscle tout le corps, et en plus n'abîme pas les genoux. »\\

\zh{老陈：} 说得好！我也想学游泳，下次你能教我吗？\\
\emph{Shuō de hǎo! Wǒ yě xiǎng xué yóuyǒng, xiàcì nǐ néng jiào wǒ ma?}\\
« Bien dit ! Moi aussi je veux apprendre à nager, la prochaine fois tu peux m'enseigner ? »\\

\zh{小王：} 没问题，我们周末一起去游泳馆吧！\\
\emph{Méi wèntí, wǒmen zhōumò yīqǐ qù yóuyǒngguǎn ba!}\\
« Pas de problème, on va ensemble à la piscine le week-end, d'accord ? »
\end{textebox}

\coursec{Formules fonctionnelles : présenter, opiner, relancer}

Après avoir travaillé le \emph{Dialogue modèle 1} où deux camarades échangeaient sur leurs sports préférés, nous passons maintenant à l'outillage linguistique indispensable pour structurer un exposé oral et animer un débat. En classe de Première, l'examen oral (épreuve de LV2) attend de vous que vous sachiez \textbf{présenter un sujet}, \textbf{développer des arguments}, \textbf{réagir aux interventions des autres} et \textbf{conclure} de façon cohérente. Cette section vous donne les « briques » prêtes à l'emploi, avec leur structure grammaticale, leurs nuances de registre et les pièges à éviter pour un francophone.

\courssub{Présenter son sujet : les formules d'introduction}

Avant de développer, il faut annoncer clairement le thème. En chinois, l'introduction suit souvent le schéma \textbf{Sujet + Verbe d'intention + Objet (sujet du discours)}.

\begin{textebox}[Modèles d'introduction pour un exposé sur le sport]
今天我要谈一谈运动的重要性。\\
Jīntiān wǒ yào tán yi tán yùndòng de zhòngyàoxìng.\\
Aujourd'hui, je vais parler de l'importance du sport.\\

我想说说体育锻炼的好处。\\
Wǒ xiǎng shuōshuo tǐyù duànliàn de hǎochù.\\
Je voudrais parler des bienfaits de l'exercice physique.\\

这个话题很有意义，我们一起来讨论吧。\\
Zhège huàtí hěn yǒu yìyì, wǒmen yīqǐ lái tǎolùn ba.\\
Ce sujet est très significatif, discutons-en ensemble.
\end{textebox}

\begin{propriete}[Structure des verbes d'intention : 想 / 要 / 打算 + V]
\begin{itemize}
    \item \textbf{想 xiǎng} : « vouloir, avoir l'intention de » — registre neutre, très courant à l'oral.
    \item \textbf{要 yào} : « devoir, aller devoir, avoir l'intention ferme » — plus décidé, marque une volonté claire.
    \item \textbf{打算 dǎsuàn} : « prévoir, compter faire » — insiste sur le projet, le plan.
\end{itemize}
\textbf{Schéma} : Sujet + (想 / 要 / 打算) + Verbe (+ Objet).\\
\textbf{Exemples} :
\begin{itemize}
    \item \zh{我想介绍一下我的运动习惯。} \emph{Wǒ xiǎng jièshào yīxià wǒ de yùndòng xíguàn.} (Je veux présenter mes habitudes sportives.)
    \item \zh{我们打算讨论“每天锻炼”的利弊。} \emph{Wǒmen dǎsuàn tǎolùn “měitiān duànliàn” de lìbì.} (Nous comptons discuter les pour et les contre du « sport quotidien ».)
\end{itemize}
\end{propriete}

\begin{aretenir}[Vocabulaire clé : introduire un exposé]
\begin{center}
\begin{tabularx}{\linewidth}{|X|X|X|X|}
\hline
\rowcolor{popblueL} Caractères & Pinyin & Nature & Traduction \\ \hline
谈一谈 & tán yi tán & verbe (redupl.) & parler un peu de, évoquer \\ \hline
说说 & shuōshuo & verbe (redupl.) & parler de, discuter brièvement \\ \hline
介绍 & jièshào & verbe & présenter, introduire \\ \hline
重要性 & zhòngyàoxìng & nom & importance \\ \hline
好处 & hǎochù & nom & avantage, bienfait \\ \hline
利弊 & lìbì & nom & pour et contre, avantages et inconvénients \\ \hline
话题 & huàtí & nom & sujet de conversation, thème \\ \hline
意义 & yìyì & nom & sens, signification, importance \\ \hline
\end{tabularx}
\end{center}
\end{aretenir}

\courssub{Donner son avis et le justifier : 觉得 vs 认为, 因为…所以…}

Une fois le sujet lancé, il faut exprimer son opinion personnelle. Deux verbes dominent : \textbf{觉得 juéde} (ressenti, impression) et \textbf{认为 rènwéi} (jugement réfléchi, opinion argumentée).

\begin{exemplebox}[觉得 vs 认为 — nuances de registre]
\begin{itemize}
    \item \zh{我觉得游泳很有趣。} \emph{Wǒ juéde yóuyǒng hěn yǒuqù.} — « Je trouve que la natation est amusante. » (Impression personnelle, subjectif, courant entre amis.)
    \item \zh{我认为游泳能锻炼全身。} \emph{Wǒ rènwéi yóuyǒng néng duànliàn quánshēn.} — « Je pense / j'estime que la natation permet de muscler tout le corps. » (Opinion fondée, plus formel, adapté à un exposé.)
    \item \zh{对我们来说，健康最重要。} \emph{Duì wǒmen láishuō, jiànkāng zuì zhòngyào.} — « Pour nous, la santé est primordiale. » (Marque le point de vue du locuteur, très utile en débat.)
    \item \zh{我的看法是：运动要坚持。} \emph{Wǒ de kànfǎ shì: yùndòng yào jiānchí.} — « Mon avis est : il faut persévérer dans le sport. » (Formule solennelle pour poser sa thèse.)
\end{itemize}
\end{exemplebox}

\begin{propriete}[Justifier avec 因为…所以… yīnwèi… suǒyǐ…]
La structure \textbf{因为…所以…} (« parce que… donc… ») est la colonne vertébrale de l'argumentation en chinois. Contrairement au français où « parce que » et « donc » sont souvent redondants (on évite « parce que… donc » à l'écrit), \textbf{en chinois les deux termes coexistent obligatoirement} dans la phrase complexe pour marquer la relation de cause à effet.
\begin{center}
\begin{tabularx}{\linewidth}{|X|X|X|}
\hline
\rowcolor{popblueL} Structure & Exemple (caractères + pinyin) & Traduction \\ \hline
因为 + Cause, 所以 + Conséquence & \zh{因为运动能强身健体，所以我每天跑步。} \emph{Yīnwèi yùndòng néng qiángshēn jiàntǐ, suǒyǐ wǒ měitiān pǎobù.} & Comme le sport renforce le corps, je cours tous les jours. \\ \hline
Cause en tête, conséquence après & \zh{因为下雨，所以比赛取消了。} \emph{Yīnwèi xiàyǔ, suǒyǐ bǐsài qǔxiāo le.} & Parce qu'il a plu, le match a été annulé. \\ \hline
Accent sur la cause : 是…因为… & \zh{我每天跑步，是因为我想保持健康。} \emph{Wǒ měitiān pǎobù, shì yīnwèi wǒ xiǎng bǎochí jiànkāng.} & Je cours tous les jours, \emph{c'est parce que} je veux rester en bonne santé. \\ \hline
\end{tabularx}
\end{center}
\textbf{Attention} : Ne dites jamais \zh{因为…，…} sans \zh{所以} (sauf style télégraphique) ni \zh{…，所以…} sans \zh{因为} (sauf réponse à une question « pourquoi »).
\end{propriete}

\begin{exemplebox}[Arguments types pour « l'importance du sport »]
\begin{itemize}
    \item \zh{因为运动可以让我们身体健康，所以我很重视锻炼。} \emph{Yīnwèi yùndòng kěyǐ ràng wǒmen shēntǐ jiànkāng, suǒyǐ wǒ hěn zhòngshì duànliàn.} (Parce que le sport peut nous rendre en bonne santé, j'accorde beaucoup d'importance à l'exercice.)
    \item \zh{因为缺乏运动会导致肥胖，所以我们应该多动。} \emph{Yīnwèi quēfá yùndòng huì dǎozhì fèipàng, suǒyǐ wǒmen yīnggāi duō dòng.} (Parce que le manque de sport entraîne l'obésité, nous devrions bouger davantage.)
    \item \zh{对学生来说，因为学习压力大，所以运动能缓解压力。} \emph{Duì xuésheng lái shuō, yīnwèi xuéxī yālì dà, suǒyǐ yùndòng néng huǐjiě yālì.} (Pour les étudiants, comme la pression des études est forte, le sport permet de la relâcher.)
\end{itemize}
\end{exemplebox}

\courssub{Exprimer l'accord et le désaccord : politesse et précision}

Dans un débat, réagir aux propos de l'autre est aussi important que parler. Le chinois distingue clairement l'accord franc, l'accord partiel et le désaccord courtois.

\begin{propriete}[Accord : 我同意 / 你说得对 / 我也是]
\begin{itemize}
    \item \zh{我同意你的看法。} \emph{Wǒ tóngyì nǐ de kànfǎ.} — « Je suis d'accord avec ton avis. » (Standard, clair.)
    \item \zh{你说得对。} \emph{Nǐ shuō de duì.} — « Tu as raison. » (Reconnaissance de la justesse de l'argument.)
    \item \zh{我也是这么想的。} \emph{Wǒ yě shì zhème xiǎng de.} — « Je pense la même chose. » (Identification totale.)
    \item \zh{完全同意。} \emph{Wánquán tóngyì.} — « Tout à fait d'accord. » (Très bref, oral.)
\end{itemize}
\textbf{Piège majeur} : \textbf{Ne jamais dire \zh{我是同意} (wǒ shì tóngyì)}. En chinois, \zh{同意} est un verbe à part entière, il ne prend pas le copule \zh{是}. C'est une erreur de calque sur « je \emph{suis} d'accord ».
\end{propriete}

\begin{propriete}[Désaccord poli : 对不起，我不同意 / 我觉得不一定]
\begin{itemize}
    \item \zh{对不起，我不同意。} \emph{Duìbuqǐ, wǒ bù tóngyì.} — « Désolé, je ne suis pas d'accord. » (Formule passe-partout, courtoise.)
    \item \zh{我不太赞同。} \emph{Wǒ bù tài zàntóng.} — « Je ne suis pas tout à fait d'accord. » (\zh{赞同} = approuver, plus formel que \zh{同意}.)
    \item \zh{我觉得不一定。} \emph{Wǒ juéde bù yídìng.} — « Je ne pense pas que ce soit forcément le cas. » (Désaccord nuancé, ouvre la discussion.)
    \item \zh{这倒不一定，比如……} \emph{Zhè dào bù yídìng, bǐrú…} — « Ce n'est pas certain, par exemple… » (Introduction d'un contre-exemple.)
\end{itemize}
\end{propriete}

\begin{aretenir}[Tableau récapitulatif : Accord / Désaccord]
\begin{center}
\begin{tabularx}{\linewidth}{|X|X|X|X|}
\hline
\rowcolor{popblueL} Fonction & Chinois (caractères) & Pinyin & Registre / Nuance \\ \hline
Accord total & 我同意 & wǒ tóngyì & Neutre, standard \\ \hline
Accord total & 你说得对 & nǐ shuō de duì & Reconnaît la justesse \\ \hline
Accord total & 我也是这么想的 & wǒ yě shì zhème xiǎng de & Identification personnelle \\ \hline
Accord fort & 完全同意 & wánquán tóngyì & Oral, enthousiaste \\ \hline
Désaccord poli & 对不起，我不同意 & duìbuqǐ, wǒ bù tóngyì & Courtois, standard \\ \hline
Désaccord formel & 我不太赞同 & wǒ bù tài zàntóng & Débat structuré, exposé \\ \hline
Désaccord nuancé & 我觉得不一定 & wǒ juéde bù yídìng & Ouvre la discussion \\ \hline
Contre-exemple & 这倒不一定，比如… & zhè dào bù yídìng, bǐrú… & Argumentatif \\ \hline
\end{tabularx}
\end{center}
\end{aretenir}

\courssub{Relancer l'échange : questions de relance et invitations à développer}

Un débat vivant suppose que chaque participant \emph{relance} l'autre. Les formules ci-dessous sont vos « passe-plats » pour ne pas laisser le dialogue tomber.

\begin{exemplebox}[Formules de relance — du plus simple au plus élaboré]
\begin{itemize}
    \item \zh{你呢？} \emph{Nǐ ne?} — « Et toi ? » (Retour miroir immédiat.)
    \item \zh{你觉得怎么样？} \emph{Nǐ juéde zěnmeyàng?} — « Qu'en penses-tu ? » (Invitation à donner son avis.)
    \item \zh{为什么？} \emph{Wèishénme?} — « Pourquoi ? » (Demande de justification.)
    \item \zh{你能举个例子吗？} \emph{Nǐ néng jǔ gè lìzi ma?} — « Peux-tu donner un exemple ? » (Demande de concrétisation.)
    \item \zh{你能具体说说吗？} \emph{Nǐ néng jùtǐ shuōshuo ma?} — « Peux-tu être plus précis / développer ? » (Invitation à élaborer.)
    \item \zh{还有别的看法吗？} \emph{Hái yǒu bié de kànfǎ ma?} — « Y a-t-il d'autres avis ? » (Ouverture au groupe.)
\end{itemize}
\end{exemplebox}

\begin{propriete}[Particule 吗 ma vs particules de relance 呢 ne / 吧 ba]
\begin{itemize}
    \item \zh{吗} : transforme une phrase déclarative en question fermée (oui/non). Ex. \zh{你喜欢跑步吗？}
    \item \zh{呢} : après un nom/pronom, renvoie la question à l'interlocuteur (« et toi ? »). Ex. \zh{我喜欢游泳，你呢？}
    \item \zh{吧} : suggestion douce, recherche d'accord. Ex. \zh{我们一起去打球吧？} (On va jouer au ballon ensemble, d'accord ?)
\end{itemize}
\end{propriete}

\courssub{Travail sur les tons et la prononciation — Points critiques}

\begin{propriete}[Rappel : Sandhi du 3e ton \& Tons neutres fréquents]
\begin{itemize}
    \item \textbf{3e ton + 3e ton} : le premier devient 2e ton. Ex. \zh{我觉得} \rightarrow \emph{wó juéde} (pas \emph{wǒ juéde}) ; \zh{很好} \rightarrow \emph{hén hǎo}.
    \item \textbf{Tons neutres (qīngshēng)} : particules \zh{了 le, 吗 ma, 呢 ne, 吧 ba, 的 de, 地 de, 得 de} ; suffixes \zh{们 men, 子 zi, 头 tou} ; 2e syllabe de certains mots bisyllabiques \zh{大家 dàjiā, 东西 dōngxi, 时候 shíhou, 一下 yīxià, 起来 qǐlai}.
    \item \textbf{Paires minimales à travailler} (écoute / répétition chorale puis individuelle) :
\end{itemize}
\begin{center}
\begin{tabularx}{\linewidth}{|X|X|X|X|}
\hline
\rowcolor{popblueL} Caractères & Pinyin & Sens & Piège \\ \hline
重点 & zhòngdiǎn & point clé & vs \zh{种点} zhǒngdiǎn (semer) \\ \hline
运动 & yùndòng & sport & vs \zh{云动} yúndòng (mouv. nuages) \\ \hline
锻炼 & duànliàn & exercice & vs \zh{短练} duǎnliàn (court entraînement) \\ \hline
重要 & zhòngyào & important & 3e+4e ton, pas 4e+4e \\ \hline
大家 & dàjiā & tout le monde & 4e + neutre (pas 4e+1e) \\ \hline
一下 & yīxià & un peu / une fois & 1e + neutre \\ \hline
\end{tabularx}
\end{center}
\end{propriete}

\begin{experience}[Exercice de prononciation guidée — « Chaîne de tons »]
L'énonciateur lit une phrase, la classe répète en chœur en exagérant les tons, puis 3 élèves la disent individuellement.
\begin{enumerate}
    \item \zh{因为运动能强身健体，所以我每天跑步。} (Focus : \zh{yīnwèi} 1-4, \zh{suǒyǐ} 3-3 $\rightarrow$ 2-3, \zh{qiángshēn} 2-1, \zh{jiàntǐ} 4-3)
    \item \zh{我觉得不一定，你能举个例子吗？} (Focus : \zh{juéde} 3-2 $\rightarrow$ 2-2, \zh{yídìng} 2-4, \zh{lìzi} 4-neutre, \zh{ma} neutre)
    \item \zh{总之，坚持锻炼，健康生活。谢谢大家！} (Focus : \zh{zǒngzhī} 3-1, \zh{jiānchí} 1-2, \zh{dàjiā} 4-neutre)
\end{enumerate}
\end{experience}

\courssub{Dialogue modèle 2 : Mini-débat « Faut-il faire du sport tous les jours ? »}

\begin{textebox}[Mini-débat modélisé — Pour / Contre]
\zh{主持人：} 今天我们讨论“学生要不要每天运动”。小李，你先说。\\
\emph{Zhǔchírén: Jīntiān wǒmen tǎolùn “xuésheng yào bu yào měitiān yùndòng”. Xiǎo Lǐ, nǐ xiān shuō.}\\
« Aujourd'hui nous discutons "Les étudiants doivent-ils faire du sport tous les jours ?" Xiao Li, vous commencez. »\\

\zh{小李（赞成）：} 我认为学生应该每天运动。因为运动能缓解压力，所以学习效率更高。而且身体健康，就不容易请假。\\
\emph{Xiǎo Lǐ: Wǒ rènwéi xuésheng yīnggāi měitiān yùndòng. Yīnwèi yùndòng néng huǐjiě yālì, suǒyǐ xuéxī xiàolǜ gèng gāo. Érqiě shēntǐ jiànkāng, jiù bù róngyì qǐngjià.}\\
« J'estime que les étudiants devraient faire du sport chaque jour. Car le sport soulage le stress, donc l'efficacité d'étude est plus élevée. De plus, corps sain, donc moins de congés maladie. »\\

\zh{小张（反对）：} 对不起，我不同意“每天都要运动”。因为学生作业太多，所以每天运动不现实。我觉得一周三四次比较合理。\\
\emph{Xiǎo Zhāng: Duìbuqǐ, wǒ bù tóngyì “měitiān dōu yào yùndòng”. Yīnwèi xuésheng zuòyè tài duō, suǒyǐ měitiān yùndòng bù xiànshí. Wǒ juéde yī zhōu sān sì cì bǐjiào hélǐ.}\\
« Désolé, je ne suis pas d'accord avec "faire du sport tous les jours". Car les étudiants ont trop de devoirs, donc le sport quotidien n'est pas réaliste. Je trouve que 3-4 fois par semaine est plus raisonnable. »\\

\zh{小李：} 你的观点有道理。但是我觉得“不现实”可以克服，比如利用课间操。你觉得呢？\\
\emph{Xiǎo Lǐ: Nǐ de guāndiǎn yǒu dàoli. Dànshì wǒ juéde “bù xiànshí” kěyǐ kèfú, bǐrú lìyòng kèjiān cāo. Nǐ juéde ne?}\\
« Ton point de vue a du sens. Mais je pense que "pas réaliste" peut être surmonté, par exemple en profitant de la gymnastique intercours. Qu'en penses-tu ? »\\

\zh{小张：} 课间时间太短，强度不够。还是建议安排专门时间。对吧？\\
\emph{Xiǎo Zhāng: Kèjiān shíjiān tài duǎn, qiángdù bù gòu. Hái shì jiànyì ānpái zhuānmén shíjiān. Duì ba?}\\
« Le temps intercours est trop court, intensité insuffisante. Mieux vaut prévoir un temps dédié. Non ? »\\

\zh{主持人：} 总之，双方都有道理。建议同学们根据自己的课表安排运动。谢谢两位同学！\\
\emph{Zhǔchírén: Zǒngzhī, shuāngfāng dōu yǒu dàoli. Jiànyì tóngxuémen gēnjù zìjǐ de kèbiǎo ānpái yùndòng. Xièxie liǎng wèi tóngxué!}\\
« En résumé, les deux côtés ont raison. Je conseille aux camarades d'organiser le sport selon leur emploi du temps. Merci aux deux camarades ! »
\end{textebox}

\courssub{Conclure un exposé : 所以 / 总之 / 谢谢大家}

La conclusion doit résumer la thèse et remercier l'auditoire. Trois marqueurs sont essentiels.

\begin{exemplebox}[Modèles de conclusion]
\begin{itemize}
    \item \zh{所以，我认为运动对每个人都很重要。谢谢大家！} \emph{Suǒyǐ, wǒ rènwéi yùndòng duì měige rén dōu hěn zhòngyào. Xièxie dàjiā!} (Donc, je pense que le sport est important pour chacun. Merci à tous !)
    \item \zh{总之，坚持锻炼，健康生活。谢谢大家聆听。} \emph{Zǒngzhī, jiānchí duànliàn, jiànkāng shēnghuó. Xièxie dàjiā língtīng.} (En résumé, persévérez dans l'exercice, vivez en santé. Merci à tous pour votre écoute.)
    \item \zh{希望大家都能养成运动的习惯。谢谢！} \emph{Xīwàng dàjiā dōu néng yǎngchéng yùndòng de xíguàn. Xièxie!} (J'espère que tout le monde pourra prendre l'habitude de faire du sport. Merci !)
\end{itemize}
\end{exemplebox}

\begin{aretenir}[Marqueurs de conclusion]
\begin{center}
\begin{tabularx}{\linewidth}{|X|X|X|}
\hline
\rowcolor{popblueL} Marqueur & Pinyin & Emploi \\ \hline
所以 & suǒyǐ & « donc, par conséquent » — lie la conclusion aux arguments (cause→effet) \\ \hline
总之 & zǒngzhī & « en résumé, en un mot » — synthèse globale, style plus formel \\ \hline
希望 & xīwàng & « j'espère que » — ouvre sur un vœu, une recommandation \\ \hline
谢谢大家 / 谢谢大家聆听 & xièxie dàjiā / língtīng & Formules de politesse de clôture obligatoires à l'oral \\ \hline
\end{tabularx}
\end{center}
\end{aretenir}

\courssub{Boîte à outils du débat : organigramme de synthèse}

La figure ci-dessous regroupe les trois familles de formules vues ci-dessus. Utilisez-la comme « fiche mémo » pendant vos entraînements à l'oral.

\begin{popfigure}
\begin{tikzpicture}[
    node distance=1.2cm and 2.5cm,
    box/.style={rectangle, rounded corners=6pt, draw=popblue, thick, fill=popblueL, text width=4.2cm, align=center, minimum height=4.2cm},
    title/.style={rectangle, rounded corners=4pt, fill=popblue, text=white, font=\bfseries\large, minimum height=1cm, text width=4.2cm, align=center},
    arrow/.style={-Stealth, poporange, thick}
]
% Colonnes
\node[title] (t1) {Présenter};
\node[box, below=0pt of t1, align=center] (b1){
    \textbf{今天我要谈一谈…}\\
    \textbf{我想说说…}\\
    \textbf{我打算介绍…}\\
    \textbf{这个话题很有意义}
};

\node[title, right=of t1] (t2) {Donner son avis / Réagir};
\node[box, below=0pt of t2, align=center] (b2){
    \textbf{Opiner :}\\
    我觉得… / 我认为…\\
    对我来说… / 我的看法是…\\
    \textbf{Justifier :}\\
    因为…所以…\\
    \textbf{Accord :}\\
    我同意 / 你说得对\\
    \textbf{Désaccord :}\\
    对不起，我不同意\\
    我觉得不一定
};

\node[title, right=of t2] (t3) {Relancer / Conclure};
\node[box, below=0pt of t3, align=center] (b3){
    \textbf{Relancer :}\\
    你呢？ / 你觉得怎么样？\\
    为什么？ / 能举个例子吗？\\
    \textbf{Conclure :}\\
    所以，我认为…\\
    总之，…\\
    谢谢大家！
};

% Flèches horizontales entre colonnes
\draw[arrow] (b1.east) -- (b2.west) node[midway, above, font=\small, poporange] {Développer};
\draw[arrow] (b2.east) -- (b3.west) node[midway, above, font=\small, poporange] {Échanger / Clôturer};
% Flèche retour (boucle débat)
\draw[arrow, bend left=30] (b3.north east) to node[above right, font=\small, poporange] {Nouvel intervenant} (b1.north west);
\end{tikzpicture}
\legende{Organigramme « Boîte à outils du débat » — trois colonnes : Présenter, Donner son avis / Réagir, Relancer / Conclure. Les flèches indiquent le flux naturel d'un exposé-débat.}
\end{popfigure}

\courssub{Méthode : structurer un mini-exposé oral (2–3 minutes)}

\begin{methode}[Préparer et prononcer un mini-exposé sur le sport]
\begin{enumerate}
    \item \textbf{Accroche (1 phrase)} : Utilisez \zh{今天我要谈一谈…} ou \zh{我想说说…} + le thème précis (ex. \zh{运动对青少年的重要性}).
    \item \textbf{Thèse claire (1 phrase)} : \zh{我认为运动非常必要。} ou \zh{我的看法是：每天锻炼有益健康。}
    \item \textbf{2–3 arguments avec 因为…所以…} : Préparez 2 à 3 paires cause-effet. Ex. \zh{因为运动能增强免疫力，所以我们不容易生病。} / \zh{因为缺乏运动会导致肥胖，所以青少年更要动起来。}
    \item \textbf{Exemple concret (1 phrase)} : \zh{比如，我同学小李每天跑步，身体很棒。} (Par exemple, mon camarade Xiao Li court chaque jour, il est en grande forme.)
    \item \textbf{Conclusion (1 phrase + politesse)} : \zh{所以，我建议大家每天运动半小时。谢谢大家！}
    \item \textbf{Répétition à voix haute} : Travaillez les tons (3e ton, tons neutres), les liaisons, le débit. Enregistrez-vous, corrigez, recommencez.
\end{enumerate}
\end{methode}

\courssub{Pièges fréquents des francophones — Attention !}

\begin{attention}[Erreurs classiques à ne pas commettre]
\begin{itemize}
    \item \textbf{« Je suis d'accord » $\neq$ \zh{我是同意} } — \zh{同意} est un verbe, pas un adjectif. Dites \zh{我同意} ou \zh{我赞同}.
    \item \textbf{« Je pense que » $\neq$ \zh{我想那} } — \zh{想} seul ne prend pas de proposition complétive. Utilisez \zh{我觉得} ou \zh{我认为} + phrase complète.
    \item \textbf{Oublier \zh{所以} dans 因为…所以…} — En français on dit « Parce qu'il pleut, le match est annulé » (pas de « donc »). En chinois, \zh{所以} est \emph{obligatoire} dans la phrase complète.
    \item \textbf{Confondre \zh{觉得} (impression) et \zh{认为} (opinion argumentée)} — Dans un exposé noté, privilégiez \zh{我认为} pour les arguments principaux.
    \item \textbf{Ne pas conclure par \zh{谢谢大家}} — L'omission est perçue comme un manque de politesse (socioculturel fort).
    \item \textbf{Tons : \zh{重点} zhòngdiǎn (point clé) vs \zh{种点} zhǒngdiǎn (semer des points) ; \zh{运动} yùndòng (sport) vs \zh{云动} yúndòng (mouvement des nuages).} Travaillez les paires minimales.
    \item \textbf{Classificateur oublié} : \zh{一个例子} (yí gè lìzi), pas \zh{一例子} ; \zh{三次} (sān cì) pour « trois fois », pas \zh{三个次}.
\end{itemize}
\end{attention}

\courssub{Exemples traduits commentés}

\begin{exemplebox}[Exemples complets : avis + justification + relance]
\begin{itemize}
    \item \zh{对我们来说，健康最重要。因为运动能预防疾病，所以我们要坚持锻炼。你觉得呢？}\\
    \emph{Duì wǒmen láishuō, jiànkāng zuì zhòngyào. Yīnwèi yùndòng néng yùfáng jíbìng, suǒyǐ wǒmen yào jiānchí duànliàn. Nǐ juéde ne?}\\
    « Pour nous, la santé est le plus important. Parce que le sport peut prévenir les maladies, il faut persévérer dans l'exercice. Qu'en penses-tu ? »

    \item \zh{我认为游泳最好，因为它锻炼全身而且不伤膝盖。你能举个例子吗？}\\
    \emph{Wǒ rènwéi yóuyǒng zuì hǎo, yīnwèi tā duànliàn quánshēn érqiě bù shāng xīgài. Nǐ néng jǔ gè lìzi ma?}\\
    « J'estime que la natation est la meilleure, car elle muscle tout le corps et n'abîme pas les genoux. Peux-tu donner un exemple ? »

    \item \zh{你说得对，运动很有用。但是我觉得不一定要每天练，三次一周也可以。对吧？}\\
    \emph{Nǐ shuō de duì, yùndòng hěn yǒuyòng. Dànshì wǒ juéde bù yídìng yào měitiān liàn, sān cì yī zhōu yě kěyǐ. Duì ba?}\\
    « Tu as raison, le sport est très utile. Mais je ne pense pas qu'il faille forcément s'entraîner tous les jours, trois fois par semaine ça va aussi. Non ? »
\end{itemize}
\end{exemplebox}

\courssub{Exercice résolu : transformer un brouillon en exposé structuré}

\begin{exoresolu}[De l'idée brute à l'exposé cohérent]
\textbf{Énoncé} : Voici des notes prises par un élève pour un exposé de 2 minutes sur « Le sport et les études ». Réorganisez-les en un discours oral complet en utilisant les formules de cette leçon (introduction, thèse, 2 arguments \zh{因为…所以…}, exemple, conclusion, politesse). Écrivez le texte en caractères, pinyin et français.

\textbf{Notes de l'élève} :
\begin{itemize}
    \item Sujet : sport et études
    \item Mon avis : sport aide études
    \item Argument 1 : sport → moins stress → mieux concentré
    \item Argument 2 : sport → corps fort → moins malade → pas rater cours
    \item Exemple : moi, je cours 30 min → meilleurs notes
    \item Conclusion : faut faire sport
\end{itemize}

\textbf{Solution détaillée} :

\textbf{Texte chinois (caractères)} :\\
今天我要谈一谈体育锻炼和学习的关系。\\
我认为运动对学习很有帮助。\\
因为运动能缓解压力，所以上课时更容易集中注意力。\\
因为坚持锻炼身体强壮，所以不容易生病，就不会耽误课程。\\
比如我自己，每天跑步三十分钟，成绩明显进步了。\\
所以，我建议同学们都要安排时间运动。谢谢大家！

\textbf{Pinyin} :\\
Jīntiān wǒ yào tán yi tán tǐyù duànliàn hé xuéxí de guānxì.\\
Wǒ rènwéi yùndòng duì xuéxí hěn yǒu bāngzhù.\\
Yīnwèi yùndòng néng huǐjiě yālì, suǒyǐ shàngkè shí gèng róngyì jízhōng zhùyìlì.\\
Yīnwèi jiānchí duànliàn shēntǐ qiángzhuàng, suǒyǐ bù róngyì shēngbìng, jiù bù huì dānwù kèchéng.\\
Bǐrú wǒ zìjǐ, měitiān pǎobù sānshí fēnzhōng, chéngjì míngxiǎn jìnbù le.\\
Suǒyǐ, wǒ jiànyì tóngxuémen dōu yào ānpái shíjiān yùndòng. Xièxie dàjiā!

\textbf{Traduction française} :\\
« Aujourd'hui je vais parler de la relation entre l'exercice physique et les études.\\
J'estime que le sport aide beaucoup les études.\\
Parce que le sport peut soulager le stress, en cours on se concentre plus facilement.\\
Parce que persévérer dans l'exercice rend le corps robuste, on tombe moins malade, donc on ne rate pas les cours.\\
Par exemple moi, je cours 30 minutes chaque jour, mes notes ont nettement progressé.\\
Donc, je conseille aux camarades d'organiser du temps pour le sport. Merci à tous ! »

\textbf{Analyse pédagogique} : L'exposé suit la méthode en 6 étapes. On note l'usage de \zh{今天我要谈一谈…} (intro), \zh{我认为…} (thèse formelle), deux structures \zh{因为…所以…} bien distinctes, \zh{比如…} (exemple personnel), \zh{所以…建议…} (conclusion + recommandation) et \zh{谢谢大家} (clôture polie).
\end{exoresolu}

\courssub{Le saviez-vous ? — Étiquette du débat en Chine}

\begin{savaistu}[Politesse et « face » dans le débat chinois]
En Chine, la notion de \textbf{面子 miànzi} (« face », honneur, dignité sociale) imprègne toute interaction orale. Interrompre brutalement quelqu'un, contredire frontalement sans précaution ou hausser le ton font \emph{perdre la face} à l'interlocuteur — et à soi-même par manque de \zh{修养 xiūyǎng} (éducation, culture de soi).

C'est pourquoi les formules \zh{对不起，我不同意} (« Désolé, je ne suis pas d'accord ») ou \zh{我不太赞同} (« Je ne suis pas tout à fait d'accord ») sont \textbf{obligatoires} dans un débat courtois, même entre amis proches. Elles montrent que vous respectez la parole de l'autre avant d'exprimer la vôtre.

De même, \zh{你说得对} (« Tu as raison ») ou \zh{有道理} (« Il y a du sens / C'est raisonnable ») sont utilisés \emph{avant} d'apporter une nuance : \zh{你说得对，但是我觉得…} (« Tu as raison, mais je pense que… »). Cette structure « accord partiel + pivot + opinion propre » est la marque d'un locuteur cultivé.

En classe, lors des jeux de rôle, entraînez-vous à \textbf{toujours commencer votre réplique par une formule de politesse} (accord ou désaccord poli) avant d'argumenter. C'est un critère d'évaluation implicite à l'oral du baccalauréat.
\end{savaistu}

\begin{exercice}[Entraînement oral : mini-débat « Faut-il faire du sport tous les jours ? »]
Formez des groupes de 4. Chaque groupe tire un rôle : \textbf{Pour} (2 élèves), \textbf{Contre} (2 élèves). Préparez 3 minutes en utilisant \textbf{au moins 3 formules de chaque colonne} de l'organigramme (Présenter, Opiner/Justifier, Relancer/Conclure). Déroulez le débat 5 minutes. L'enseignant note : justesse des tons, fluidité, richesse des connecteurs (\zh{因为…所以…}, \zh{但是}, \zh{而且}), politesse du désaccord. Changez de sujet : « Le sport à l'école doit-il être noté ? » / « Les jeux vidéo sont-ils du sport ? ».
\end{exercice}

\begin{corrige}[Exercice n°1 — Mini-débat]
\textbf{Grille d'auto-évaluation (à remplir après le jeu de rôle)} :\\
\begin{itemize}
    \item [ ] Ai-je utilisé \zh{今天我要谈一谈…} ou équivalent pour ouvrir ?
    \item [ ] Ai-je produit au moins deux \zh{因为…所以…} corrects (avec \zh{所以}) ?
    \item [ ] Ai-je exprimé un désaccord par \zh{对不起，我不同意} ou \zh{我觉得不一定} (pas \zh{我是不同意}) ?
    \item [ ] Ai-je relancé mon partenaire par \zh{你呢？} / \zh{为什么？} / \zh{能举个例子吗？} ?
    \item [ ] Ai-je conclu par \zh{所以…} ou \zh{总之…} + \zh{谢谢大家} ?
    \item [ ] Tons : ai-je vérifié les 3e tons (\zh{我觉得 wǒ juéde} $\rightarrow$ \zh{wó juéde} avant 4e ton) et les tons neutres (\zh{大家 dàjiā}, \zh{一下 yīxià}) ?
\end{itemize}
\textbf{Exemple de réplique « Contre » bien construite} :\\
\zh{对不起，我不同意“每天都要运动”。因为作业太多，所以每天运动不现实。我觉得一周三次比较合理。你觉得呢？}\\
\emph{Duìbuqǐ, wǒ bù tóngyì “měitiān dōu yào yùndòng”. Yīnwèi zuòyè tài duō, suǒyǐ měitiān yùndòng bù xiànshí. Wǒ juéde yī zhōu sān cì bǐjiào hélǐ. Nǐ juéde ne?}
\end{corrige}

\coursec{Travail sur les tons et la prononciation}

Après avoir acquis les formules fonctionnelles pour présenter un exposé et animer un débat, il est indispensable de soigner la prononciation. En chinois, \cle{le ton fait partie du sens} : une même syllabe prononcée avec un ton différent devient un autre mot, parfois sans aucun lien sémantique. Dans un exposé oral sur le sport, confondre \zh{跑} (courir) et \zh{泡} (tremper) ou mal articuler \zh{运动} (sport) peut rendre le message incompréhensible. Cette section propose un entraînement systématique : révision des tons sur le vocabulaire du thème, travail sur les paires minimales, règles de sandhi (modification tonale en contexte) et activités de discrimination et de \emph{shadowing} (répétition en écho).

\courssub{Révision des quatre tons et du ton neutre : vocabulaire du sport et de la santé}

Rappelons les quatre tons de base et le ton neutre (léger, court, sans contour défini) en les ancrant dans le lexique de la leçon. Observez le tableau ci-dessous : chaque mot-clé du thème « sport et santé » illustre un ton. Nous privilégions les mots disyllabiques (vocabulaire HSK 1-3) pour une prononciation plus naturelle.

\begin{center}
\begin{tabularx}{\linewidth}{|X|X|X|X|}
\hline
\rowcolor{popblueL} \textbf{Caractères} & \textbf{Pinyin} & \textbf{Nature} & \textbf{Traduction} \\
\hline
身体 & shēntǐ (1-3) & nom & corps, santé physique \\
\hline
健康 & jiànkāng (4-1) & adj / nom & en bonne santé, santé \\
\hline
运动 & yùndòng (4-4) & nom / verbe & sport, exercice / bouger \\
\hline
锻炼 & duànliàn (4-4) & verbe / nom & exercice physique, s'entraîner \\
\hline
强壮 & qiángzhuàng (2-4) & adjectif & robuste, costaud \\
\hline
练习 & liànxí (4-2) & verbe / nom & pratiquer, exercice, entraînement \\
\hline
重要 & zhòngyào (4-4) & adjectif & important \\
\hline
的 & de (ton neutre) & particule & (possessif / attributif) \\
\hline
\end{tabularx}
\end{center}

\begin{propriete}[Les quatre tons du mandarin standard]
Le chinois mandarin utilise quatre tons pleins (contours de hauteur) et un ton neutre. La hauteur relative se situe dans l'espace vocalique du locuteur (généralement 5 niveaux, de 1 = bas à 5 = haut).
\begin{itemize}
    \item \textbf{1ʳᵉ ton (Yīnpíng 阳平) :} Haut et plat \textbf{(55)}. Exemple : \zh{身 shēn} (dans \zh{身体}). La voix reste haute et stable.
    \item \textbf{2ᵉ ton (Yángpíng 陽平) :} Montant \textbf{(35)}. Exemple : \zh{强 qiáng} (dans \zh{强壮}). La voix monte du milieu vers le haut (comme une question en français : « Quoi ? »).
    \item \textbf{3ᵉ ton (Shǎngshēng 上声) :} Descendant-montant \textbf{(214)}. Exemple : \zh{体 tǐ} (dans \zh{身体}). La voix descend bas puis remonte légèrement. \textbf{Attention :} en pratique courante (sauf en isolation), il se réalise souvent comme un ton « demi-tiers » bas \textbf{(21)} sans la remontée finale.
    \item \textbf{4ᵉ ton (Qùshēng 去声) :} Descendant \textbf{(51)}. Exemple : \zh{健 jiàn}, \zh{动 dòng} (dans \zh{运动}). La voix chute brutalement du haut vers le bas (comme un ordre sec : « Non ! »).
    \item \textbf{Ton neutre (Qīngshēng 轻声) :} Court, léger, sans contour propre. Sa hauteur dépend du ton précédent. Exemple : \zh{的 de} dans \zh{身体的 shēntǐ de}.
\end{itemize}
\end{propriete}

\begin{popfigure}
\begin{tikzpicture}[
    scale=1.2,
    tone/.style={very thick, popblue},
    tone2/.style={very thick, poporange},
    tone3/.style={very thick, popgreen},
    tone4/.style={very thick, poppink},
    axis/.style={thin, popmuted},
    label/.style={font=\small, popdark},
    word/.style={font=\normalsize\bfseries, popdark, anchor=west}
]
% Axes
\draw[axis] (0,0) -- (0,5) node[above, label] {Haut (5)};
\draw[axis] (0,0) -- (6,0) node[right, label] {Temps};
\foreach \y/\label in {1/1, 2/2, 3/3, 4/4, 5/5} {
    \draw[axis, densely dashed] (0,\y) -- (5.5,\y);
    \node[label, left] at (0,\y) {\label};
}

% Ton 1 : Haut plat (55) - 身 shēn
\draw[tone] (0.5,4.5) -- (2.5,4.5);
\node[word] at (2.7,4.5) {1ʳᵉ ton : \zh{身 shēn} (55)};

% Ton 2 : Montant (35) - 强 qiáng
\draw[tone2] (0.5,2) .. controls (1.5,3) and (2.5,4) .. (2.5,4.5);
\node[word] at (2.7,3.2) {2ᵉ ton : \zh{强 qiáng} (35)};

% Ton 3 : Descendant-montant (214) - 体 tǐ
\draw[tone3] (3.5,2) .. controls (4.2,1) and (4.8,1) .. (5.2,2.5);
\node[word] at (5.4,1.5) {3ᵉ ton : \zh{体 tǐ} (214)};

% Ton 4 : Descendant (51) - 动 dòng
\draw[tone4] (3.5,4.5) -- (5.2,1);
\node[word] at (5.4,3) {4ᵉ ton : \zh{动 dòng} (51)};

% Points de repère pour les mots du vocabulaire
\node[label, anchor=east] at (0.3,4.5) {shēn};
\node[label, anchor=east] at (0.3,2) {qiáng};
\node[label, anchor=east] at (3.3,2) {tǐ};
\node[label, anchor=east] at (3.3,4.5) {dòng};
\end{tikzpicture}
\legende{Schéma des contours tonaux des quatre tons avec mots-exemples du thème (1ʳᵉ : \zh{身}, 2ᵉ : \zh{强}, 3ᵉ : \zh{体}, 4ᵉ : \zh{动}). La hauteur 5 est le plus aigu, 1 le plus grave.}
\end{popfigure}

\begin{methode}[Comment travailler un ton isolé puis en phrase]
\begin{enumerate}
    \item \textbf{Écoute active :} L'enseignant (ou l'audio) prononce le mot isolé 3 fois. L'élève visualise la courbe (main levée pour le 1ʳᵉ, main qui monte pour le 2ᵉ, main qui descend puis remonte pour le 3ᵉ, main qui coupe vers le bas pour le 4ᵉ).
    \item \textbf{Exagération lente :} L'élève prononce le mot très lentement en exagérant le contour (monter très haut, descendre très bas). Ex : \zh{shēēēn}, \zh{qiiiiáng}, \zh{tǐǐǐ (bas)}, \zh{jiàààn (sec)}.
    \item \textbf{Vitesse normale isolée :} Prononcer à vitesse naturelle, ton précis.
    \item \textbf{Insertion en phrase :} Dire le mot dans une phrase complète pour sentir l'influence des tons voisins (sandhi). Ex : \zh{我身体很好。 Wǒ shēntǐ hěn hǎo.}
\end{enumerate}
\end{methode}

\courssub{Paires minimales à risque : discrimination auditive}

Les \cle{paires minimales} sont des mots qui ne diffèrent que par un seul trait phonologique (ici le ton). Les francophones ont tendance à ignorer le ton ou à l'aplatir sous l'influence de l'intonation française (montante à la fin de phrase). Voici deux paires critiques pour le thème du sport.

\begin{exemplebox}[Paire minimale 1 : \zh{跑 pǎo} vs \zh{泡 pào}]
\begin{center}
\begin{tabularx}{\linewidth}{|X|X|X|X|}
\hline
\rowcolor{poporangeL} \textbf{Caractère} & \textbf{Pinyin} & \textbf{Nature} & \textbf{Traduction} \\
\hline
跑 & pǎo (3ᵉ) & verbe & courir \\
\hline
泡 & pào (4ᵉ) & verbe & tremper, faire tremper (thé, pieds) \\
\hline
\end{tabularx}
\end{center}
\textbf{Contexte :} \zh{我喜欢跑步。 Wǒ xǐhuan pǎobù.} (J'aime courir.) \\
\textbf{Erreur :} \zh{*我喜欢泡步。 Wǒ xǐhuan pàobù.} (Sans sens : « J'aime tremper l'étape/pied »).
\end{exemplebox}

\begin{attention}[Un mauvais ton change le sens !]
\zh{跑步 pǎobù} = courir (sport). Si vous prononcez \zh{*pāobù} (1ʳᵉ + 4ᵉ) ou \zh{*pàobù} (4ᵉ + 4ᵉ), le mot \zh{跑} n'est pas reconnu. L'interlocuteur entendra \zh{泡 pào} (tremper) ou ne comprendra tout simplement pas. En chinois, \textbf{le ton est aussi important que la voyelle ou la consonne}.
\end{attention}


\begin{exemplebox}[Paire minimale 2 : \zh{练 liàn} vs \zh{连 lián}]
\begin{center}
\begin{tabularx}{\linewidth}{|X|X|X|X|}
\hline
\rowcolor{poporangeL} \textbf{Caractère} & \textbf{Pinyin} & \textbf{Nature} & \textbf{Traduction} \\
\hline
练 & liàn (4ᵉ) & verbe & s'entraîner, pratiquer \\
\hline
连 & lián (2ᵉ) & adv / prép & même, relier, y compris (structure \zh{连...都/也...}) \\
\hline
\end{tabularx}
\end{center}
\textbf{Contexte :} \zh{我每天练习。 Wǒ měitiān liànxí.} (Je m'entraîne tous les jours.) \\
\textbf{Erreur :} \zh{*我每天连习。 Wǒ měitiān liánxí.} (Incorrect grammaticalement, \zh{连} ne s'utilise pas seul ainsi, mais la confusion auditive est forte).
\end{exemplebox}

\begin{attention}[Piège tonal : 4ᵉ ton (court, sec) vs 2ᵉ ton (montant, long)]
\zh{练 liàn} (4ᵉ) chute brutalement : imaginez un coup de sifflet d'arbitre. \zh{连 lián} (2ᵉ) monte comme une question surprise « Ah bon ? ». Pour \zh{锻炼 duànliàn} (exercice physique), les deux syllabes sont en 4ᵉ ton : il faut les « claquer » distinctement sans les traîner.
\end{attention}


\courssub{Le sandhi du troisième ton : deux tons 3 qui se suivent}

C'est la règle de modification tonale la plus fréquente et la plus cruciale pour la fluidité. En chinois, \textbf{deux troisièmes tons ne se suivent jamais à l'oral} : le premier devient un deuxième ton (montant).

\begin{propriete}[Règle du sandhi du 3ᵉ ton (变调 biàndiào)]
\begin{itemize}
    \item \textbf{Structure :} Ton 3 + Ton 3 $\rightarrow$ \textbf{Ton 2 (oral) + Ton 3 (oral)}.
    \item \textbf{Ecriture pinyin :} \textbf{Ne change pas.} On écrit toujours le 3ᵉ ton (caron ǎ, ǐ, ǔ, ǚ, ǒ) pour le premier mot dans le dictionnaire et les manuels.
    \item \textbf{Application :} Valable à l'intérieur d'un mot (ex: \zh{很好 hěn hǎo}) ou à la frontière de deux mots (ex: \zh{我可以 wǒ kěyǐ}).
\end{itemize}
\end{propriete}

\begin{exemplebox}[Applications sur le thème du sport et expressions fréquentes]
\begin{center}
\begin{tabularx}{\linewidth}{|X|X|X|X|}
\hline
\rowcolor{popgreenL} \textbf{Écrit (Pinyin standard)} & \textbf{Prononciation réelle (Sandhi)} & \textbf{Signification} & \textbf{Contexte} \\
\hline
hěn hǎo (很好) & \textbf{hén hǎo} (2+3) & très bien & \zh{我很好。 Wǒ hěn hǎo.} \\
\hline
wǒ kěyǐ (我可以) & \textbf{wó kěyǐ} (2+3) & je peux & \zh{我可以打篮球。 Wǒ kěyǐ dǎ lánqiú.} \\
\hline
nǐ hǎo (你好) & \textbf{ní hǎo} (2+3) & bonjour & \zh{你好！ Nǐ hǎo!} \\
\hline
shuǐguǒ (水果) & \textbf{shuíguǒ} (2+3) & fruit & \zh{吃水果好。 Chī shuǐguǒ hǎo.} \\
\hline
\end{tabularx}
\end{center}
\textbf{Exemples phrases complets :}
\begin{itemize}
    \item \zh{运动很好。 Yùndòng hěn hǎo.} $\rightarrow$ Oral : \zh{Yùndòng hén hǎo.} (Le sport, c'est très bien.)
    \item \zh{我有时间。 Wǒ yǒu shíjiān.} $\rightarrow$ Pas de sandhi (3 + 2).
    \item \zh{我要游泳。 Wǒ yào yóuyǒng.} $\rightarrow$ Pas de sandhi (4 + 2 + 3).
\end{itemize}
\end{exemplebox}

\begin{exoresolu}[Discrimination et application du sandhi]
\textbf{Énoncé :} L'enseignant lit les séquences suivantes. L'élève note le pinyin \textbf{écrit} (standard) et indique la prononciation \textbf{réelle} (avec sandhi si applicable) en cochant la case correspondante.

\begin{enumerate}
    \item \zh{很好} (hěn hǎo) $\rightarrow$ Prononcé : \fbox{hén hǎo} (2+3) \quad \fbox{hěn hǎo} (3+3)
    \item \zh{我可以} (wǒ kěyǐ) $\rightarrow$ Prononcé : \fbox{wó kěyǐ} (2+3) \quad \fbox{wǒ kěyǐ} (3+3)
    \item \zh{身体} (shēntǐ) $\rightarrow$ Prononcé : \fbox{shēntǐ} (1+3) \quad \fbox{shén tǐ} (2+3)
    \item \zh{锻炼} (duànliàn) $\rightarrow$ Prononcé : \fbox{duànliàn} (4+4) \quad \fbox{duánliàn} (2+4)
\end{enumerate}

\tcblower
\textbf{Solution détaillée :}
\begin{enumerate}
    \item \zh{很好} : \textbf{hén hǎo}. Règle : 3+3 $\rightarrow$ 2+3. Le premier \zh{很} devient montant.
    \item \zh{我可以} : \textbf{wó kěyǐ}. Règle : 3+3 $\rightarrow$ 2+3. \zh{我} (wǒ) devient montant \zh{wó}.
    \item \zh{身体} : \textbf{shēntǐ}. Pas de sandhi : 1ʳᵉ ton + 3ᵉ ton. Le 3ᵉ ton reste 3ᵉ (souvent réalisé en demi-ton bas 21).
    \item \zh{锻炼} : \textbf{duànliàn}. Pas de sandhi : 4ᵉ ton + 4ᵉ ton. Les deux restent 4ᵉ ton (voir section suivante).
\end{enumerate}
\end{exoresolu}

\courssub{Deux quatrièmes tons consécutifs : articulation et détachement}

Contrairement aux 3ᵉ tons, deux 4ᵉ tons ne subissent \textbf{pas de modification tonale obligatoire} (pas de sandhi standard). Cependant, ils posent un problème de \cle{prosodie} : la chute brutale de la première syllabe suivie immédiatement d'une seconde chute crée un effet « saccadé » ou « avalé » si l'on ne marque pas la frontière syllabique. Le vocabulaire du sport en est truffé.

\begin{definition}[Mots à double 4ᵉ ton du thème (HSK 2-3)]
\begin{center}
\begin{tabularx}{\linewidth}{|X|X|X|X|}
\hline
\rowcolor{poppinkL} \textbf{Caractères} & \textbf{Pinyin} & \textbf{Nature} & \textbf{Traduction} \\
\hline
运动 & yùndòng & nom / verbe & sport, exercice / bouger \\
\hline
锻炼 & duànliàn & verbe / nom & exercice physique, s'entraîner \\
\hline
重要 & zhòngyào & adjectif & important \\
\hline
决定 & juédìng & verbe / nom & décider, décision \\
\hline
\end{tabularx}
\end{center}
\end{definition}

\begin{propriete}[Articulation des double 4ᵉ tons]
Pour ne pas « avaler » la deuxième syllabe :
\begin{enumerate}
    \item \textbf{Marquer le pic de hauteur :} La première syllabe (ex: \zh{yùn}) doit partir haut (niveau 5) et chuter net (niveau 1).
    \item \textbf{Remonter (reset) :} Avant la deuxième syllabe (\zh{dòng}), la voix doit \textbf{remonter invisiblement} au niveau haut (5) pour pouvoir rechuter. C'est un micro-souffle, une micro-pause laryngée.
    \item \textbf{Éviter le « glissement » :} Ne pas faire \zh{yùndooòng} (une seule chute longue). Faire \zh{yùn. Dòng.} (deux chutes distinctes).
\end{enumerate}
\end{propriete}

\begin{exemplebox}[Entraînement ciblé : \zh{运动}、\zh{锻炼}、\zh{重要}、\zh{决定}]
\begin{itemize}
    \item \zh{我喜欢运动。 Wǒ xǐhuan yùndòng.} $\rightarrow$ \zh{yùn} (chute nette) \textbf{[reset]} \zh{dòng} (chute nette).
    \item \zh{锻炼身体。 Duànliàn shēntǐ.} $\rightarrow$ \zh{duàn} (sec) \textbf{[reset]} \zh{liàn} (sec) \zh{shēn} (plat) \zh{tǐ} (bas).
    \item \zh{这很重要。 Zhè hěn zhòngyào.} $\rightarrow$ \zh{zhòng} (chute) \textbf{[reset]} \zh{yào} (chute).
    \item \zh{我决定跑步。 Wǒ juédìng pǎobù.} $\rightarrow$ \zh{jué} (chute) \textbf{[reset]} \zh{dìng} (chute).
\end{itemize}
\end{exemplebox}

\begin{attention}[Erreur fréquente : « avaler » le 2ᵉ 4ᵉ ton]
Dire \zh{*yùndooòng} (une seule syllabe longue descendante) fait perdre le sens de \zh{动} (action). L'auditeur entendra peut-être \zh{云} (nuage) + son indéfini. \textbf{Exercice :} Tapez dans les mains sur chaque 4ᵉ ton : \textbf{CLAP} \zh{yùn} \textbf{CLAP} \zh{dòng}.
\end{attention}


\courssub{Exercices de discrimination et de production orale (Shadowing)}

La maîtrise des tons ne s'acquiert pas seulement par l'analyse, mais par l'imitation massive et la boucle audio-motrice.

\begin{methode}[Technique du Shadowing (répétition en écho) pour le dialogue modèle]
Le \emph{shadowing} consiste à répéter ce que l'on entend \textbf{quasi simultanément} (décalage de 0,5 à 1 seconde), en imitant non seulement les tons, mais l'intonation globale, le rythme, les pauses et l'intensité.
\begin{enumerate}
    \item \textbf{Écoute globale :} Écoutez le dialogue modèle (Section 3 ou 5) une fois sans parler. Visualisez la « mélodie » de la phrase.
    \item \textbf{Shadowing segmenté :} L'enseignant lit une phrase (ex: \zh{我认为运动很重要。 Wǒ rènwéi yùndòng hěn zhòngyào.}). Les élèves répètent en écho immédiat, calqués sur le modèle.
    \item \textbf{Shadowing continu :} L'enseignant (ou l'audio) lit le dialogue complet à vitesse normale. Les élèves parlent \textbf{par-dessus} (voix basse), essayant de coller parfaitement au flux.
    \item \textbf{Autocorrection :} Enregistrez-vous (téléphone) lors de l'étape 3. Réécoutez et comparez au modèle original. Notez les tons défaillants (souvent les 3ᵉ tons non « sablés » ou les 4ᵉ tons collés).
\end{enumerate}
\end{methode}

\begin{experience}[Jeu de discrimination : « Levez la carte du ton »]
\textbf{Matériel :} 4 cartes par élève (ou 4 doigts levés : 1=index, 2=index+majeur, 3=3 doigts, 4=4 doigts) représentant les 4 tons.
\textbf{Déroulement :}
\begin{enumerate}
    \item L'énonciateur (prof ou élève désigné) lit un \textbf{mot isolé} du vocabulaire de la leçon (liste ci-dessous).
    \item Les élèves lèvent \textbf{immédiatement} la carte/doigts correspondant au ton de la \textbf{dernière syllabe} (ou de l'unique syllabe).
    \item Vérification collective : l'enseignant montre la bonne carte. Discussion si divergence.
\end{enumerate}
\textbf{Liste de mots pour l'exercice (mélangés) :}
\begin{multicols}{2}
\begin{itemize}
    \item \zh{身 shēn} (1)
    \item \zh{强 qiáng} (2)
    \item \zh{体 tǐ} (3)
    \item \zh{健 jiàn} (4)
    \item \zh{跑 pǎo} (3)
    \item \zh{泡 pào} (4)
    \item \zh{练 liàn} (4)
    \item \zh{连 lián} (2)
    \item \zh{动 dòng} (4)
    \item \zh{运 yùn} (4)
    \item \zh{重 zhòng} (4)
    \item \zh{要 yào} (4)
    \item \zh{好 hǎo} (3)
    \item \zh{很 hěn} (3 $\rightarrow$ contexte sandhi)
    \item \zh{我 wǒ} (3 $\rightarrow$ contexte sandhi)
    \item \zh{可 kě} (3)
\end{itemize}
\end{multicols}
\textbf{Variante difficile :} L'enseignant lit un \textbf{disyllabe} (ex: \zh{运动 yùndòng}). Les élèves lèvent \textbf{deux cartes successives} (ou montrent deux doigts puis quatre doigts) pour \zh{yùn} (4) puis \zh{dòng} (4).
\end{experience}

\begin{experience}[Atelier Shadowing guidé : Mini-dialogue « Sport et Santé »]
\textbf{Support :} Dialogue court (à projeter ou distribuer) intégrant les points travaillés (sandhi 3+3, double 4, paires minimales).
\begin{textebox}[Dialogue d'entraînement phonétique]
A: \zh{你好！你身体好吗？} (Nǐ hǎo! Nǐ shēntǐ hǎo ma?) \\
B: \zh{我很好。我每天运动。} (Wǒ hěn hǎo. Wǒ měitiān yùndòng.) \\
A: \zh{你跑步吗？} (Nǐ pǎobù ma?) \\
B: \zh{不，我练习太极拳。} (Bù, wǒ liànxí tàijíquán.) \\
A: \zh{锻炼很重要。} (Duànliàn hěn zhòngyào.) \\
B: \zh{对，坚持最重要。} (Duì, jiānchí zuì zhòngyào.)
\end{textebox}
\textbf{Consignes :}
\begin{enumerate}
    \item Repérez les \textbf{sandhi 3+3} : \zh{很 好} (hén hǎo), \zh{你 好} (ní hǎo).
    \item Repérez les \textbf{double 4} : \zh{运动 yùndòng}, \zh{锻炼 duànliàn}, \zh{重要 zhòngyào}. Préparez le « reset » de hauteur.
    \item Repérez la \textbf{paire minimale} : \zh{跑 pǎo} (3) vs \zh{泡 pào} (4) — ici c'est \zh{跑}.
    \item \textbf{Phase 1 :} Lecture chorale lente, main traçant les courbes dans l'air.
    \item \textbf{Phase 2 :} Shadowing individuel (casque audio ou prof modèle) $\times$ 3.
    \item \textbf{Phase 3 :} Jeu de rôle en binômes sans texte, en se forçant à utiliser \zh{运动}、\zh{锻炼}、\zh{重要} avec la bonne prosodie.
\end{enumerate}
\end{experience}

\begin{aretenir}[Check-list prononciation pour l'exposé oral (Section 6)]
Avant de passer à l'exposé et au débat, vérifiez que vous maîtrisez :
\begin{itemize}
    \item[\checkmark] Les 4 tons sur les mots-clés : \zh{身 shēn}, \zh{体 tǐ}, \zh{健 jiàn}, \zh{康 kāng}, \zh{运 yùn}, \zh{动 dòng}, \zh{练 liàn}, \zh{重 zhòng}, \zh{要 yào}.
    \item[\checkmark] La distinction \zh{跑 pǎo} (3) / \zh{泡 pào} (4) et \zh{练 liàn} (4) / \zh{连 lián} (2).
    \item[\checkmark] L'application automatique du \textbf{sandhi 3+3} : \zh{很好 hén hǎo}, \zh{我可以 wó kěyǐ}, \zh{你好 ní hǎo}.
    \item[\checkmark] L'articulation nette des \textbf{double 4ᵉ tons} : \zh{yùndòng}, \zh{duànliàn}, \zh{zhòngyào}, \zh{juédìng} (reset de hauteur entre les deux syllabes).
    \item[\checkmark] L'intonation de phrase : déclaration (descendante finale), question \zh{吗 ma} (montante légère sur la particule neutre), insistance (ton haut sur l'adverbe \zh{很 hěn}, \zh{最 zuì}).
\end{itemize}
\end{aretenir}

\coursec{Leçon 11 — Expression orale : importance de la pratique du sport}

\courssub{Vocabulaire du sport et de la santé}

\begin{definition}[Lexique de base (HSK 2-3) — Sports, santé, fréquence et opinion]
Le tableau ci-dessous regroupe le vocabulaire essentiel pour s'exprimer sur la pratique sportive, ses bienfaits et ses contraintes. Maîtrisez les classificateurs (\zh{个} pour \zh{小时}, \zh{场} pour \zh{比赛}, \zh{次} pour la fréquence) et la place des adverbes de fréquence (\zh{每天}, \zh{经常}) avant le verbe.
\end{definition}

\begin{center}
\begin{tabularx}{\linewidth}{|X|X|X|X|}
\hline
\rowcolor{popblueL} \textbf{Caractères} & \textbf{Pinyin} & \textbf{Nature} & \textbf{Traduction} \\ \hline
\zh{运动} & yùndòng & VO / N & Faire du sport / Le sport, l'exercice \\ \hline
\zh{锻炼} & duànliàn & V & S'entraîner, faire de l'exercice physique \\ \hline
\zh{跑步} & pǎobù & VO & Courir (pied + étape) \\ \hline
\zh{游泳} & yóuyǒng & V & Nager \\ \hline
\zh{打篮球} & dǎ lánqiú & VO & Jouer au basket (frapper + panier-ballon) \\ \hline
\zh{踢足球} & tī zúqiú & VO & Jouer au foot (donner un coup de pied + foot-ballon) \\ \hline
\zh{打乒乓球} & dǎ pīngpāngqiú & VO & Jouer au ping-pong \\ \hline
\zh{练瑜伽} & liàn yújiā & VO & Faire du yoga \\ \hline
\zh{健康} & jiànkāng & Adj / N & En bonne santé / La santé \\ \hline
\zh{强身健体} & qiángshēn jiàntǐ & Id. & Renforcer le corps (objectif du sport scolaire) \\ \hline
\zh{放松} & fàngsōng & V / Adj & Se détendre / Détendu \\ \hline
\zh{累} & lèi & Adj & Fatigué \\ \hline
\zh{受伤} & shòushāng & V & Se blesser \\ \hline
\zh{膝盖} & xīgài & N & Genou \\ \hline
\zh{肺} & fèi & N & Poumon \\ \hline
\zh{每天} & měitiān & Adv & Tous les jours \\ \hline
\zh{每周} & měizhōu & Adv & Chaque semaine \\ \hline
\zh{经常} & jīngcháng & Adv & Souvent \\ \hline
\zh{有时候} & yǒushíhou & Adv & Parfois \\ \hline
\zh{小时} & xiǎoshí & N / M & Heure (durée) \\ \hline
\zh{分钟} & fēnzhōng & N / M & Minute \\ \hline
\zh{半个小时} & bàn ge xiǎoshí & Q & Une demi-heure \\ \hline
\zh{应该} & yīnggāi & Modal & Devoir (obligation morale, conseil) \\ \hline
\zh{同意} & tóngyì & V & Être d'accord \\ \hline
\zh{不同意} & bù tóngyì & V & Ne pas être d'accord \\ \hline
\zh{觉得} & juéde & V & Penser, trouver (avis subjectif) \\ \hline
\zh{认为} & rènwéi & V & Considérer, estimer (plus formel) \\ \hline
\end{tabularx}
\end{center}

\courssub{Dialogue modèle 1 : parler de ses sports}

\begin{textebox}[Dialogue : Quels sports aimes-tu ?]
A: 你喜欢什么运动？\\
Nǐ xǐhuan shénme yùndòng?\\
Quel sport aimes-tu ?

B: 我喜欢跑步和游泳。你呢？\\
Wǒ xǐhuan pǎobù hé yóuyǒng. Nǐ ne?\\
J'aime courir et nager. Et toi ?

A: 我经常打篮球。有时候踢足球。\\
Wǒ jīngcháng dǎ lánqiú. Yǒushíhou tī zúqiú.\\
Je joue souvent au basket. Parfois au foot.

B: 运动对身体好吗？\\
Yùndòng duì shēntǐ hǎo ma?\\
Le sport est-il bon pour le corps ?

A: 当然，运动能强身健体。\\
Dāngrán, yùndòng néng qiángshēn jiàntǐ.\\
Bien sûr, le sport renforce le corps.

B: 那我们周末一起去游泳吧。\\
Nà wǒmen zhōumò yīqǐ qù yóuyǒng ba.\\
Alors, on va nager ensemble ce week-end ?

A: 好主意！\\
Hǎo zhǔyi!\\
Bonne idée !
\end{textebox}

\begin{methode}[Exploitation du dialogue modèle 1]
\begin{enumerate}
    \item \textbf{Compréhension globale :} Identifier les interlocuteurs (deux lycéens), le lieu (cour de récréation / messagerie) et le sujet (préférences sportives + invitation).
    \item \textbf{Repérage phonétique :} Surlignez les tons 3 (\zh{运} yùn, \zh{动} dòng, \zh{游} yóu, \zh{泳} yǒng, \zh{打} dǎ, \zh{踢} tī, \zh{身} shēn, \zh{体} tǐ, \zh{好} hǎo, \zh{周} zhōu, \zh{末} mò, \zh{一} yī, \zh{起} qǐ, \zh{主} zhǔ, \zh{意} yì). Appliquez la règle du \textbf{ton 3 + ton 3} : \zh{运动} $\rightarrow$ \zh{yóu} (ton 2) + \zh{dòng} ; \zh{游泳} $\rightarrow$ \zh{yóu} (ton 2) + \zh{yǒng} ; \zh{强身健体} $\rightarrow$ \zh{qiáng} (2) + \zh{shén} (2) + \zh{jiàn} (4) + \zh{tǐ} (3).
    \item \textbf{Structures cibles :} \zh{喜欢 + Nom/Verbe} ; \zh{经常 / 有时候 + Verbe} (adverbes de fréquence AVANT le verbe) ; \zh{对 ... 好} (bon pour...) ; \zh{能 + Verbe} (capacité/possibilité) ; \zh{那 ... 吧} (suggestion).
    \item \textbf{Entraînement guidé :} En binôme, remplacez les sports (\zh{打乒乓球}, \zh{练瑜伽}) et la fréquence (\zh{每天}, \zh{每周两次}). Inversez les rôles.
\end{enumerate}
\end{methode}

\courssub{Formules fonctionnelles : Présenter, opiner, relancer}

\begin{propriete}[Boîte à outils communicatifs pour l'oral (Niveau A2/B1)]
Ces formules sont les « briques » de toute interaction orale structurée. Classez-les par fonction dans votre cahier.
\end{propriete}

\begin{center}
\begin{tabularx}{\linewidth}{|X|X|X|X|}
\hline
\rowcolor{popblueL} \textbf{Fonction} & \textbf{Chinois (Caractères)} & \textbf{Pinyin} & \textbf{Français} \\ \hline
\multicolumn{4}{|c|}{\cellcolor{poporangeL}\textbf{1. S'informer / Présenter un sujet}} \\ \hline
Demander le sport préféré & 你喜欢什么运动？ & Nǐ xǐhuan shénme yùndòng? & Quel sport aimes-tu ? \\ \hline
Demander la fréquence & 你经常锻炼吗？ & Nǐ jīngcháng duànliàn ma? & T'entraînes-tu souvent ? \\ \hline
Présenter son avis & 我觉得 / 我认为 & Wǒ juéde / Wǒ rènwéi & Je pense que / À mon avis \\ \hline
\multicolumn{4}{|c|}{\cellcolor{popgreenL}\textbf{2. Exprimer l'accord / le désaccord}} \\ \hline
Être d'accord & 我同意 / 对 / 好的 / 行 & Wǒ tóngyì / Duì / Hǎo de / Xíng & Je suis d'accord / Oui / D'accord / Ça marche \\ \hline
Ne pas être d'accord & 我不同意 / 我不觉得... / 不一定 & Wǒ bù tóngyì / Wǒ bù juéde... / Bù yídìng & Pas d'accord / Je ne pense pas... / Pas nécessairement \\ \hline
\multicolumn{4}{|c|}{\cellcolor{poppurpleL}\textbf{3. Argumenter / Nuancer}} \\ \hline
Donner une raison & 因为... 所以... & Yīnwèi... Suǒyǐ... & Parce que... donc... \\ \hline
Concéder & 你说得对 / 有道理 / 是真的 & Nǐ shuō de duì / Yǒu dàoli / Shì zhēn de & Tu as raison / C'est sensé / C'est vrai \\ \hline
S'opposer (doux) & 可是 / 不过 & Kěshì / Bùguò & Mais / Toutefois \\ \hline
S'opposer (fort) & 但是 & Dànshì & Mais / Cependant \\ \hline
\multicolumn{4}{|c|}{\cellcolor{poppinkL}\textbf{4. Proposer / Conclure / Relancer}} \\ \hline
Proposer un compromis & 那 / 那么 / 要不 & Nà / Nàme / Yàobù & Alors / Dans ce cas / Sinon \\ \hline
Faire une suggestion & 我们可以... 吧 / 要不... & Wǒmen kěyǐ... ba / Yàobù... & On peut... / Et si... \\ \hline
Demander confirmation & 好吗？ / 怎么样？ / 行吗？ & Hǎo ma? / Zěnme yàng? / Xíng ma? & D'accord ? / Qu'en penses-tu ? / Ça va ? \\ \hline
Clore l'échange & 那就这么说定了 / 好的，再见 & Jiù nàme shuō dìng le / Hǎo, zàijiàn & C'est entendu / Au revoir \\ \hline
\end{tabularx}
\end{center}

\courssub{Travail sur les tons et la prononciation}

\begin{propriete}[Rappel des règles de changement de tons (Sandhi) — Indispensable pour l'oral]
\begin{enumerate}
    \item \textbf{Deux tons 3 consécutifs :} Le premier devient ton 2 (montant). Ex: \zh{运动} (yùn+dòng $\rightarrow$ \textbf{yún} dòng), \zh{游泳} (yǒu+yǒng $\rightarrow$ \textbf{yóu} yǒng), \zh{强身健体} (qiǎng+shēn+jiàn+tǐ $\rightarrow$ \textbf{qiáng} shén jiàn tǐ).
    \item \textbf{\zh{一} (yī) :} Ton 1 seul / à la fin / avant chiffre. \textbf{Ton 2 (yí)} avant ton 4 (\zh{一个} yí ge). \textbf{Ton 4 (yì)} avant tons 1, 2, 3 (\zh{一天} yì tiān, \zh{一年} yì nián, \zh{一起} yì qǐ).
    \item \textbf{\zh{不} (bù) :} Ton 4 normal. \textbf{Ton 2 (bú)} avant ton 4 (\zh{不是} bú shì, \zh{不要} bú yào).
    \item \textbf{Ton neutre (轻声) :} Particules (\zh{了, 吧, 呢, 吗, 的, 地, 得}), classificateurs (\zh{个, 场, 次} souvent), 2e syllabe de certains mots (\zh{东西} dōngxi, \zh{时候} shíhou, \zh{分钟} fēnzhōng $\rightarrow$ fēnzhong).
\end{enumerate}
\end{propriete}

\begin{experience}[Atelier prononciation : « Le tribunal des tons »]
\begin{enumerate}
    \item L'enseignant dicte 10 mots/phrases du vocabulaire (ex: \zh{每天锻炼}, \zh{游泳池}, \zh{不累}, \zh{一小时}, \zh{强身健体}).
    \item Les élèves écrivent le pinyin \textbf{avec les tons réels prononcés} (après sandhi), pas les tons dictionnaire.
    \item Correction collective au tableau : un élève au tableau, la classe valide ou corrige.
    \item \textbf{Focus :} \zh{半个小时} (bàn ge xiǎoshí $\rightarrow$ 4 - neutre - 3-2) ; \zh{一起去} (yì qǐ qù $\rightarrow$ 4 - 3 - 4).
\end{enumerate}
\end{experience}

% --- DEBUT DU BLOC COLLEGUE CORRIGÉ ET AUGMENTÉ ---

\coursec{Dialogue modèle 2 : mini-débat « faut-il faire du sport tous les jours ? »}

Après avoir travaillé la prononciation des tons et le vocabulaire de base dans les sections précédentes, nous passons maintenant à la mise en situation communicative complexe. Ce dialogue modèle illustre un échange d'opinions structuré : exprimer un point de vue, manifester un désaccord, argumenter, concéder un point à l'adversaire, proposer un compromis et aboutir à un accord. C'est la structure même de l'argumentation orale attendue au baccalauréat et dans la vie courante.

\courssub{Présentation et écoute du dialogue}

\begin{textebox}[Dialogue : Faut-il faire du sport tous les jours ?]
A: 我觉得我们应该每天锻炼。\\
Wǒ juéde wǒmen yīnggāi měitiān duànliàn.\\
Je pense que nous devrions nous entraîner tous les jours.

B: 我不同意。每天锻炼太累了。\\
Wǒ bù tóngyì. Měitiān duànliàn tài lèi le.\\
Je ne suis pas d'accord. S'entraîner tous les jours est trop fatigant.

A: 可是运动对身体很好，可以让我们更健康。\\
Kěshì yùndòng duì shēntǐ hěn hǎo, kěyǐ ràng wǒmen gèng jiànkāng.\\
Mais le sport est très bon pour le corps, il peut nous rendre plus en santé.

B: 你说得对，可是我们也要学习。\\
Nǐ shuō de duì, kěshì wǒmen yě yào xuéxí.\\
Tu as raison, mais nous devons aussi étudier.

A: 那我们可以每天运动半个小时。\\
Nà wǒmen kěyǐ měitiān yùndòng bàn ge xiǎoshí.\\
Alors nous pouvons faire du sport une demi-heure par jour.

B: 好，我同意！\\
Hǎo, wǒ tóngyì!\\
D'accord, je suis d'accord !
\end{textebox}

\begin{methode}[Comment travailler ce dialogue en binôme]
\begin{enumerate}
    \item \textbf{Écoute globale :} L'enseignant lit le dialogue à vitesse normale. Les élèves identifient le sujet (le sport quotidien) et l'issue (accord sur 30 min).
    \item \textbf{Repérage des tons et sandhi :} Sur le texte, surlignez les tons 3 (basculants) et les tons 4 (chutants) dans les mots-clés : \zh{锻炼} (duànliàn, 4-4), \zh{累} (lèi, 4), \zh{健康} (jiànkāng, 4-1), \zh{半个小时} (bàn ge xiǎoshí, 4 - neutre - 3-2). Attention au changement de ton de \zh{一} (yī $\rightarrow$ yí avant ton 4, yì avant tons 1/2/3) dans \zh{一个} (yí ge) mais \zh{半个} (bàn ge) ne change pas. \zh{更} (gèng, 4) dans \zh{更健康}.
    \item \textbf{Lecture chorale puis individuelle :} Insistez sur l'intonation montante de \zh{我觉得} (opinion) et l'intonation ferme de \zh{我不同意}. Marquez la pause après \zh{那} (pivot).
    \item \textbf{Jeu de rôle :} Inversez les rôles. Puis changez le sport : \zh{游泳} (yóuyǒng, nager), \zh{跑步} (pǎobù, courir), \zh{打篮球} (dǎ lánqiú, basket).
\end{enumerate}
\end{methode}

\courssub{Analyse structurelle : les étapes du débat}

Le dialogue suit une progression logique que tout orateur doit maîtriser. Observez le schéma ci-dessous qui modélise l'enchaînement des actes de parole.

\begin{popfigure}
\begin{tikzpicture}[
    node distance=1.1cm and 1.5cm,
    box/.style={rectangle, rounded corners, draw=popblue, thick, fill=popblueL, text width=2.8cm, align=center, minimum height=1.1cm, font=\small},
    arrow/.style={-Stealth, thick, popdark},
    labelarrow/.style={midway, fill=white, inner sep=2pt, font=\tiny, text=popdark, align=center}
]
% Nœuds
\node[box, align=center] (opinion){Opinion\\ \zh{我觉得我们应该...}\\ (Je pense que nous devrions...)};
\node[box, below=of opinion, align=center] (desaccord){Désaccord\\ \zh{我不同意。太...了}\\ (Pas d'accord. Trop...)};
\node[box, below=of desaccord, align=center] (argument){Argument\\ \zh{可是...对...很好}\\ (Mais... est bon pour...)};
\node[box, below=of argument, align=center] (concession){Concession\\ \zh{你说得对，可是...}\\ (Tu as raison, mais...)};
\node[box, below=of concession, align=center] (compromis){Compromis\\ \zh{那我们可以...}\\ (Alors nous pouvons...)};
\node[box, below=of compromis, align=center] (accord){Accord\\ \zh{好，我同意！}\\ (D'accord, j'accepte !)};

% Flèches
\draw[arrow] (opinion) -- (desaccord) node[left] {réplique};
\draw[arrow] (desaccord) -- (argument) node[left] {argumentation};
\draw[arrow] (argument) -- (concession) node[left] {contre-argument};
\draw[arrow] (concession) -- (compromis) node[left] {pivot (\zh{那})};
\draw[arrow] (compromis) -- (accord) node[left] {résolution};
\end{tikzpicture}
\legende{Organigramme de la structure argumentative du mini-débat}
\end{popfigure}

\courssub{Formules fonctionnelles : Exprimer l'opinion, l'accord et le désaccord}

\begin{propriete}[Vocabulaire fonctionnel du débat oral — Niveau approfondi]
\begin{center}
\begin{tabularx}{\linewidth}{|X|X|X|X|}
\hline
\rowcolor{popblueL} \textbf{Fonction} & \textbf{Chinois (Caractères)} & \textbf{Pinyin} & \textbf{Français} \\ \hline
Exprimer son avis & 我觉得 / 我认为 & wǒ juéde / wǒ rènwéi & Je pense que / À mon avis \\ \hline
Exprimer une obligation morale & 应该 / 应当 & yīnggāi / yīngdāng & Devoir / Il faudrait \\ \hline
Être d'accord & 我同意 / 对 / 好的 / 行 & wǒ tóngyì / duì / hǎo de / xíng & Je suis d'accord / Oui / D'accord / Ça marche \\ \hline
Ne pas être d'accord & 我不同意 / 我不觉得... / 不一定 & wǒ bù tóngyì / wǒ bù juéde... / bù yídìng & Je ne suis pas d'accord / Je ne pense pas... / Pas nécessairement \\ \hline
Concéder un point & 你说得对 / 有道理 / 确实 & nǐ shuō de duì / yǒu dàoli / quèshí & Tu as raison / C'est sensé / Effectivement \\ \hline
Introduire une objection & 可是 / 但是 / 不过 & kěshì / dànshì / bùguò & Mais / Cependant / Toutefois \\ \hline
Proposer un compromis & 那 / 那么 / 要不 & nà / nàme / yàobù & Alors / Dans ce cas / Sinon \\ \hline
Accepter le compromis & 好 / 行 / 可以 / 成 & hǎo / xíng / kěyǐ / chéng & D'accord / Ça marche / Peut se faire / C'est bon (oral) \\ \hline
\end{tabularx}
\end{center}
\end{propriete}

\courssub{Points de grammaire clés du dialogue}

\begin{propriete}[Structure : \textbf{应该} (yīnggāi) + Verbe — Exprimer le devoir, la recommandation]
L'auxiliaire modal \cle{应该} (ou \cle{应当} yīngdāng, plus formel, écrit) signifie « devoir », « il faut », « il serait bon de ». Il exprime une obligation morale, un conseil ou une forte probabilité. Il ne s'utilise pas pour une contrainte physique (ex: « je dois partir » contrainte horaire $\rightarrow$ \zh{得} děi / \zh{必须} bìxū).
\medskip
\textbf{Structure :} Sujet + 应该 / 应当 + Verbe (+ Complément)
\medskip
\begin{center}
\begin{tabularx}{\linewidth}{|X|X|X|}
\hline
\rowcolor{poporangeL} \textbf{Chinois} & \textbf{Pinyin} & \textbf{Français} \\ \hline
我们应该每天锻炼。 & Wǒmen yīnggāi měitiān duànliàn. & Nous devrions faire de l'exercice tous les jours. \\ \hline
学生应当努力学习。 & Xuéshēng yīngdāng nǔlì xuéxí. & Les étudiants doivent étudier dur. \\ \hline
你应该早点睡觉。 & Nǐ yīnggāi zǎo diǎn shuìjiào. & Tu devrais te coucher plus tôt. \\ \hline
\end{tabularx}
\end{center}
\textbf{Négation :} Sujet + \textbf{不应该} + Verbe. (Ex: \zh{你不应该吸烟。} Nǐ bù yīnggāi xīyān. — Tu ne devrais pas fumer.)
\textbf{Interrogation :} Sujet + 应该 + Verbe + \zh{吗} ? Ou \zh{应该不应该}... (A-not-A formel).
\end{propriete}

\begin{attention}[Piège majeur : Position de \textbf{应该} et négation]
En français, « nous devons nous entraîner » place le verbe modal « devons » avant l'infinitif. En chinois, \zh{应该} se place \textbf{impérativement AVANT le verbe principal}, jamais après.
\begin{itemize}
    \item \textbf{Correct :} \zh{我们应该锻炼。} (Sujet + 应该 + Verbe)
    \item \textbf{Incorrect :} \zh{*我们锻炼应该。} (Influence de la structure française ou placement final erroné).
\end{itemize}
De même, la négation \zh{不} précède \zh{应该} : \zh{不应该去} (ne pas devoir aller), jamais \zh{应该不去} (qui signifierait « devoir ne pas aller » — nuance forte d'interdiction).
\end{attention}

\begin{propriete}[Conjonctions adversatives : \textbf{可是} vs \textbf{但是} vs \textbf{不过}]
Ces trois mots traduisent « mais », mais leur registre et leur force diffèrent. Le choix influence le ton du débat.
\begin{center}
\begin{tabularx}{\linewidth}{|X|X|X|X|}
\hline
\rowcolor{poppurpleL} \textbf{Mot} & \textbf{Pinyin} & \textbf{Registre / Nuance} & \textbf{Exemple} \\ \hline
\zh{可是} & kěshì & Courant, oral, transition douce, lien énonciatif & \zh{我想去，可是没时间。} (Je veux y aller, mais pas le temps.) \\ \hline
\zh{但是} & dànshì & Plus formel, écrit, opposition forte, marque une rupture & \zh{他很忙，但是他来了。} (Il est occupé, mais il est venu.) \\ \hline
\zh{不过} & bùguò & Oral, familier, « simplement / toutefois », souvent après concession & \zh{贵是贵，不过很好。} (C'est cher, ceci dit c'est bon.) \\ \hline
\end{tabularx}
\end{center}
Dans notre dialogue, \zh{可是} est utilisé deux fois par A et B pour maintenir un ton dialogué (conversationnel) et peu conflictuel, typique d'une discussion entre camarades.
\end{propriete}

\begin{propriete}[La particule \textbf{那} (nà) / \textbf{那么} (nàme) en début de phrase : « Alors / Dans ce cas »]
C'est un marqueur discursif essentiel pour \textbf{proposer un compromis} ou tirer une conséquence logique de ce qui vient d'être dit. Il marque le pivot du débat (transition du conflit vers la solution).
\medskip
\textbf{Structure :} \zh{那 / 那么} + Sujet + Verbe (+ Proposition de solution)
\medskip
\begin{center}
\begin{tabularx}{\linewidth}{|X|X|X|}
\hline
\rowcolor{popgreenL} \textbf{Chinois} & \textbf{Pinyin} & \textbf{Français} \\ \hline
\zh{那我们可以每天运动半个小时。} & Nà wǒmen kěyǐ měitiān yùndòng bàn ge xiǎoshí. & Alors, nous pouvons faire du sport 30 min par jour. \\ \hline
\zh{你不喜欢跑步？那我们去游泳吧。} & Nǐ bù xǐhuan pǎobù? Nà wǒmen qù yóuyǒng ba. & Tu n'aimes pas courir ? Alors allons nager. \\ \hline
\zh{那么，我们下周再见。} & Nàme, wǒmen xià zhōu zàijiàn. & Dans ce cas, à la semaine prochaine. \\ \hline
\end{tabularx}
\end{center}
\cle{那} (nà) est plus bref, très oral, familier. \cle{那么} (nàme) est légèrement plus formel, plus insistant sur la conséquence logique.
\end{propriete}

\begin{propriete}[Structure \textbf{「Adj + 了」} : Changement d'état / Excès — \zh{太累了}]
L'adverbe de degré \cle{太} (tài, « trop ») s'associe obligatoirement à la particule \cle{了} (le) en fin de phrase (ou de proposition) pour marquer l'excès ou un nouvel état ressenti par le locuteur.
\medskip
\textbf{Structure :} Sujet + 太 + Adjectif + 了
\medskip
\begin{center}
\begin{tabularx}{\linewidth}{|X|X|X|}
\hline
\rowcolor{poppinkL} \textbf{Chinois} & \textbf{Pinyin} & \textbf{Français} \\ \hline
\zh{每天锻炼太累了。} & Měitiān duànliàn tài lèi le. & S'entraîner tous les jours est trop fatigant. \\ \hline
\zh{这个作业太难了。} & Zhège zuòyè tài nán le. & Ce devoir est trop difficile. \\ \hline
\zh{今天太热了！} & Jīntiān tài rè le! & Aujourd'hui il fait trop chaud ! \\ \hline
\end{tabularx}
\end{center}
\end{propriete}

\begin{attention}[Oublier le \textbf{了} final — Erreur classique HSK 2/3]
\zh{*太累} (sans \zh{了}) est une phrase inachevée en chinois standard. Le \zh{了} marque la réalisation de l'état « trop fatigué » dans la réalité présente (particule modale de phrase). Ne le confondez pas avec le \zh{了} d'accompli (verbe + \zh{了} = action terminée). Ici, pas de verbe d'action, c'est une phrase adjectivale.
\end{attention}


\courssub{Exemples traduits et tableaux de substitution}

Pour varier le débat, utilisez le tableau de substitution ci-dessous. Remplacez les éléments soulignés dans le dialogue original.

\begin{center}
\begin{tabularx}{\linewidth}{|X|X|X|}
\hline
\rowcolor{popblueL} \textbf{Sujet du débat} & \textbf{Argument « Pour » (A)} & \textbf{Argument « Contre » (B)} \\ \hline
\zh{游泳} (yóuyǒng) — Nager & \zh{游泳对肺很好。} (Bon pour les poumons.) & \zh{游泳池太远了。} (Piscine trop loin.) \\ \hline
\zh{跑步} (pǎobù) — Courir & \zh{跑步不花钱。} (Courir ne coûte rien.) & \zh{跑步伤膝盖。} (Abîme les genoux.) \\ \hline
\zh{打篮球} (dǎ lánqiú) — Basket & \zh{打篮球能长高。} (Peut faire grandir.) & \zh{太危险，容易受伤。} (Trop dangereux, facile de se blesser.) \\ \hline
\zh{练瑜伽} (liàn yújiā) — Yoga & \zh{瑜伽让人放松。} (Le yoga détend.) & \zh{动作太难了。} (Mouvements trop durs.) \\ \hline
\end{tabularx}
\end{center}

\begin{exemplebox}[Application : Adapter le dialogue (Sujet : \zh{跑步} pǎobù)]
\textbf{A:} 我觉得我们应该每天\zh{跑步}。 (Je pense qu'on devrait courir tous les jours.)\\
\textbf{B:} 我不同意。每天\zh{跑步}\zh{伤膝盖}。 (Pas d'accord. Courir tous les jours abîme les genoux.)\\
\textbf{A:} 可是\zh{跑步}\zh{不花钱}，可以\zh{强身健体}。 (Mais courir ne coûte rien, ça renforce le corps.)\\
\textbf{B:} 你说得对，可是\zh{天气太热}。 (Tu as raison, mais il fait trop chaud.)\\
\textbf{A:} \zh{那}我们可以\zh{晚上}跑\zh{半个小时}。 (Alors on peut courir 30 min le soir.)\\
\textbf{B:} 好，\zh{我同意}！ (D'accord, j'accepte !)
\end{exemplebox}

\courssub{Exercice résolu : Réorganiser le débat}

\begin{exoresolu}[Remettre les répliques dans l'ordre logique]
\textbf{Énoncé :} Les phrases suivantes sont mélangées. Réorganisez-les pour reconstituer un mini-débat cohérent sur le thème « 学中文难不难 ? » (Le chinois est-il difficile ?). Identifiez l'étape (Opinion, Désaccord, Argument, Concession, Compromis, Accord).

\begin{enumerate}
    \item 那我们可以每天学十五分钟。
    \item 我觉得学中文很难。
    \item 好，我同意！
    \item 我不觉得难。
    \item 可是中文很有用，以后去中国工作方便。
    \item 你说得对，可是汉字太多了。
\end{enumerate}

\tcblower
\textbf{Solution détaillée :}

L'ordre logique suit l'organigramme vu plus haut (A = Locuteur 1, B = Locuteur 2) :

1. \textbf{Opinion (A) :} \zh{我觉得学中文很难。} (Wǒ juéde xué Zhōngwén hěn nán.) — « Je trouve le chinois difficile. »
2. \textbf{Désaccord (B) :} \zh{我不觉得难。} (Wǒ bù juéde nán.) — « Moi, je ne trouve pas que ce soit difficile. » (Désaccord franc).
3. \textbf{Argument (A) :} \zh{可是中文很有用，以后去中国工作方便。} (Kěshì Zhōngwén hěn yǒuyòng, yǐhòu qù Zhōngguó gōngzuò fāngbiàn.) — « Mais le chinois est utile, aller travailler en Chine plus tard sera pratique. »
4. \textbf{Concession (B) :} \zh{你说得对，可是汉字太多了。} (Nǐ shuō de duì, kěshì Hànzì tài duō le.) — « Tu as raison, mais les caractères sont trop nombreux. »
5. \textbf{Compromis (A) :} \zh{那我们可以每天学十五分钟。} (Nà wǒmen kěyǐ měitiān xué shíwǔ fēnzhōng.) — « Alors on peut étudier 15 minutes par jour. »
6. \textbf{Accord (B) :} \zh{好，我同意！} (Hǎo, wǒ tóngyì!) — « D'accord, je suis d'accord ! »
\end{exoresolu}

\courssub{Élément socioculturel : La culture du sport en Chine}

\begin{savaistu}[Les exercices du matin : \zh{早操} (zǎocāo) et la réforme du sport scolaire]
Dans la quasi-totalité des écoles chinoises (primaire, collège, lycée), la journée commence par \zh{早操} (zǎocāo), littéralement « exercices du matin ». Vers 7h30 ou 8h00, tous les élèves se rassemblent sur le terrain de sport ou dans la cour, par classes, au son de la musique diffusée par haut-parleurs. Ils effectuent en synchronisation une série de mouvements d'étirement, de gymnastique rythmique (\zh{广播体操} guǎngbō tǐcāo) ou de course légère, encadrés par le professeur principal (\zh{班主任} bānzhǔrèn).
\medskip
Cette tradition, héritée de l'éducation physique soviétique des années 50 et adaptée, vise à \zh{强身健体} (qiángshēn jiàntǐ, renforcer le corps), à réveiller la concentration avant les cours et à renforcer l'esprit de collectif (\zh{集体主义} jítǐ zhǔyì). Le lundi matin, place souvent à la \zh{升旗仪式} (shēngqí yíshì, cérémonie de lever du drapeau) suivie d'un discours du directeur sous le drapeau national.
\medskip
\textbf{Réforme récente (2021+) :} Le ministère de l'Éducation chinois impose désormais \zh{每天一小时校园体育活动} (une heure d'activité sportive quotidienne sur le campus) et interdit d'occuper les cours d'EPS par d'autres matières. L'EPS entre au \zh{中考} (zhōngkǎo, brevet) et \zh{高考} (gāokǎo, bac) avec une note comptant pour le total.
\medskip
\textbf{Comparaison Cameroun / Chine :} Au Cameroun, l'EPS (Éducation Physique et Sportive) est une matière hebdomadaire (souvent 2h/semaine) avec un professeur spécialisé. En Chine, l'activité physique quotidienne (souvent 30-45 min le matin + cours EPS) est une routine institutionnelle obligatoire pour tous, intégrée à l'emploi du temps matinal, bien distincte du cours d'EPS hebdomadaire. Au Cameroun, le sport scolaire est souvent lié aux compétitions inter-écoles (FENASCO) ; en Chine, c'est une hygiène de vie quotidienne évaluée.
\end{savaistu}

\courssub{Exercices de consolidation}

\begin{exercice}[Production orale guidée : Mini-débat en binôme — Évaluation formative]
\textbf{Consigne :} En binôme, choisissez l'un des trois sujets ci-dessous. Suivez \textbf{impérativement} la structure du dialogue modèle (Opinion $\rightarrow$ Désaccord $\rightarrow$ Argument $\rightarrow$ Concession $\rightarrow$ Compromis $\rightarrow$ Accord). Utilisez \zh{应该}, \zh{可是/但是}, \zh{那}, \zh{太...了}, \zh{因为...所以...}. Préparez 3 minutes (notes en pinyin autorisées), puis présentez devant la classe (1 min 30 max par binôme). Évaluation : prononciation (tons), fluidité, respect de la structure, richesse lexicale.

\begin{enumerate}
    \item \textbf{Sujet 1 :} \zh{周末要不要补课？} (Faut-il faire des cours de rattrapage le week-end ?)
    \item \textbf{Sujet 2 :} \zh{中学生能不能带手机上学？} (Les collégiens/lycéens peuvent-ils apporter leur téléphone à l'école ?)
    \item \textbf{Sujet 3 :} \zh{吃零食好不好？} (Manger des snacks est-ce bien ?)
\end{enumerate}

\textbf{Aide lexicale :}
\zh{补课} (bǔkè, cours de rattrapage), \zh{带手机} (dài shǒujī, apporter portable), \zh{分心} (fēnxīn, se distraire), \zh{联系方便} (liánxì fāngbiàn, contact facile), \zh{零食} (língshí, snacks), \zh{长胖} (zhǎngpàng, grossir), \zh{解馋} (jiěchán, calmer une envie), \zh{视力下降} (shìlì xiàjiàng, baisse de vue), \zh{近视} (jìnshì, myopie).
\end{exercice}

\begin{exercice}[Traduction : Thème (Français $\rightarrow$ Chinois) — Entraînement écrit]
Traduisez les phrases suivantes en chinois (caractères + pinyin). Attention à la place de \zh{应该}, à la structure \zh{太...了}, aux adverbes de fréquence et à la ponctuation chinoise.

\begin{enumerate}
    \item Je pense que les étudiants devraient faire du sport une heure par jour.
    \item Courir tous les matins est trop fatigant.
    \item Le sport est bon pour la santé, mais nous n'avons pas de temps.
    \item Tu as raison, mais le devoir est trop difficile.
    \item Alors, nous pouvons faire du sport le week-end, d'accord ?
\end{enumerate}
\end{exercice}


\begin{corrige}[Exercice de traduction — Corrigé détaillé]
\begin{enumerate}
    \item \zh{我觉得学生应该每天运动一个小时。} (Wǒ juéde xuéshēng yīnggāi měitiān yùndòng \textbf{yí} ge xiǎoshí.) \textit{Note : yī $\rightarrow$ yí avant ton 4 (ge).}
    \item \zh{每天早上跑步太累了。} (Měitiān zǎoshang pǎobù tài lèi le.) \textit{Note : zǎoshang (matin) comme complément de temps en début de phrase.}
    \item \zh{运动对健康很好，可是我们没时间。} (Yùndòng duì jiànkāng hěn hǎo, kěshì wǒmen méi shíjiān.) \textit{Note : « n'avons pas » $\rightarrow$ \zh{没} (méi) + nom, pas \zh{不}. \zh{对...好} structure.}
    \item \zh{你说得对，可是作业太难了。} (Nǐ shuō de duì, kěshì zuòyè tài nán le.) \textit{Note : \zh{说得对} (dire + correctement) = avoir raison.}
    \item \zh{那我们可以周末运动，好吗？} (Nà wǒmen kěyǐ zhōumò yùndòng, hǎo ma?) \textit{Note : \zh{周末} (week-end) comme complément de temps sans préposition. \zh{好吗} pour demander l'accord.}
\end{enumerate}
\end{corrige}


\begin{exercice}[Expression écrite : Rédaction d'un paragraphe d'opinion (80-100 caractères)]
\textbf{Consigne :} Rédigez un court paragraphe (environ 5-6 phrases) pour exprimer votre opinion personnelle sur l'importance du sport pour les lycéens. Utilisez au moins : \zh{应该}, \zh{因为...所以...}, \zh{可是/但是}, \zh{那/那么}, \zh{太...了}. Écrivez les caractères, le pinyin et la traduction française.
\end{exercice}


\begin{corrige}[Modèle de rédaction (exemple de production attendue)]
\textbf{Caractères :}
我觉得中学生应该每天锻炼。因为学习压力大，所以运动很重要。运动能强身健体，也能让我们放松。可是有些同学觉得太累了，不想动。那么，学校应该安排更有趣的体育课。那样，大家都会喜欢运动。

\textbf{Pinyin :}
Wǒ juéde zhōngxuéshēng yīnggāi měitiān duànliàn. Yīnwèi xuéxí yālì dà, suǒyǐ yùndòng hěn zhòngyào. Yùndòng néng qiángshēn jiàntǐ, yě néng ràng wǒmen fàngsōng. Kěshì yǒuxiē tóngxué juéde tài lèi le, bù xiǎng dòng. Nàme, xuéxiào yīnggāi ānpái gèng yǒuqù de tǐyùkè. Nàyàng, dàjiā dōu huì xǐhuan yùndòng.

\textbf{Traduction :}
Je pense que les lycéens devraient s'entraîner tous les jours. Parce que la pression des études est forte, le sport est donc important. Le sport peut renforcer le corps et aussi nous détendre. Mais certains camarades trouvent que c'est trop fatigant et ne veulent pas bouger. Dans ce cas, l'école devrait organiser des cours d'EPS plus amusants. Ainsi, tout le monde aimera le sport.
\end{corrige}


\coursec{Leçon 11 — Expression orale : importance de la pratique du sport}

\courssub{Vocabulaire du sport et de la santé}

\begin{textebox}[Lexique de base — 体育与健康词汇]
\zh{运动 (yùndòng) — sport, exercice physique} \\
\zh{锻炼身体 (duànliàn shēntǐ) — faire de l'exercice, renforcer son corps} \\
\zh{打篮球 (dǎ lánqiú) — jouer au basket} \\
\zh{踢足球 (tī zúqiú) — jouer au football} \\
\zh{游泳 (yóuyǒng) — nager} \\
\zh{跑步 (pǎobù) — courir, footing} \\
\zh{健康 (jiànkāng) — santé, en bonne santé} \\
\zh{强身健体 (qiángshēn jiàntǐ) — renforcer le corps (expr. formelle)} \\
\zh{缓解压力 (huǐjiě yālì) — soulager le stress} \\
\zh{心情放松 (xīnqíng fàngsōng) — humeur détendue} \\
\zh{专心 (zhuānxīn) — concentré, attentif} \\
\zh{平衡 (pínghéng) — équilibre} \\
\zh{自律 (zìlǜ) — autodiscipline} \\
\zh{劳逸结合 (láoyì jiéhé) — alterner travail et repos (chéngyǔ)} \\
\zh{坚持 (jiānchí) — persévérer} \\
\zh{习惯 (xíguàn) — habitude}
\end{textebox}

\begin{center}
\begin{tabularx}{\linewidth}{|X|X|X|X|}
\hline
\rowcolor{popblueL} \textbf{Caractères} & \textbf{Pinyin} & \textbf{Nature} & \textbf{Traduction} \\
\hline
运动 & yùndòng & n./v. & sport / faire du sport \\
\hline
锻炼身体 & duànliàn shēntǐ & v. & faire de l'exercice physique \\
\hline
打篮球 & dǎ lánqiú & v. & jouer au basket \\
\hline
踢足球 & tī zúqiú & v. & jouer au football \\
\hline
游泳 & yóuyǒng & v. & nager \\
\hline
跑步 & pǎobù & v. & courir, faire du footing \\
\hline
健康 & jiànkāng & adj. & en bonne santé \\
\hline
强身健体 & qiángshēn jiàntǐ & v. (expr.) & renforcer le corps \\
\hline
缓解压力 & huǐjiě yālì & v. & soulager le stress \\
\hline
心情放松 & xīnqíng fàngsōng & expr. & humeur détendue \\
\hline
专心 & zhuānxīn & adj. & concentré \\
\hline
平衡 & pínghéng & n./v. & équilibre / équilibrer \\
\hline
自律 & zìlǜ & n./adj. & autodiscipline / autodiscipliné \\
\hline
劳逸结合 & láoyì jiéhé & chéngyǔ & alterner travail et repos \\
\hline
坚持 & jiānchí & v. & persévérer \\
\hline
习惯 & xíguàn & n. & habitude \\
\hline
\end{tabularx}
\end{center}

\courssub{Dialogue modèle 1 : parler de ses sports — 记者采访运动员}

\begin{textebox}[Interview sportive — 校园广播采访]
\zh{记者：你好，我是校园广播的记者。请问你叫什么名字？} \\
\zh{Jìzhě: Nǐ hǎo, wǒ shì xiàoyuán guǎngbō de jìzhě. Qǐngwèn nǐ jiào shénme míngzi?} \\
\zh{Journaliste : Bonjour, je suis journaliste de la radio scolaire. Comment t'appelles-tu ?} \\

\zh{运动员：我叫王强，读高一。} \\
\zh{Yùndòngyuán: Wǒ jiào Wáng Qiáng, dú gāoyī.} \\
\zh{Athlète : Je m'appelle Wang Qiang, je suis en Première.} \\

\zh{记者：王强同学，你最喜欢什么运动？} \\
\zh{Jìzhě: Wáng Qiáng tóngxué, nǐ zuì xǐhuan shénme yùndòng?} \\
\zh{Journaliste : Camarade Wang Qiang, quel est ton sport préféré ?} \\

\zh{运动员：我最喜欢打篮球。因为打篮球可以锻炼身体，还可以交朋友。} \\
\zh{Yùndòngyuán: Wǒ zuì xǐhuan dǎ lánqiú. Yīnwèi dǎ lánqiú kěyǐ duànliàn shēntǐ, hái kěyǐ jiāo péngyou.} \\
\zh{Athlète : J'aime le plus le basket. Parce que le basket permet de faire de l'exercice et aussi de se faire des amis.} \\

\zh{记者：你每周锻炼几次？每次多长时间？} \\
\zh{Jìzhě: Nǐ měi zhōu duànliàn jǐ cì? Měi cì duō cháng shíjiān?} \\
\zh{Journaliste : Combien de fois par semaine t'entraînes-tu ? Combien de temps à chaque fois ?} \\

\zh{运动员：我每周打三次球，每次大概一个半小时。周末有时候打两小时。} \\
\zh{Yùndòngyuán: Wǒ měi zhōu dǎ sān cì qiú, měi cì dàgài yī gè bàn xiǎoshí. Zhōumò yǒushíhou dǎ liǎng xiǎoshí.} \\
\zh{Athlète : Je joue trois fois par semaine, environ une heure et demie chaque fois. Le week-end parfois deux heures.} \\

\zh{记者：你觉得运动对学习有帮助吗？} \\
\zh{Jìzhě: Nǐ juéde yùndòng duì xuéxí yǒu bāngzhù ma?} \\
\zh{Journaliste : Penses-tu que le sport aide les études ?} \\

\zh{运动员：当然有帮助。运动以后，我心情放松，上课更专心。所以我觉得运动很重要。} \\
\zh{Yùndòngyuán: Dāngrán yǒu bāngzhù. Yùndòng yǐhòu, wǒ xīnqíng fàngsōng, shàngkè gèng zhuānxīn. Suǒyǐ wǒ juéde yùndòng hěn zhòngyào.} \\
\zh{Athlète : Bien sûr que ça aide. Après le sport, je suis détendu, je suis plus concentré en classe. Donc je trouve le sport très important.} \\

\zh{记者：最后，给同学们一个建议吧。} \\
\zh{Jìzhě: Zuìhòu, gěi tóngxuémen yī gè jiànyì ba.} \\
\zh{Journaliste : Enfin, un conseil pour les camarades.} \\

\zh{运动员：每天动一动，身体棒棒！谢谢大家。} \\
\zh{Yùndòngyuán: Měitiān dòng yī dòng, shēntǐ bàng bàng! Xièxie dàjiā.} \\
\zh{Athlète : Bougez un peu chaque jour, le corps sera au top ! Merci à tous.}
\end{textebox}

\begin{center}
\begin{tabularx}{\linewidth}{|X|X|X|X|}
\hline
\rowcolor{popblueL} \textbf{Caractères} & \textbf{Pinyin} & \textbf{Nature} & \textbf{Traduction} \\
\hline
记者 & jìzhě & n. & journaliste \\
\hline
运动员 & yùndòngyuán & n. & athlète, sportif \\
\hline
校园广播 & xiàoyuán guǎngbō & n. & radio scolaire \\
\hline
锻炼身体 & duànliàn shēntǐ & v. & faire de l'exercice physique \\
\hline
交朋友 & jiāo péngyou & v. & se faire des amis \\
\hline
大概 & dàgài & adv. & environ, à peu près \\
\hline
心情 & xīnqíng & n. & humeur, état d'esprit \\
\hline
放松 & fàngsōng & adj./v. & détendu / se détendre \\
\hline
专心 & zhuānxīn & adj. & concentré, attentif \\
\hline
建议 & jiànyì & n./v. & conseil / conseiller \\
\hline
棒棒 & bàng bàng & adj. (fam.) & super, au top \\
\hline
\end{tabularx}
\end{center}

\begin{exoresolu}[Transformer l'interview en compte rendu écrit]
\textbf{Énoncé :} À partir du dialogue ci-dessus, écris un court paragraphe (50–60 caractères) en chinois pour le journal mural de la classe. Utilise \zh{据…介绍} (selon l'introduction de…) et \zh{表示} (déclarer/indiquer).

\tcblower
\textbf{Solution détaillée :}
\zh{据记者介绍，王强同学最喜欢打篮球。他每周锻炼三次，每次一个半小时。王强表示，运动能让心情放松，上课更专心，所以运动很重要。他建议同学们每天动一动，身体棒棒。} \\
\zh{Jù jìzhě jièshào, Wáng Qiáng tóngxué zuì xǐhuan dǎ lánqiú. Tā měi zhōu duànliàn sān cì, měi cì yī gè bàn xiǎoshí. Wáng Qiáng biǎoshì, yùndòng néng ràng xīnqíng fàngsōng, shàngkè gèng zhuānxīn, suǒyǐ yùndòng hěn zhòngyào. Tā jiànyì tóngxuémen měitiān dòng yī dòng, shēntǐ bàng bàng.} \\
« Selon le journaliste, l'élève Wang Qiang aime le plus le basket. Il s'entraîne trois fois par semaine, une heure et demie chaque fois. Wang Qiang a déclaré que le sport détend l'humeur, rend plus concentré en classe, donc le sport est très important. Il conseille aux camarades de bouger un peu chaque jour pour avoir un corps au top. »
\end{exoresolu}

\courssub{Formules fonctionnelles : présenter, opiner, relancer}

\begin{aretenir}[Formules fonctionnelles indispensables pour l'oral]
\begin{center}
\begin{tabularx}{\linewidth}{|X|X|X|}
\hline
\rowcolor{popblueL} \textbf{Fonction} & \textbf{Chinois (caractères + pinyin)} & \textbf{Français} \\
\hline
Se présenter & \zh{我叫…，我是高一学生。} (Wǒ jiào…, wǒ shì gāoyī xuésheng.) & Je m'appelle…, je suis en Première. \\
\hline
Présenter le sujet & \zh{今天我们讨论…} (Jīntiān wǒmen tǎolùn…) & Aujourd'hui nous discutons de… \\
\hline
Demander l'avis & \zh{你觉得怎么样？ / 你怎么看？} (Nǐ juéde zěnmeyàng? / Nǐ zěnme kàn?) & Qu'en penses-tu ? / Comment tu vois ça ? \\
\hline
Donner son avis & \zh{我认为… / 我觉得… / 按我看…} (Wǒ rènwéi… / Wǒ juéde… / Àn wǒ kàn…) & Je pense que… / À mon avis… \\
\hline
Accord & \zh{我赞同。/ 有道理。/ 说得好。} (Wǒ zàntóng. / Yǒu dàolǐ. / Shuō de hǎo.) & Je suis d'accord. / C'est juste. / Bien dit. \\
\hline
Désaccord poli & \zh{我不同意。/ 这不完全对。/ 我有不同看法。} (Wǒ bù tóngyì. / Zhè bù wánquán duì. / Wǒ yǒu bùtóng kànfǎ.) & Je ne suis pas d'accord. / Ce n'est pas tout à fait exact. / J'ai un autre avis. \\
\hline
Argumenter (cause-effet) & \zh{因为…，所以…} (Yīnwèi…, suǒyǐ…) & Parce que…, donc… \\
\hline
Concession & \zh{虽然…，但是…} (Suīrán…, dànshì…) & Bien que…, mais… \\
\hline
Relancer / Passer la parole & \zh{那么…呢？ / 请另一方回应。} (Nàme… ne? / Qǐng lìng yīfāng huíyìng.) & Et… ? / Que l'autre partie réponde. \\
\hline
Gagner du temps (fillers) & \zh{嗯… / 这个… / 让我想想…} (Èn… / Zhège… / Ràng wǒ xiǎngxiǎng…) & Euh… / Alors… / Laisse-moi réfléchir… \\
\hline
Conclure & \zh{所以，我们认为… / 综上所述…} (Suǒyǐ, wǒmen rènwéi… / Zōngshàng shùshù…) & Donc, nous pensons que… / En résumé… \\
\hline
\end{tabularx}
\end{center}
\end{aretenir}

\courssub{Travail sur les tons et la prononciation}

\begin{propriete}[Rappel des quatre tons et du ton neutre]
Le chinois standard (putonghua) possède 4 tons pleins + 1 ton neutre. La maîtrise des tons est \textbf{obligatoire} pour être compris à l'oral.
\begin{itemize}
\item \textbf{1er ton (─) :} haut et plat. Ex. \zh{天} (tiān, ciel), \zh{高} (gāo, haut), \zh{书} (shū, livre).
\item \textbf{2e ton (/) :} montant. Ex. \zh{人} (rén), \zh{来} (lái).
\item \textbf{3e ton (V) :} basculant (bas → haut). Ex. \zh{我} (wǒ), \zh{好} (hǎo). \textbf{Sandhi :} 3e + 3e → 2e + 3e (\zh{你好} → ní hǎo).
\item \textbf{4e ton (\textbackslash{}) :} tombant fort. Ex. \zh{是} (shì), \zh{对} (duì).
\item \textbf{Ton neutre (·) :} bref, sans accent, suit la syllabe précédente. Ex. \zh{妈妈} (māma), \zh{棒棒} (bàng·bang).
\end{itemize}
\end{propriete}

\begin{attention}[Pièges tonaux fréquents pour francophones (mots-clés du thème)]
\begin{itemize}
\item \zh{重要} (zhòngyào) : 4e + 4e. Ne pas dire *zhóngyáo (2e+2e).
\item \zh{锻炼} (duànliàn) : 4e + 4e. Attention au \zh{liàn} (4e) vs \zh{lián} (2e, 连).
\item \zh{平衡} (pínghéng) : 2e + 2e. Ne pas confondre avec \zh{平行} (píngxíng, 2e+2e).
\item \zh{自律} (zìlǜ) : 4e + 4e. Le \zh{lü} (ü) s'arrondit bien.
\item \zh{缓解} (huǐjiě) : 3e + 3e → \textbf{huí jiě} (2e + 3e) par sandhi.
\item \zh{挤时间} (jǐ shíjiān) : 3e + 2e + 1e. \zh{jǐ} (3e) bascule bas.
\item \zh{劳逸结合} (láoyì jiéhé) : 2e + 4e + 2e + 2e.
\end{itemize}
\end{attention}

\begin{experience}[Exercices de prononciation guidée (10 min)]
\begin{enumerate}
\item \textbf{Écoute-discrimination :} L'enseignant lit paires minimales ; élèves lèvent carte « Même / Différent ».
\begin{itemize}
\item \zh{重要} vs \zh{种植} (zhòngzhí)
\item \zh{锻炼} vs \zh{短练} (duǎnliàn)
\item \zh{平衡} vs \zh{凭恒} (pínghéng - homophones)
\item \zh{自律} vs \zh{字律} (zìlǜ)
\end{itemize}
\item \textbf{Répétition chorale puis individuelle} des phrases-clés du dialogue 1 :
\zh{因为打篮球可以锻炼身体} (Yīnwèi dǎ lánqiú kěyǐ duànliàn shēntǐ) — insister sur \zh{duànliàn}.
\zh{运动能让心情放松} (Yùndòng néng ràng xīnqíng fàngsōng) — \zh{fàngsōng} (4e+1e).
\item \textbf{Lecture à vue du dialogue 2} (extrait radio) en binôme : l'un lit l'animateur, l'autre les invités. L'enseignant circule et corrige les tons 3/4.
\end{enumerate}
\end{experience}

\courssub{Dialogue modèle 2 : mini-débat « Faut-il faire du sport tous les jours ? » — 广播辩论}

\begin{textebox}[Extrait type de débat radio — 喀麦隆之声 校园版]
\zh{主持人：亲爱的听众朋友们，欢迎收听“喀麦隆之声”校园版。今天我们讨论一个热门话题：运动重要，还是学习重要？请两位嘉宾先自我介绍。} \\
\zh{Zhǔchírén: Qīn'ài de tīngzhòng péngyoumen, huānyíng shōutīng "Kāmàilóng zhī shēng" xiàoyuán bǎn. Jīntiān wǒmen tǎolùn yī gè rèmén huàtí: yùndòng zhòngyào, háishì xuéxí zhòngyào? Qǐng liǎng wèi jiābīn xiān zìwǒ jièshào.} \\
\zh{Animateur : Chers auditeurs, bienvenue sur "Voix du Cameroun" version scolaire. Aujourd'hui nous débattons d'un sujet brûlant : le sport ou les études, lequel est plus important ? Que les deux invités se présentent.} \\

\zh{嘉宾A（赞成运动）：我是高一的李明。我认为运动更重要。因为身体是革命的本钱。没有健康的身体，怎么学习？运动还能缓解压力，让我们交到好朋友。} \\
\zh{Jiābīn A (zànchéng yùndòng): Wǒ shì gāoyī de Lǐ Míng. Wǒ rènwéi yùndòng gèng zhòngyào. Yīnwèi shēntǐ shì gémìng de běnqián. Méiyǒu jiànkāng de shēntǐ, zěnme xuéxí? Yùndòng hái néng huǐjiě yālì, ràng wǒmen jiāo dào hǎo péngyou.} \\
\zh{Invité A (pour le sport) : Je suis Li Ming de Première. Je pense que le sport est plus important. Parce que le corps est le capital de la révolution. Sans corps en bonne santé, comment étudier ? Le sport soulage aussi le stress et nous permet de se faire de bons amis.} \\

\zh{嘉宾B（偏向学习）：我是高一的陈静。我赞同运动有好处，但学习是学生的本分。高一课程重，作业多，如果每天花两小时运动，复习就不够了。考试考砸了，以后上大学很难。} \\
\zh{Jiābīn B (piānxiàng xuéxí): Wǒ shì gāoyī de Chén Jìng. Wǒ zàntóng yùndòng yǒu hǎochù, dàn xuéxí shì xuésheng de fènfèn. Gāoyī kèchéng zhòng, zuòyè duō, rúguǒ měitiān huā liǎng xiǎoshí yùndòng, fùxí jiù bù gòu le. Kǎoshì kǎozá le, yǐhòu shàng dàxué hěn nán.} \\
\zh{Invité B (penché études) : Je suis Chen Jing de Première. Je suis d'accord que le sport a des avantages, mais les études sont le devoir de l'élève. En Première, le programme est lourd, beaucoup de devoirs ; si on passe deux heures par jour au sport, la révision ne suffit plus. Si on rate les examens, entrer à l'université sera difficile.} \\

\zh{主持人：两位观点都很有道理。李明同学，你怎么回应陈静说的“时间不够”？} \\
\zh{Zhǔchírén: Liǎng wèi guāndiǎn dōu hěn yǒu dàolǐ. Lǐ Míng tóngxué, nǐ zěnme huíyìng Chén Jìng shuō de "shíjiān bù gòu"?} \\
\zh{Animateur : Les deux points de vue sont très fondés. Camarade Li Ming, comment réponds-tu à Chen Jing qui dit « pas assez de temps » ?} \\

\zh{嘉宾A：其实不需要两小时。每天三十分钟慢跑或打球，效率反而更高。因为运动后大脑清醒，记单词更快。所以“挤时间”也是一种方法。} \\
\zh{Jiābīn A: Qíshí bù xūyào liǎng xiǎoshí. Měitiān sānshí fēnzhōng mànpǎo huò dǎ qiú, xiàolǜ fǎn'ér gèng gāo. Yīnwèi yùndòng hòu dànǎo qīngxǐng, jì dāncí gèng kuài. Suǒyǐ "jǐ shíjiān" yě shì yī zhǒng fāngfǎ.} \\
\zh{Invité A : En fait, il ne faut pas deux heures. Trente minutes de footing ou de ballon par jour, l'efficacité est au contraire plus grande. Parce qu'après le sport le cerveau est clair, on mémorise les mots plus vite. Donc « se faire du temps » est aussi une méthode.} \\

\zh{主持人：陈静同学，你同意“挤时间”吗？} \\
\zh{Zhǔchírén: Chén Jìng tóngxué, nǐ tóngyì "jǐ shíjiān" ma?} \\
\zh{Animateur : Camarade Chen Jing, es-tu d'accord avec « se faire du temps » ?} \\

\zh{嘉宾B：部分同意。但要看人。有的同学自制力差，一玩就停不下来。建议学校安排固定的体育课，大家一起动，比较公平、也安全。} \\
\zh{Jiābīn B: Bùfēn tóngyì. Dàn yào kàn rén. Yǒu de tóngxué zìzhìlì chà, yī wán jiù tíng bù xiàlái. Jiànyì xuéxiào ānpái gùdìng de tǐyùkè, dàjiā yīqǐ dòng, bǐjiào gōngpíng, yě ānquán.} \\
\zh{Invité B : Partiellement d'accord. Mais ça dépend des personnes. Certains élèves manquent d'autodiscipline, une fois qu'ils jouent ils ne s'arrêtent plus. Je suggère que l'école organise des cours d'EPS fixes, tout le monde bouge ensemble, c'est plus équitable et plus sûr.} \\

\zh{主持人：谢谢两位精彩的交锋。请报告员总结。} \\
\zh{Zhǔchírén: Xièxie liǎng wèi jīngcǎi de jiāofēng. Qǐng bàogàoyuán zǒngjié.} \\
\zh{Animateur : Merci aux deux pour cet échange passionnant. Rapporteur, concluons.} \\

\zh{报告员：通过讨论，我们认为：运动和学习不是对立的，而是互相促进。关键是“平衡”和“自律”。所以，我们建议：每天适量运动，合理安排学习时间，劳逸结合。} \\
\zh{Bàogàoyuán: Tōngguò tǎolùn, wǒmen rènwéi: yùndòng hé xuéxí bùshì duìlì de, érshì hùxiāng cùjìn. Guānjiàn shì "pínghéng" hé "zìlǜ". Suǒyǐ, wǒmen jiànyì: měitiān shìliàng yùndòng, hélǐ ānpái xuéxí shíjiān, láoyì jiéhé.} \\
\zh{Rapporteur : À travers la discussion, nous pensons : le sport et les études ne sont pas opposés, ils se favorisent mutuellement. L'essentiel est « l'équilibre » et « l'autodiscipline ». Donc nous conseillons : sport modéré chaque jour, organiser raisonnablement le temps d'étude, alterner travail et repos.}
\end{textebox}

\begin{center}
\begin{tabularx}{\linewidth}{|X|X|X|X|}
\hline
\rowcolor{popblueL} \textbf{Caractères} & \textbf{Pinyin} & \textbf{Nature} & \textbf{Traduction} \\
\hline
主持人 & zhǔchírén & n. & animateur, présentateur \\
\hline
听众 & tīngzhòng & n. & auditeurs \\
\hline
嘉宾 & jiābīn & n. & invité (émission) \\
\hline
热门话题 & rèmén huàtí & n. & sujet brûlant, thème chaud \\
\hline
革命的本钱 & gémìng de běnqián & expr. & le capital de la révolution (santé = base) \\
\hline
缓解压力 & huǐjiě yālì & v. & soulager le stress \\
\hline
本分 & fènfèn & n. & devoir, rôle propre \\
\hline
考砸 & kǎozá & v. (fam.) & rater un examen \\
\hline
效率 & xiàolǜ & n. & efficacité \\
\hline
大脑清醒 & dànǎo qīngxǐng & expr. & esprit clair, tête fraîche \\
\hline
挤时间 & jǐ shíjiān & v. & se faire du temps, caser \\
\hline
自制力 & zìzhìlì & n. & autodiscipline, maîtrise de soi \\
\hline
停不下来 & tíng bù xiàlái & v. & ne pas pouvoir s'arrêter \\
\hline
体育课 & tǐyùkè & n. & cours d'EPS \\
\hline
公平 & gōngpíng & adj. & équitable, juste \\
\hline
对立 & duìlì & v./adj. & s'opposer / opposé \\
\hline
互相促进 & hùxiāng cùjìn & v. & se favoriser mutuellement \\
\hline
平衡 & pínghéng & n./v. & équilibre / équilibrer \\
\hline
自律 & zìlǜ & n./adj. & autodiscipline / autodiscipliné \\
\hline
劳逸结合 & láoyì jiéhé & chéngyǔ & alterner travail et repos \\
\hline
\end{tabularx}
\end{center}

\courssub{Jeux de rôle et débat guidé en classe}

\begin{methode}[Déroulement de l'interview — Jeu de rôle 1]
\begin{enumerate}
\item \textbf{Préparation (5 min) :} par binôme, l'élève A (journaliste) prépare 4 questions sur une fiche \emph{mots-clés seulement} (ex. : sport ? fréquence ? pourquoi important ? conseil ?). L'élève B (champion) prépare 3 arguments avec \zh{因为…所以…}.
\item \textbf{Enregistrement / présentation (3 min) :} l'interview se déroule debout, sans lire les notes. Le journaliste utilise les formules de relance : \zh{请问…}，\zh{那么…呢？}，\zh{您觉得呢？}.
\item \textbf{Retour classe (2 min) :} deux binômes volontaires rejouent devant la classe ; l'enseignant note sur la grille d'observation (voir plus bas).
\end{enumerate}
\end{methode}

\begin{experience}[Mise en place du débat radio — Jeu de rôle 2]
\begin{enumerate}
\item \textbf{Distribution des rôles (4 élèves) :}
\begin{itemize}
\item \textbf{Animateur(trice) :} lance le sujet, relance avec \zh{你觉得怎么样？}，\zh{还有谁想补充？}，\zh{我们听听另一方的观点。}
\item \textbf{Invité POUR (运动更重要) :} 3 arguments : santé, détente, esprit d'équipe.
\item \textbf{Invité CONTRE / NUANCÉ (学习为主) :} 3 arguments : temps limité, fatigue, priorité aux examens.
\item \textbf{Rapporteur :} prend des notes, conclut avec \zh{所以，我们认为…} ou \zh{综上所述…}.
\end{itemize}
\item \textbf{Temps de parole :} 1 min par argument, 30 s pour la conclusion du rapporteur. Total : ~10 min.
\item \textbf{Public :} les autres élèves notent sur une fiche « Argument POUR / Argument CONTRE / Formule entendue ».
\end{enumerate}
\end{experience}

\begin{methode}[Réussir un débat oral en chinois — 5 étapes]
\begin{enumerate}
\item \textbf{Préparer 3 arguments structurés} avec \zh{因为…，所以…} (car… donc…). Ex. : \zh{因为运动能强身健体，所以我们应该每天锻炼。}
\item \textbf{Écouter activement} l'adversaire : noter un mot-clé de son argument pour y répondre précisément (pas de hors-sujet).
\item \textbf{Utiliser au moins une formule d'accord} (\zh{我赞同…}，\zh{有道理}) \textbf{et une de désaccord} (\zh{我不同意…}，\zh{这不完全对}) par intervention.
\item \textbf{Relancer le débat} si tu es animateur : \zh{你怎么看？}，\zh{还有不同意见吗？}，\zh{请另一方回应。}
\item \textbf{Conclure collectivement} : le rapporteur synthétise avec \zh{所以，我们认为…} ou \zh{综上所述，…}.
\end{enumerate}
\end{methode}

\begin{attention}[Pièges fréquents des francophones en débat oral]
\begin{itemize}
\item \textbf{Lire ses notes mot à mot} → \emph{Interdit}. Prépare une fiche \textbf{3–4 mots-clés chinois seulement} (ex. : \zh{健康、交友、缓解压力、自律}). Parle \emph{à partir} de ces mots.
\item \textbf{Oublier les tons} sous le stress → répète les 3 arguments \emph{à voix haute} avant le débat ; concentre-toi sur les tons 3 (basculant) et 4 (tombant) des mots-clés : \zh{重\underline{要} (zhòngyào)}, \zh{锻\underline{炼} (duànliàn)}, \zh{平\underline{衡} (pínghéng)}.
\item \textbf{Ordre des mots français} : \emph{« Je pense que le sport est important »} $\neq$ \zh{我想运动重要} (faux). Structure correcte : \zh{我认为运动很重要} (Sujet + 认为 + Objet + 很 + Adj).
\item \textbf{Confondre 得 / 的 / 地} : \zh{跑得快} (courir vite — complément de degré), \zh{快的跑} (inexistant), \zh{快地跑} (vite courir — adverbial rare). En débat : \zh{运动能让我们健康地成长} (grandir sainement).
\item \textbf{Ne pas céder la parole} : attends la fin de la phrase de l'autre (marqueur \zh{。} ou intonation descendante) avant de prendre la parole avec \zh{那个…} ou \zh{我想补充一点}.
\end{itemize}
\end{attention}

\begin{exemplebox}[Utilisation des « mots-outils de remplissage » pour gagner du temps]
En conversation réelle, les Chinois utilisent des particules pour « tenir le plancher » le temps de formuler sa pensée. Apprends-les pour ne pas basculer en français (\emph{euh…}, \emph{ben…}) :

\begin{center}
\begin{tabularx}{\linewidth}{|X|X|X|}
\hline
\rowcolor{popblueL} \textbf{Particule} & \textbf{Pinyin} & \textbf{Usage + Exemple} \\
\hline
\zh{嗯} & èn (ton 4) & Accord / réflexion : \zh{嗯… 我觉得也是。} (Èn… wǒ juéde yě shì. — « Euh… je pense aussi. ») \\
\hline
\zh{这个…} & zhège… & « Ça… / Alors… » : \zh{这个… 我们可以这样说…} (Zhège… wǒmen kěyǐ zhèyàng shuō… — « Alors… on peut dire ça… ») \\
\hline
\zh{然后…} & ránhòu… & « Et puis… / Ensuite… » pour enchaîner : \zh{然后，还有健康的问题。} \\
\hline
\zh{就是…} & jiùshì… & « C'est-à-dire… / En fait… » pour reformuler : \zh{就是… 运动不一定要很久。} \\
\hline
\end{tabularx}
\end{center}
\textbf{Exercice :} Réécris la phrase « Je pense que le sport c'est bien mais j'ai pas le temps » en insérant \zh{这个…} et \zh{嗯…} pour faire naturel. \\
\textbf{Réponse :} \zh{嗯… 这个… 我觉得运动挺好，但是没有时间。} (Èn… zhège… wǒ juéde yùndòng tǐng hǎo, dànshì méiyǒu shíjiān.)
\end{exemplebox}

\courssub{Carte des rôles — Schéma TikZ}

\begin{popfigure}
\begin{tikzpicture}[
  node distance=2.2cm and 1.8cm,
  role/.style={draw=popblue, thick, rounded corners=6pt, fill=popblueL, minimum width=2.8cm, minimum height=1.4cm, align=center, font=\small},
  arrow/.style={-Stealth, thick, popdark},
  labelstyle/.style={font=\tiny, text=popmuted, align=center}
]
% Nœud central : Animateur
\node[role, align=center] (anim){Animateur\\主持人};
\node[labelstyle, below=0.1cm of anim, align=center] (animlab){Relance : \zh{你觉得怎么样？}\\ \zh{还有谁想补充？}};

% Camp POUR
\node[role, left=of anim, xshift=-1.5cm, fill=popgreenL, draw=popgreen, align=center] (pour){POUR\\运动很重要};
\node[labelstyle, below=0.1cm of pour, align=center] (pourlab){Args : \zh{健康、放松、交友}\\Formule : \zh{因为…所以…}};

% Camp CONTRE
\node[role, right=of anim, xshift=1.5cm, fill=poporangeL, draw=poporange, align=center] (contre){CONTRE / NUANCÉ\\学习为主};
\node[labelstyle, below=0.1cm of contre, align=center] (contrelab){Args : \zh{时间、累、考试}\\Formule : \zh{虽然…但是…}};

% Rapporteur
\node[role, below=of anim, yshift=-0.8cm, fill=poppurpleL, draw=poppurple, align=center] (rap){Rapporteur\\报告员};
\node[labelstyle, below=0.1cm of rap, align=center] (raplab){Conclut : \zh{所以，我们认为…}\\ \zh{综上所述…}};

% Flèches
\draw[arrow] (anim) -- node[above, labelstyle] {lance / relance} (pour);
\draw[arrow] (anim) -- node[above, labelstyle] {lance / relance} (contre);
\draw[arrow] (pour) -- node[left, labelstyle] {argument} (rap);
\draw[arrow] (contre) -- node[right, labelstyle] {argument} (rap);
\draw[arrow, dashed] (rap) -- node[below, labelstyle] {synthèse finale} (anim);

% Public
\node[draw=poppink, dashed, thick, rounded corners=8pt, fit=(pour) (contre) (rap), inner sep=6pt, label={[font=\tiny, text=poppink]above:Public — note arguments et formules}] {};
\end{tikzpicture}
\legende{Carte des rôles pour le débat guidé : l'animateur au centre distribue la parole vers les deux camps (POUR / CONTRE) et reçoit la synthèse du rapporteur.}
\end{popfigure}

\courssub{Grille d'observation de l'enseignant (à remplir pendant les jeux de rôle)}

\begin{center}
\begin{tabularx}{\linewidth}{|X|X|X|X|X|}
\hline
\rowcolor{popblueL} \textbf{Critère} & \textbf{4 — Excellent} & \textbf{3 — Bon} & \textbf{2 — Passable} & \textbf{1 — Insuffisant} \\
\hline
Prononciation des tons & Tons 1–4 justes, naturel & 1–2 erreurs non gênantes & Erreurs fréquentes, compréhensible & Incompréhensible (tons plats) \\
\hline
Formules fonctionnelles & $\ge$ 4 formules variées & 3 formules correctes & 1–2 formules, répétitives & Aucune formule, français seul \\
\hline
Fluidité / débit & Débit naturel, pauses logiques & Quelques hésitations & Hésitations longues, « euh » français & Lecture mot à mot, très lent \\
\hline
Respect du tour de parole & Écoute active, relance l'autre & Attend son tour, répond au sujet & Coupe parfois la parole & Monologue, ignore les autres \\
\hline
Contenu / arguments & 3 arguments + exemple perso & 2 arguments clairs & 1 argument vague & Hors sujet / silence \\
\hline
\end{tabularx}
\end{center}

\courssub{Exercices de consolidation}

\begin{exercice}[Préparation fiche mots-clés]
Pour le débat « Le sport est-il plus important que les études ? », prépare \textbf{deux} fiches (recto-verso) de 3–4 mots-clés chinois chacune :
\begin{enumerate}
\item Fiche CAMP POUR (arguments santé, détente, amitié).
\item Fiche CAMP CONTRE/NUANCÉ (arguments temps, fatigue, priorité études).
\end{enumerate}
Écris uniquement les mots-clés en caractères (pas de phrases). Échange avec ton voisin et fais-lui deviner ton camp à partir des mots.
\end{exercice}

\begin{corrige}[Exercice 1]
\textbf{Fiche POUR (exemple) :} \zh{强身健体、缓解压力、交朋友、劳逸结合} \\
\textbf{Fiche CONTRE (exemple) :} \zh{时间不够、太累、考试重要、自制力差} \\
\emph{Variantes acceptées :} \zh{健康、放松、团队、平衡} / \zh{作业多、复习、上大学、自律}.
\end{corrige}

\begin{exercice}[Mini-débat éclair — 3 minutes]
En binôme, tirez un sort : l'un est POUR, l'autre CONTRE. Sujet : \zh{中学生每天必须锻炼一小时} (« Les lycéens doivent s'entraîner 1 h par jour »). Utilise \textbf{au moins} : une formule d'accord, une de désaccord, \zh{因为…所以…}, et un mot de remplissage (\zh{嗯} ou \zh{这个}). L'enseignant chronomètre et note sur la grille.
\end{exercice}

\begin{corrige}[Exercice 2 — Exemple de trace orale attendue]
\zh{A: 嗯… 我认为中学生每天必须锻炼一小时。因为身体是革命的本钱，所以健康最重要。} \\
\zh{B: 这个… 我不同意。虽然健康重要，但高一作业太多，一小时太长。挤时间可以，但不一定每天。} \\
\zh{A: 有道理。那你觉得多少时间合适？} \\
\zh{B: 三十分钟吧。既锻炼了，又不耽误复习。所以，我们认为：适量运动，平衡最重要。}
\end{corrige}

\begin{exercice}[Expression écrite : compte rendu de débat]
Rédige en chinois (80–100 caractères) le compte rendu du mini-débat ci-dessus pour l'affichage en classe. Structure : \zh{主题} — \zh{正方观点} — \zh{反方观点} — \zh{共识}. Utilise \zh{认为}、\zh{表示}、\zh{建议}.
\end{exercice}

\begin{corrige}[Exercice 3]
\zh{主题：中学生每天必须锻炼一小时。正方认为：身体是本钱，运动能缓解压力，所以必须锻炼。反方表示：作业重，时间不够，一小时太长，建议三十分钟。双方达成共识：适量运动，劳逸结合，平衡最重要。}
\end{corrige}

\begin{aretenir}[Check-list avant l'oral noté (bac / contrôle continu)]
\begin{itemize}
\item[☐] Je connais mes 3 arguments \emph{avec} \zh{因为…所以…}.
\item[☐] J'ai 2 formules d'accord et 2 de désaccord prêtes.
\item[☐] Je sais placer \zh{嗯…} / \zh{这个…} sans passer au français.
\item[☐] Mes tons sur les mots-clés (\zh{重要}、\zh{锻炼}、\zh{平衡}、\zh{自律}) sont stables.
\item[☐] Ma fiche mots-clés tient sur un post-it (max 4 mots par camp).
\item[☐] Je regarde mes camarades quand je parle (pas mes chaussures).
\end{itemize}
\end{aretenir}

\begin{savaistu}[Culture sportive : Cameroun vs Chine]
Au Cameroun, les \textbf{journées sportives inter-écoles} (souvent en mars-avril) et les matchs de quartier — le fameux football « bœuf » (pieds nus, ballon en chiffons, buts en pierres) — sont des moments forts de la vie des jeunes. Ils forgent l'esprit d'équipe et la résilience, exactement comme les \textbf{compétitions de 乒乓球 (pīngpāngqiú, tennis de table)} dans les écoles chinoises, où le ping-pong est sport national, enseigné dès le primaire, avec des entraînements quotidiens de 30 min. Dans les deux pays, le sport scolaire n'est pas une option : c'est un \emph{rite de passage} et un vecteur de cohésion sociale.
\end{savaistu}

\newpage

\coursec{Activité d'intégration}

\courssub{Objectifs de l'activité}

À la fin de cette activité d'intégration, tu seras capable de :

\begin{itemize}
\item mobiliser dans une situation orale authentique le \cle{vocabulaire du sport et de la santé} étudié dans la leçon ;
\item utiliser les \cle{formules fonctionnelles} pour présenter un sujet, donner son avis, réagir et conclure ;
\item prononcer correctement les \cle{tons} dans des phrases complètes, et pas seulement dans des mots isolés ;
\item respecter le \cle{tour de parole} dans un débat organisé ;
\item écouter activement tes camarades pour voter pour l'argument le plus convaincant.
\end{itemize}

\courssub{Situation}

Le lycée de Yaoundé organise la \emph{Semaine de la santé}. La radio scolaire, « La Voix du Lycée », a décidé d'enregistrer une émission spéciale \textbf{en chinois} pour les élèves de la section chinois. Le thème choisi est : \zh{运动重要吗？} (\emph{Yùndòng zhòngyào ma?} « Le sport, est-ce important ? »).

Tu es invité(e) à participer à cette émission. La classe est divisée en quatre groupes : un(e) \textbf{animateur/animatrice} qui présente le sujet et donne la parole, une \textbf{équipe POUR} (le sport est très important), une \textbf{équipe NUANCÉE} (le sport est important, mais pas tous les jours), et un(e) \textbf{rapporteur/rapporteuse} qui conclut l'émission. Le reste de la classe forme le public et vote à la fin.

Voici le script d'ouverture que l'animateur peut lire ou adapter :

\begin{textebox}[Script d'ouverture de l'émission]
大家好！欢迎收听《校园之声》。\\
\emph{Dàjiā hǎo! Huānyíng shōutīng « Xiàoyuán zhī Shēng ».}\\
« Bonjour à tous ! Bienvenue à l'écoute de "La Voix du Campus". »\\[4pt]
今天我们要谈一谈运动的重要性。\\
\emph{Jīntiān wǒmen yào tán yi tan yùndòng de zhòngyàoxìng.}\\
« Aujourd'hui, nous allons parler de l'importance du sport. »\\[4pt]
运动重要吗？应该每天运动吗？\\
\emph{Yùndòng zhòngyào ma? Yīnggāi měitiān yùndòng ma?}\\
« Le sport, est-ce important ? Faut-il faire du sport tous les jours ? »\\[4pt]
我们先请第一组的同学说一说。\\
\emph{Wǒmen xiān qǐng dì-yī zǔ de tóngxué shuō yi shuō.}\\
« Nous invitons d'abord les élèves du premier groupe à prendre la parole. »
\end{textebox}

\courssub{Consignes}

\begin{enumerate}
\item \textbf{Préparation en équipes (10 min).} Dans ton équipe, relis le vocabulaire du sport et choisis \textbf{trois arguments} pour ta position. Écris chaque argument en une phrase chinoise courte, avec le pinyin au-dessus si nécessaire. Exemple d'argument POUR : \zh{运动对身体很好。} (\emph{Yùndòng duì shēntǐ hěn hǎo.} « Le sport est très bon pour la santé. »).

\item \textbf{Distribution des rôles (5 min).} Chaque équipe choisit deux porte-parole. L'animateur et le rapporteur répètent leur script. Vérifie la prononciation des tons avec ton voisin : lis ta phrase à voix haute, ton voisin corrige.

\item \textbf{L'émission (15 min).} L'animateur ouvre le débat avec le script. Chaque équipe présente ses arguments à tour de rôle. Règle du jeu : chaque porte-parole doit utiliser \textbf{au moins trois formules fonctionnelles} différentes, par exemple \zh{我觉得…}, \zh{我同意。}, \zh{你说得对，可是…}, \zh{那我们可以…}. On ne coupe pas la parole : pour réagir, on commence par une formule de réaction.

\item \textbf{Conclusion (5 min).} Le rapporteur résume les deux positions et conclut, par exemple : \zh{所以，我们认为运动对每个人都很重要。谢谢大家！} (\emph{Suǒyǐ, wǒmen rènwéi yùndòng duì měi ge rén dōu hěn zhòngyào. Xièxie dàjiā!} « Donc, nous pensons que le sport est très important pour chacun. Merci à tous ! »).

\item \textbf{Vote du public (5 min).} Le public vote à main levée pour \textbf{l'argument le plus convaincant} et justifie son vote avec une phrase en chinois : \zh{我觉得这个理由很好，因为…} (\emph{Wǒ juéde zhège lǐyóu hěn hǎo, yīnwèi…} « Je trouve que cet argument est bon, parce que… »).

\item \textbf{Auto-évaluation.} À l'aide de la grille d'observation, note ta propre performance : ai-je bien prononcé les tons ? ai-je utilisé trois formules ? ai-je parlé sans lire ?
\end{enumerate}

\courssub{Ressources à mobiliser}

\textbf{Vocabulaire utile.}

\begin{center}
\begin{tabularx}{\linewidth}{|X|X|X|X|}
\hline
\rowcolor{popblueL} Caractères & Pinyin & Nature & Traduction \\
\hline
运动 & yùndòng & n. / v. & sport ; faire du sport \\
\hline
重要 & zhòngyào & adj. & important \\
\hline
身体 & shēntǐ & n. & corps, santé \\
\hline
健康 & jiànkāng & adj. / n. & sain ; santé \\
\hline
每天 & měitiān & expr. temp. & tous les jours \\
\hline
跑步 & pǎobù & v. & courir \\
\hline
踢足球 & tī zúqiú & v. + n. & jouer au football \\
\hline
游泳 & yóuyǒng & v. & nager \\
\hline
累 & lèi & adj. & fatigué \\
\hline
应该 & yīnggāi & v. modal & devoir (il faut) \\
\hline
理由 & lǐyóu & n. & raison, argument \\
\hline
\end{tabularx}
\end{center}

\textbf{Formules fonctionnelles du débat.}

\begin{center}
\begin{tabularx}{\linewidth}{|X|X|X|}
\hline
\rowcolor{popblueL} Fonction & Formule (caractères + pinyin) & Traduction \\
\hline
Présenter le sujet & 今天我们要谈一谈… \emph{Jīntiān wǒmen yào tán yi tan…} & Aujourd'hui, nous allons parler de… \\
\hline
Donner son avis & 我觉得… \emph{Wǒ juéde…} & Je pense que… / Je trouve que… \\
\hline
Être d'accord & 我同意。 \emph{Wǒ tóngyì.} & Je suis d'accord. \\
\hline
Ne pas être d'accord & 我不同意。 \emph{Wǒ bù tóngyì.} & Je ne suis pas d'accord. \\
\hline
Nuancer & 你说得对，可是… \emph{Nǐ shuō de duì, kěshì…} & Tu as raison, mais… \\
\hline
Proposer & 那我们可以… \emph{Nà wǒmen kěyǐ…} & Alors nous pouvons… \\
\hline
Conclure & 所以，我们认为… \emph{Suǒyǐ, wǒmen rènwéi…} & Donc, nous pensons que… \\
\hline
\end{tabularx}
\end{center}

\begin{attention}[Pièges de prononciation à surveiller]
\begin{itemize}
\item \zh{运动} \emph{yùndòng} : deux quatrièmes tons descendants, bien marquer la chute des deux syllabes ; ne pas confondre avec \emph{yǔndòng}.
\item \zh{重要} \emph{zhòngyào} : deux quatrièmes tons ; les francophones ont tendance à adoucir le second ton.
\item \zh{身体} \emph{shēntǐ} : premier ton haut et plat sur \emph{shēn}, troisième ton sur \emph{tǐ}.
\item \zh{我觉得} \emph{wǒ juéde} : troisième + deuxième ton ; enchaîner sans pause.
\item Dans \zh{你说得对，可是…}, attention à \emph{de} (ton neutre) après le verbe \emph{shuō} : ne pas le prononcer fort.
\end{itemize}
\end{attention}

\courssub{Proposition de production et corrigé commenté}

\begin{exoresolu}[Modèle d'émission radiophonique complète]
Énoncé : rédige et joue une émission de débat d'environ dix répliques sur le thème \zh{运动重要吗？}, avec un animateur, deux équipes et un rapporteur.
\tcblower
\textbf{Production modèle.}

\textbf{Animateur :} 大家好！欢迎收听《校园之声》。今天我们要谈一谈运动的重要性。运动重要吗？应该每天运动吗？我们先请第一组的同学说一说。\\
\emph{Dàjiā hǎo! Huānyíng shōutīng « Xiàoyuán zhī Shēng ». Jīntiān wǒmen yào tán yi tan yùndòng de zhòngyàoxìng. Yùndòng zhòngyào ma? Yīnggāi měitiān yùndòng ma? Wǒmen xiān qǐng dì-yī zǔ de tóngxué shuō yi shuō.}\\
« Bonjour à tous ! Bienvenue dans "La Voix du Campus". Aujourd'hui, nous allons parler de l'importance du sport. Le sport, est-ce important ? Faut-il faire du sport tous les jours ? Nous invitons d'abord le premier groupe à parler. »

\textbf{Équipe POUR (élève A) :} 我觉得运动非常重要。运动对身体很好，也对学习很好。\\
\emph{Wǒ juéde yùndòng fēicháng zhòngyào. Yùndòng duì shēntǐ hěn hǎo, yě duì xuéxí hěn hǎo.}\\
« Je pense que le sport est très important. Le sport est bon pour la santé, et aussi bon pour les études. »

\textbf{Équipe NUANCÉE (élève B) :} 你说得对，可是每天运动太累了。我觉得一个星期运动三次就可以了。\\
\emph{Nǐ shuō de duì, kěshì měitiān yùndòng tài lèi le. Wǒ juéde yí ge xīngqī yùndòng sān cì jiù kěyǐ le.}\\
« Tu as raison, mais faire du sport tous les jours, c'est trop fatigant. Je pense que trois fois par semaine, ça suffit. »

\textbf{Équipe POUR (élève C) :} 我不同意。每天运动不一定很累。我们可以每天走路、跑步，也可以踢足球。\\
\emph{Wǒ bù tóngyì. Měitiān yùndòng bù yídìng hěn lèi. Wǒmen kěyǐ měitiān zǒulù, pǎobù, yě kěyǐ tī zúqiú.}\\
« Je ne suis pas d'accord. Faire du sport tous les jours n'est pas forcément fatigant. On peut marcher, courir tous les jours, ou jouer au football. »

\textbf{Équipe NUANCÉE (élève D) :} 你说得对，可是我们还要学习，还要帮助爸爸妈妈。那我们可以周末多运动。\\
\emph{Nǐ shuō de duì, kěshì wǒmen hái yào xuéxí, hái yào bāngzhù bàba māma. Nà wǒmen kěyǐ zhōumò duō yùndòng.}\\
« Tu as raison, mais nous devons aussi étudier et aider papa et maman. Alors nous pouvons faire plus de sport le week-end. »

\textbf{Rapporteur :} 第一组说运动对身体和学习都很好；第二组说运动很重要，可是不用每天运动。所以，我们认为运动对每个人都很重要。谢谢大家！\\
\emph{Dì-yī zǔ shuō yùndòng duì shēntǐ hé xuéxí dōu hěn hǎo; dì-èr zǔ shuō yùndòng hěn zhòngyào, kěshì bú yòng měitiān yùndòng. Suǒyǐ, wǒmen rènwéi yùndòng duì měi ge rén dōu hěn zhòngyào. Xièxie dàjiā!}\\
« Le premier groupe dit que le sport est bon pour la santé et les études ; le deuxième groupe dit que le sport est important, mais qu'il n'est pas nécessaire d'en faire tous les jours. Donc, nous pensons que le sport est très important pour chacun. Merci à tous ! »

\textbf{Commentaire des choix de langue.}
\begin{itemize}
\item \textbf{Formules respectées :} chaque porte-parole utilise au moins trois formules fonctionnelles (\zh{我觉得}, \zh{我不同意}, \zh{你说得对，可是…}, \zh{那我们可以…}), conformément à la consigne.
\item \textbf{Structure d'opinion :} \zh{我觉得 + phrase complète} ; on ne met jamais \zh{是} devant l'adjectif : on dit \zh{运动非常重要} et non \zh{运动是非常重要}.
\item \textbf{Préposition \zh{对} :} \zh{运动对身体很好} (« le sport est bon \emph{pour} la santé ») ; structure Sujet + 对 + complément + adverbe + adjectif.
\item \textbf{\zh{不一定} :} « pas forcément », négation nuancée très utile dans un débat ; \zh{不用} signifie « pas besoin de ».
\item \textbf{Mesure de fréquence :} \zh{一个星期运动三次} : le complément de fréquence \zh{三次} se place après le verbe.
\item \textbf{Tons travaillés :} \zh{yùndòng} (4+4), \zh{zhòngyào} (4+4), \zh{shēntǐ} (1+3), \zh{lǐyóu} (3+2).
\end{itemize}
\end{exoresolu}

\courssub{Bilan de l'activité}

\begin{aretenir}[Ce que cette activité t'a appris]
\begin{itemize}
\item Tu sais maintenant \textbf{animer et participer à un débat} en chinois : présenter un sujet, donner ton avis, réagir poliment et conclure.
\item Les \cle{formules fonctionnelles} (\zh{我觉得}, \zh{我同意/不同意}, \zh{你说得对，可是…}, \zh{那我们可以…}, \zh{所以，我们认为…}) sont des outils réutilisables dans toute discussion orale.
\item La correction des \cle{tons} dans des phrases complètes est aussi importante que le vocabulaire : un ton faux peut changer le sens.
\item Un débat réussi repose sur le \cle{tour de parole} : écouter, réagir avec une formule, puis argumenter.
\end{itemize}
\end{aretenir}

\textbf{Grille d'observation de l'enseignant.}

\begin{center}
\begin{tabularx}{\linewidth}{|X|X|X|X|}
\hline
\rowcolor{popblueL} Critère & Bien (2) & Moyen (1) & À retravailler (0) \\
\hline
Prononciation des tons & Tons corrects dans les phrases & Quelques erreurs & Tons souvent faux \\
\hline
Formules fonctionnelles & 3 formules ou plus & 2 formules & 0 ou 1 formule \\
\hline
Fluidité & Parle sans lire & Lit un peu & Lit tout \\
\hline
Tour de parole & Écoute et réagit & Interrompt peu & Coupe la parole \\
\hline
\end{tabularx}
\end{center}

\begin{experience}[Pour aller plus loin : la radio du quartier]
En binôme, enregistre sur un téléphone une mini-émission de deux minutes : un élève est journaliste à la radio de Douala, l'autre est un jeune footballeur camerounais qui parle de son entraînement. Utilise au moins quatre formules fonctionnelles et trois mots du vocabulaire du sport. La classe écoute les enregistrements et choisit la meilleure interview.
\end{experience}

\newpage

\coursec{Méthodes et exercices résolus}

\courssub{Vocabulaire du sport et de la santé}

\begin{center}
\begin{tabularx}{\linewidth}{|X|X|X|X|}
\hline
\rowcolor{popblueL} \textbf{Caractères} & \textbf{Pinyin} & \textbf{Nature} & \textbf{Traduction} \\ \hline
运动 & yùndòng & nom/verbe & sport, exercice ; faire du sport \\ \hline
锻炼 & duànliàn & verbe & s'entraîner, faire de l'exercice physique \\ \hline
身体 & shēntǐ & nom & corps, condition physique \\ \hline
健康 & jiànkāng & adj/nom & en bonne santé ; santé \\ \hline
足球 & zúqiú & nom & football \\ \hline
篮球 & lánqiú & nom & basket-ball \\ \hline
乒乓球 & pīngpāngqiú & nom & tennis de table \\ \hline
羽毛球 & yǔmáoqiú & nom & badminton \\ \hline
游泳 & yóuyǒng & verbe/nom & nager ; natation \\ \hline
跑步 & pǎobù & verbe & courir, faire du jogging \\ \hline
爬山 & páshān & verbe & faire de la randonnée, grimper \\ \hline
比赛 & bǐsài & nom/verbe & match, compétition ; concourir \\ \hline
体育场 & tǐyùchǎng & nom & stade \\ \hline
健身房 & jiànshēnfáng & nom & salle de sport \\ \hline
频率 & pínlǜ & nom & fréquence \\ \hline
每周 & měi zhōu & loc. temps & chaque semaine \\ \hline
每天 & měitiān & loc. temps & tous les jours \\ \hline
三次 & sān cì & quantif. & trois fois (classif. \cle{次}) \\ \hline
增强 & zēngqiáng & verbe & renforcer \\ \hline
体质 & tǐzhì & nom & constitution physique \\ \hline
预防 & yùfáng & verbe & prévenir \\ \hline
疾病 & jíbìng & nom & maladie \\ \hline
缓解 & huǎnjiě & verbe & soulager, atténuer \\ \hline
压力 & yālì & nom & pression, stress \\ \hline
心情 & xīnqíng & nom & humeur, moral \\ \hline
意志力 & yìzhìlì & nom & force de volonté, volonté \\ \hline
适度 & shìdù & adj & modéré, mesuré \\ \hline
避免 & bìmiǎn & verbe & éviter \\ \hline
受伤 & shòushāng & verbe & se blesser \\ \hline
赞成 & zànchéng & verbe & être pour, approuver \\ \hline
反对 & fǎnduì & verbe & être contre, s'opposer \\ \hline
同意 & tóngyì & verbe & être d'accord \\ \hline
坚持 & jiānchí & verbe & persévérer, insister \\ \hline
当然 & dāngrán & adv. & bien sûr, naturellement \\ \hline
总之 & zǒngzhī & conj. & en bref, en conclusion \\ \hline
首先 & shǒuxiān & adv. & premièrement \\ \hline
其次 & qícì & adv. & deuxièmement \\ \hline
最后 & zuìhòu & adv. & enfin, finalement \\ \hline
虽然 & suīrán & conj. & bien que, quoique \\ \hline
但是 & dànshì & conj. & mais, cependant \\ \hline
因为 & yīnwèi & conj. & parce que \\ \hline
所以 & suǒyǐ & conj. & donc, c'est pourquoi \\ \hline
不但 & búdàn & adv. & non seulement \\ \hline
而且 & érqiě & adv. & mais encore, de plus \\ \hline
建议 & jiànyì & verbe/nom & suggérer ; suggestion \\ \hline
参加 & cānjiā & verbe & participer \\ \hline
集体 & jítǐ & nom/adj & collectif, groupe \\ \hline
交到 & jiāo dào & verbe & se faire (des amis), réussir à se faire \\ \hline
规律 & guīlǜ & adj/nom & régulier ; règle, loi \\ \hline
培养 & péiyǎng & verbe & cultiver, former \\ \hline
动起来 & dòng qǐlai & loc. verbale & se bouger, s'activer \\ \hline
\end{tabularx}
\end{center}

\courssub{Dialogue modèle 1 : Parler de ses sports}

\begin{textebox}[Dialogue : Au stade municipal (Yaoundé)]
阿马杜：周末你去干什么了？\quad (Āmǎdù: Zhōumò nǐ qù gàn shénme le?)
玲：我去市体育场踢足球了。\quad (Líng: Wǒ qù shì tǐyùchǎng tī zúqiú le.)
阿马杜：哦，你踢得好吗？\quad (Āmǎdù: Ó, nǐ tī de hǎo ma?)
玲：一般吧。我不但喜欢踢球，而且喜欢看比赛。\quad (Líng: Yībān ba. Wǒ búdàn xǐhuan tī qiú, érqiě xǐhuan kàn bǐsài.)
阿马杜：我也喜欢看足球赛。下次我们一起去吧？\quad (Āmǎdù: Wǒ yě xǐhuan kàn zúqiú sài. Xiàcì wǒmen yīqǐ qù ba?)
玲：好主意！\quad (Líng: Hǎo zhǔyi!)
\end{textebox}

\courssub{Dialogue modèle 2 : Mini-débat « Faut-il faire du sport tous les jours ? »}

\begin{textebox}[Débat : 每天运动好不好？]
阿马杜：我认为中学生应该每天运动。因为运动能增强体质，所以很有必要。\quad (Āmǎdù: Wǒ rènwéi zhōngxuéshēng yīnggāi měitiān yùndòng. Yīnwèi yùndòng néng zēngqiáng tǐzhì, suǒyǐ hěn yǒu bìyào.)
玲：我不同意。虽然运动好，但是每天运动身体会累，而且作业太多，没有时间。\quad (Líng: Wǒ bù tóngyì. Suīrán yùndòng hǎo, dànshì měitiān yùndòng shēntǐ huì lèi, érqiě zuòyè tài duō, méiyǒu shíjiān.)
阿马杜：你说得有道理。但是我觉得一周三四次最好。\quad (Āmǎdù: Nǐ shuō de yǒu dàolǐ. Dànshì wǒ juéde yī zhōu sān sì cì zuì hǎo.)
玲：我也赞成这个频率。适度运动，避免受伤。\quad (Líng: Wǒ yě zànchéng zhège pínlǜ. Shìdù yùndòng, bìmiǎn shòushāng.)
\end{textebox}

\courssub{Formules fonctionnelles : Présenter, opiner, relancer}

\begin{methode}[Présenter un exposé oral simple : « Mon sport préféré »]
\textbf{Objectif :} Produire un exposé court (1 à 2 minutes), structuré et cohérent, sur une activité sportive personnelle.
\textbf{Démarche en 5 étapes :}
\begin{enumerate}
    \item \textbf{Accroche et civilités :} Saluer l'auditoire, se présenter (nom, classe). \textit{Exemple :} \zh{大家好，我叫王明，我是高一A班的学生。} (Dàjiā hǎo, wǒ jiào Wáng Míng, wǒ shì Gāoyī A bān de xuésheng.) — « Bonjour à tous, je m'appelle Wang Ming, je suis élève de 1ère A. »
    \item \textbf{Annonce du sujet :} Utiliser \cle{今天我要谈谈…} (Jīntiān wǒ yào tántan…) ou \cle{我的主题是…} (Wǒ de zhǔtí shì…). \textit{Exemple :} \zh{今天我要谈谈我最喜欢的运动——篮球。} (Jīntiān wǒ yào tántan wǒ zuì xǐhuan de yùndòng — lánqiú.) — « Aujourd'hui je vais parler de mon sport préféré : le basket-ball. »
    \item \textbf{Développement structuré (3 arguments) :} Utiliser les connecteurs logiques \cle{首先} (shǒuxiān, premièrement), \cle{其次} (qícì, deuxièmement), \cle{最后} (zuìhòu, enfin). Chaque argument associe \textbf{fait personnel} + \textbf{bénéfice}. Structure clé : \zh{通过…可以…} (Tōngguò… kěyǐ…, « Grâce à… on peut… »).
    \item \textbf{Conclusion et ouverture :} Résumer en une phrase avec \cle{总之} (zǒngzhī, « en bref ») et inviter à la pratique : \zh{我建议大家也试一试。} (Wǒ jiànyì dàjiā yě shì yī shì. — « Je suggère à tous d'essayer aussi. »)
    \item \textbf{Remerciements et clôture :} \zh{谢谢大家，我的发言完了。} (Xièxie dàjiā, wǒ de fāyán wánle. — « Merci à tous, mon intervention est terminée. »)
\end{enumerate}
\textbf{Structures à retenir :} \zh{我认为…因为…} (Wǒ rènwéi… yīnwèi… — « Je pense que… parce que… ») ; \zh{不但…而且…} (Búdàn… érqiě… — « Non seulement… mais encore… »).
\end{methode}

\begin{methode}[Argumenter dans un mini-débat : « Faut-il faire du sport tous les jours ? »]
\textbf{Objectif :} Participer à un échange contradictoire en exprimant un avis nuancé (accord, désaccord, concession) avec un vocabulaire fonctionnel précis.
\textbf{Démarche en 4 étapes :}
\begin{enumerate}
    \item \textbf{Prendre position clairement :} \zh{我赞成…} (Wǒ zànchéng… — « Je suis pour… ») / \zh{我反对…} (Wǒ fǎnduì… — « Je suis contre… ») / \zh{我觉得…} (Wǒ juéde… — « Je trouve que… »).
    \item \textbf{Justifier avec \cle{因为…所以…} (Yīnwèi… suǒyǐ…) :} Lier cause et conséquence. \textit{Exemple :} \zh{因为每天运动对心脏好，所以我赞成。} (Yīnwèi měitiān yùndòng duì xīnzàng hǎo, suǒyǐ wǒ zànchéng.)
    \item \textbf{Nuancer avec la concession \cle{虽然…，但是/可是…} (Suīrán…, dànshì/kěshì…) :} Indispensable pour montrer un niveau B1/B2. \textit{Exemple :} \zh{虽然每天运动好，但是身体也需要休息。} (Suīrán měitiān yùndòng hǎo, dànshì shēntǐ yě xūyào xiūxi.)
    \item \textbf{Relancer l'interlocuteur :} Utiliser les formules de \cle{relance} : \zh{你呢？} (Nǐ ne? — « Et toi ? »), \zh{你怎么看？} (Nǐ zěnme kàn? — « Qu'en penses-tu ? »), \zh{为什么？} (Wèishénme? — « Pourquoi ? »).
\end{enumerate}
\textbf{Tableau récapitulatif des formules fonctionnelles :}
\begin{center}
\begin{tabularx}{\linewidth}{|X|X|X|}
\hline
\rowcolor{popblueL} \textbf{Fonction} & \textbf{Chinois (Caractères + Pinyin)} & \textbf{Français} \\ \hline
Donner son avis & \zh{我认为…} (Wǒ rènwéi…) / \zh{我觉得…} (Wǒ juéde…) & Je pense que… / Je trouve que… \\ \hline
Être d'accord & \zh{我赞成 / 我同意} (Wǒ zànchéng / Wǒ tóngyì) & Je suis pour / J'approuve \\ \hline
Ne pas être d'accord & \zh{我反对 / 我不同意} (Wǒ fǎnduì / Wǒ bù tóngyì) & Je suis contre / Je n'approuve pas \\ \hline
Argumenter (Cause) & \zh{因为…} (Yīnwèi…) & Parce que… \\ \hline
Argumenter (Conséquence) & \zh{所以…} (Suǒyǐ…) & Donc… / C'est pourquoi… \\ \hline
Concéder / Nuancer & \zh{虽然…，但是…} (Suīrán…, dànshì…) & Bien que…, cependant… \\ \hline
Relancer & \zh{你怎么看？ / 你呢？} (Nǐ zěnme kàn? / Nǐ ne?) & Qu'en penses-tu ? / Et toi ? \\ \hline
\end{tabularx}
\end{center}
\end{methode}

\courssub{Travail sur les tons et la prononciation}

\begin{methode}[Maîtriser les tons du vocabulaire sportif et les particules à ton neutre]
\textbf{Objectif :} Corriger les erreurs tonales systématiques des francophones (sandhi du 3e ton, ton neutre obligatoire sur particules).
\textbf{Démarche en 3 étapes :}
\begin{enumerate}
    \item \textbf{Identifier le \cle{Sandhi du 3e ton} (deux 3e tons consécutifs) :} Règle absolue : 3e + 3e $\rightarrow$ 2e + 3e.
    \begin{itemize}
        \item \zh{比赛} (bǐsài) : \zh{比} (bǐ, 3e) + \zh{赛} (sài, 4e) $\rightarrow$ \textbf{Pas de sandhi} (3e + 4e). Prononciation standard : \zh{bǐsài}. En flux rapide, \zh{比} peut devenir demi-3e ton (bas), mais \textbf{jamais 2e ton} (bí).
        \item \zh{跑步} (pǎobù) : \zh{跑} (pǎo, 3e) + \zh{步} (bù, 4e) $\rightarrow$ \textbf{Pas de sandhi}. Standard : \zh{pǎobù}.
        \item \zh{运动} (yùndòng) : \zh{运} (yùn, 3e) + \zh{动} (dòng, 4e) $\rightarrow$ \textbf{Pas de sandhi}. Standard : \zh{yùndòng}. \textbf{Attention :} Ne pas prononcer *yún dòng (2e + 4e) ni *yǔn dòng (3e ton bas maintenu « haché »). En phrase, \zh{运} réalise souvent un \emph{demi-3e ton} (courbe descendante incomplète).
        \item \zh{游泳} (yóuyǒng) : 2e + 3e $\rightarrow$ Pas de sandhi.
        \item \textbf{Vrai sandhi 3e+3e :} \zh{很好} (hěn hǎo $\rightarrow$ hén hǎo), \zh{五百} (wǔ bǎi $\rightarrow$ wú bǎi).
    \end{itemize}
    \item \textbf{Maîtriser le \cle{ton neutre} (qīngshēng) sur les particules grammaticales :} Courtes, faibles, sans contour tonal. Erreur fatale : leur donner un ton plein (1, 2, 3 ou 4).
    \begin{itemize}
        \item \zh{了} (le) aspect accompli : \zh{我踢了球。} (Wǒ tī \textbf{le} qiú.) \textit{Pas *lè (4e = joie).}
        \item \zh{的} (de) possession/modification : \zh{我的球鞋。} (Wǒ \textbf{de} qiúxié.) \textit{Pas *dì (4e) ni *de (1e).}
        \item \zh{吧} (ba) suggestion : \zh{我们去跑步吧。} (Wǒmen qù pǎobù \textbf{ba}.) \textit{Pas *bā (1e = bar) ni *bà (4e = arrêter).}
        \item \zh{吗} (ma) question fermée : \zh{你喜欢足球吗？} (Nǐ xǐhuan zúqiú \textbf{ma}?) \textit{Pas *mà (4e = gronder) ni *mā (1e = mère).}
        \item \zh{呢} (ne) relance : \zh{我喜欢篮球，你呢？} (Wǒ xǐhuan lánqiú, nǐ \textbf{ne}?)
    \end{itemize}
    \item \textbf{Intonation des questions :} \cle{吗} $\rightarrow$ montante finale légère. \cle{呢} $\rightarrow$ soutenue, plate (relance). \cle{怎么样？} $\rightarrow$ montante sur \cle{么}.
\end{enumerate}
\textbf{Exercice de discrimination auditive (binôme) :} L'un lit la colonne A (correct), l'autre identifie l'erreur en B.
\begin{center}
\begin{tabularx}{\linewidth}{|X|X|}
\hline
\rowcolor{poporangeL} \textbf{A (Correct standard)} & \textbf{B (Erreur fréquente francophone)} \\ \hline
\zh{yùndòng} (运动) & *yǔndòng (3e ton bas maintenu) / *yúndòng (2e ton forcé) \\ \hline
\zh{pǎobù} (跑步) & *páobù (2e ton forcé) / *pǎo bù (3e ton bas) \\ \hline
\zh{Wǒ tī \textbf{le} qiú.} (我踢了球。) & *Wǒ tī \textbf{lè} qiú. (le $\rightarrow$ 4e ton) \\ \hline
\zh{Nǐ xǐhuan zúqiú \textbf{ma}?} (你喜欢足球吗？) & *Nǐ xǐhuan zúqiú \textbf{mà}? (ma $\rightarrow$ 4e ton) \\ \hline
\zh{Wǒmen qù yóuyǒng \textbf{ba}.} (我们去游泳吧。) & *Wǒmen qù yóuyǒng \textbf{bā}. (ba $\rightarrow$ 1er ton) \\ \hline
\end{tabularx}
\end{center}
\end{methode}

\courssub{Méthode : Relancer l'échange — Questions de suivi et réactifs}

\begin{methode}[Éviter les silences blancs : interaction orale]
\textbf{Objectif :} Montrer l'écoute active (compétence interactionnelle évaluée au Bac) par des relances et des réactifs.
\textbf{Démarche en 3 étapes :}
\begin{enumerate}
    \item \textbf{Relance symétrique \cle{你呢？} (Nǐ ne?) — « Et toi ? » :} Après une réponse courte pour retourner la parole. \textit{Exemple :} A: \zh{我喜欢篮球。} B: \zh{我喜欢足球，\textbf{你呢？} }
    \item \textbf{Relance ouverte \cle{为什么？} (Wèishénme?) / \cle{怎么？} (Zěnme?) :} Demander justification ou détail. \textit{Exemple :} A: \zh{我不喜欢跑步。} B: \zh{\textbf{为什么？} } (Ouvert) / \zh{\textbf{怎么？} } (Surpris).
    \item \textbf{Relance évaluative \cle{怎么样？} (Zěnmeyàng?) — « Comment c'est ? / Qu'en penses-tu ? » :} Pour demander un avis sur une proposition. \textit{Exemple :} A: \zh{下次一起去游泳吧。} B: \zh{\textbf{怎么样？} } (Souvent suivi de \zh{好主意！}).
    \item \textbf{Les \cle{réactifs} (marqueurs d'écoute active) :} Insérer \textbf{sans prendre la parole} :
    \begin{itemize}
        \item \zh{真的吗？} (Zhēn de ma? — « Vraiment ? » surprise/vérification)
        \item \zh{是吗？} (Shì ma? — « Ah bon ? / C'est vrai ? » intérêt poli)
        \item \zh{原来如此。} (Yuánlái rúcǐ. — « Je comprends / C'est donc ça. » compréhension)
        \item \zh{我也一样。} (Wǒ yě yīyàng. — « Moi aussi. » convergence)
    \end{itemize}
\end{enumerate}
\textbf{Attention :} Ne pas confondre \zh{你呢？} (relance symétrique « et toi ? ») et \zh{怎么样？} (demande d'évaluation « comment trouves-tu ? »).
\end{methode}

\courssub{Exercices résolus et production guidée}

\begin{exoresolu}[Rédiger le script d'un exposé oral sur le football]
\textbf{Énoncé :} Vous êtes \cle{Ngoa} (枚奥 Méi'ào), élève à Yaoundé. Préparez le texte intégral (caractères + pinyin) d'un exposé de 1 minute sur le football. Contraintes : 1) fréquence (tous les samedis), 2) lieu (stade municipal), 3) deux bienfaits (santé, amis), 4) invitation finale.
\tcblower
\textbf{Solution détaillée :}
\textbf{1. Analyse :} Registre oral standard (HSK 2-3). Vocabulaire : \zh{踢足球} (tī zúqiú), \zh{每个星期六} (měi gè xīngqīliù), \zh{市体育场} (shì tǐyùchǎng), \zh{锻炼身体} (duànliàn shēntǐ), \zh{结交朋友} (jiējiāo péngyou).
\textbf{2. Construction (Méthode 5 étapes) :}
\begin{itemize}
    \item \textit{Accroche :} \zh{大家好，我叫枚奥。} (Dàjiā hǎo, wǒ jiào Méi'ào.)
    \item \textit{Sujet :} \zh{今天我要谈谈足球。} (Jīntiān wǒ yào tántan zúqiú.)
    \item \textit{Arg 1 (Fréquence/Lieu) :} \zh{首先，我每个星期六在市体育场踢足球。} (Shǒuxiān, wǒ měi gè xīngqīliù zài shì tǐyùchǎng tī zúqiú.) \textbf{Règle :} \cle{Sujet + Temps + Lieu (在…) + Verbe}.
    \item \textit{Arg 2 (Bienfaits) :} \zh{其次，踢足球不但能锻炼身体，而且能结交很多朋友。} (Qícì, tī zúqiú búdàn néng duànliàn shēntǐ, érqiě néng jiējiāo hěnduō péngyou.) Structure \cle{不但…而且…} + \cle{能}.
    \item \textit{Conclusion :} \zh{总之，足球很有意思。} (Zǒngzhī, zúqiú hěn yǒu yìsi.)
    \item \textit{Invitation :} \zh{我建议大家也来踢球吧！} (Wǒ jiànyì dàjiā yě lái tī qiú ba!) Particule \cle{ba} suggestion douce.
    \item \textit{Clôture :} \zh{谢谢大家。} (Xièxie dàjiā.)
\end{itemize}
\textbf{3. Texte final (prêt à l'oral) :}
\begin{textebox}[Exposé : Mon sport préféré — Football]
大家好，我叫枚奥。今天我要谈谈足球。首先，我每个星期六在市体育场踢足球。其次，踢足球不但能锻炼身体，而且能结交很多朋友。总之，足球很有意思。我建议大家也来踢球吧！谢谢大家。
\end{textebox}
\textbf{Pinyin :} Dàjiā hǎo, wǒ jiào Méi'ào. Jīntiān wǒ yào tántan zúqiú. Shǒuxiān, wǒ měi gè xīngqīliù zài shì tǐyùchǎng tī zúqiú. Qícì, tī zúqiú búdàn néng duànliàn shēntǐ, érqiě néng jiējiāo hěnduō péngyou. Zǒngzhī, zúqiú hěn yǒu yìsi. Wǒ jiànyì dàjiā yě lái tī qiú ba! Xièxie dàjiā.
\textbf{Traduction :} « Bonjour à tous, je m'appelle Ngoa. Aujourd'hui je parle du football. Premièrement, chaque samedi je joue au foot au stade municipal. Deuxièmement, le foot non seulement renforce le corps, mais encore permet de se faire beaucoup d'amis. En bref, le foot est très intéressant. Je vous suggère de venir jouer aussi ! Merci à tous. »
\end{exoresolu}

\begin{exoresolu}[Transformer des opinions brutes en arguments de débat structurés]
\textbf{Énoncé :} Réécrivez en chinois (caractères + pinyin) les opinions ci-dessous, prêtes pour un débat oral, en utilisant : 1) formule d'opinion, 2) \cle{因为…所以…}, 3) \cle{虽然…但是…} (pour au moins deux).
\begin{enumerate}
    \item « Je suis pour le sport quotidien car c'est bon pour la santé, mais attention aux blessures. »
    \item « Je suis contre : on n'a pas le temps avec les devoirs. »
    \item « Je trouve que 3 fois par semaine c'est l'idéal, car le corps a besoin de repos. »
\end{enumerate}
\tcblower
\textbf{Solution détaillée :}
\textbf{Opinion 1 (Accord nuancé) :}
\begin{itemize}
    \item \textit{Vocabulaire :} \zh{受伤} (shòushāng), \zh{容易} (róngyì, facile à / susceptible de).
    \item \textit{Structure :} \zh{我赞成每天运动，因为对健康有好处，虽然容易受伤，但是好处更多。}
    \item \textit{Pinyin :} Wǒ zànchéng měitiān yùndòng, yīnwèi duì jiànkāng yǒu hǎochù, suīrán róngyì shòushāng, dànshì hǎochù gèng duō.
    \item \textit{Analyse :} \cle{对…有好处}. Concession : \cle{虽然容易受伤} (risque réel) $\rightarrow$ \cle{但是好处更多} (contre-argument).
\end{itemize}
\textbf{Opinion 2 (Désaccord ferme) :}
\begin{itemize}
    \item \textit{Vocabulaire :} \zh{作业} (zuòyè), \zh{时间} (shíjiān).
    \item \textit{Structure :} \zh{我反对每天运动，因为作业太多，没有时间。}
    \item \textit{Pinyin :} Wǒ fǎnduì měitiān yùndòng, yīnwèi zuòyè tài duō, méiyǒu shíjiān.
    \item \textit{Analyse :} Position ferme, pas de concession. \cle{没有时间} = \cle{没} + \cle{有} + objet.
\end{itemize}
\textbf{Opinion 3 (Position intermédiaire) :}
\begin{itemize}
    \item \textit{Vocabulaire :} \zh{一周三次} (yī zhōu sān cì), \zh{最好} (zuì hǎo), \zh{休息} (xiūxi).
    \item \textit{Structure :} \zh{我认为一周三次最好，因为身体需要休息，虽然每天运动也可以，但是容易累。}
    \item \textit{Pinyin :} Wǒ rènwéi yī zhōu sān cì zuì hǎo, yīnwèi shēntǐ xūyào xiūxi, suīrán měitiān yùndòng yě kěyǐ, dànshì róngyì lèi.
    \item \textit{Analyse :} \cle{我认为} (avis). Concession : \cle{虽然每天运动也可以} (concède le possible) $\rightarrow$ \cle{但是容易累} (inconvénient majeur). \cle{容易} + verbe état négatif.
\end{itemize}
\textbf{Conclusion :} La maîtrise de \cle{因为…所以…} et \cle{虽然…但是…} distingue le niveau A2 (phrases simples) du B1 (discours complexe) au Bac oral.
\end{exoresolu}

\begin{exoresolu}[Dictée tonale et correction de phrases orales]
\textbf{Énoncé :} Identifiez chaque erreur de ton, expliquez la règle (sandhi, ton neutre, confusion) et donnez le pinyin correct.
\begin{enumerate}
    \item *Wǒ \textbf{yǔn}dòng hěn hǎo. (Cible : \zh{运动很好。})
    \item *Tā \textbf{pǎo} bù \textbf{lè}. (Cible : \zh{他跑步了。})
    \item *Nǐ \textbf{xǐhuan} \textbf{mà}? (Cible : \zh{你喜欢吗？})
    \item *Wǒmen \textbf{qù} \textbf{yóu} yǒng \textbf{bā}. (Cible : \zh{我们去游泳吧。})
\end{enumerate}
\tcblower
\textbf{Correction détaillée :}
\begin{enumerate}
    \item \textbf{Erreur :} *yǔn (3e ton bas maintenu) ou *yún (2e ton forcé). \textbf{Correct :} \zh{yùn} (3e ton standard, ou demi-3e ton en flux). \textbf{Règle :} \zh{运} (3e) + \zh{动} (4e) $\rightarrow$ \textbf{pas de sandhi obligatoire} vers le 2e ton. Le 3e ton garde sa nature (descendant-ascendant en citation, demi-3e en phrase). L'erreur *yǔn rend la parole hachée ; *yún est une hypercorrection (confusion avec règle 3e+3e).
    \item \textbf{Erreur :} *lè (4e ton). \textbf{Correct :} \textbf{le} (ton neutre). \textbf{Règle :} Particule aspectuelle \cle{了} = ton neutre obligatoire. *lè = \zh{乐} (joie). Confusion graphique/phonétique.
    \item \textbf{Erreur :} *mà (4e ton). \textbf{Correct :} \textbf{ma} (ton neutre). \textbf{Règle :} Particule interrogative \cle{吗} = ton neutre obligatoire. *mà = \zh{骂} (gronder/insulter). Risque de sens inverse !
    \item \textbf{Erreur :} *bā (1er ton). \textbf{Correct :} \textbf{ba} (ton neutre). \zh{游泳} (yóuyǒng) : 2e + 3e, correct. \textbf{Règle :} Particule modale \cle{吧} = ton neutre. *bā = \zh{巴} (bar/attendre), *bà = \zh{霸} (hégémon/arrêter).
    \textbf{Phrase corrigée :} Wǒmen qù yóuyǒng ba.
\end{enumerate}
\textbf{Bilan :} 3 erreurs sur 4 concernent le \cle{ton neutre} (\cle{了, 吗, 吧}). Cause n°1 de perte de points à l'oral (intelligibilité + correction grammaticale).
\end{exoresolu}

\begin{exoresolu}[Compléter un dialogue lacunaire avec formules de relance et réactifs]
\textbf{Énoncé :} Complétez le dialogue entre \cle{Amadou} (阿马杜 Āmǎdù) et \cle{Ling} (玲 Líng) avec la boîte à outils (chaque item une fois).
\textbf{Boîte à outils :} \zh{你呢？} \quad \zh{为什么？} \quad \zh{真的吗？} \quad \zh{我也一样。} \quad \zh{原来如此。} \quad \zh{怎么样？}
\begin{textebox}[Dialogue à trous]
阿马杜：我周末去爬山了，很累但很开心。
玲：\underline{\hspace{3cm}} 我也想去。
阿马杜：下次我们一起去吧。
玲：\underline{\hspace{3cm}} 好主意！你喜欢爬什么山？
阿马杜：我喜欢喀麦隆的曼恩古巴山。
玲：\underline{\hspace{3cm}} 我没去过。
阿马杜：风景很美，\underline{\hspace{3cm}}
玲：下次带我去吧！
\end{textebox}
\tcblower
\textbf{Solution détaillée :}
\begin{enumerate}
    \item \zh{玲：\textbf{我也一样。} 我也想去。} (Réactif convergence + envie partagée). Pinyin: Wǒ yě yīyàng. Wǒ yě xiǎng qù.
    \item \zh{玲：\textbf{怎么样？} 好主意！} (Relance évaluative sur proposition $\rightarrow$ validation). Pinyin: Zěnmeyàng? Hǎo zhǔyi!
    \item \zh{玲：\textbf{真的吗？} 我没去过。} (Réactif surprise face info nouvelle $\rightarrow$ aveu). Pinyin: Zhēn de ma? Wǒ méi qùguo.
    \item \zh{阿马杜：风景很美，\textbf{你呢？}} (Relance symétrique après description : invite Ling à réagir). Pinyin: Fēngjǐng hěn měi, nǐ ne?
\end{enumerate}
\textbf{Dialogue complet corrigé :}
\begin{textebox}[Dialogue corrigé]
阿马杜：我周末去爬山了，很累但很开心。
玲：\textbf{我也一样。} 我也想去。
阿马杜：下次我们一起去吧。
玲：\textbf{怎么样？} 好主意！你喜欢爬什么山？
阿马杜：我喜欢喀麦隆的曼恩古巴山。
玲：\textbf{真的吗？} 我没去过。
阿马杜：风景很美，\textbf{你呢？}
玲：下次带我去吧！
\end{textebox}
\textbf{Justification :} \zh{去过} (qùguo) = \cle{过} expérience. \zh{带我去} (dài wǒ qù) = construction \cle{把} implicite (m'amener). Enchaînement naturel \zh{真的吗？} $\rightarrow$ \zh{我没去过}.
\end{exoresolu}

\courssub{Méthode : Passer de l'oral à l'écrit — Paragraphe argumentatif}

\begin{methode}[Structurer un paragraphe écrit « Bienfaits du sport » (80-100 caractères)]
\textbf{Objectif :} Produire un texte écrit cohérent, trace de l'oral, avec connecteurs et paragraphage.
\textbf{Démarche en 4 étapes :}
\begin{enumerate}
    \item \textbf{Thèse claire (Phrase sujet) :} \zh{坚持体育锻炼对身心健康非常重要。} (Jiānchí tǐyù duànliàn duì shēnxīn jiànkāng fēicháng zhòngyào.) Vocabulaire : \cle{坚持} (persévérer), \cle{身心健康} (santé corps-esprit).
    \item \textbf{Arguments ordonnés (Connecteurs logiques) :}
    \begin{itemize}
        \item \cle{首先} / \cle{第一} : Physiologique. \zh{首先，运动能增强体质，预防疾病。} (Yùndòng néng zēngqiáng tǐzhì, yùfáng jíbìng.) — \cle{增强}, \cle{体质}, \cle{预防}.
        \item \cle{其次} / \cle{第二} : Psycho-social. \zh{其次，参加集体运动可以缓解压力，交到朋友。} (Cānjiā jítǐ yùndòng kěyǐ huǎnjiě yālì, jiāo dào péngyou.) — \cle{缓解}, \cle{压力}, \cle{交到}.
        \item \cle{最后} : Mental. \zh{最后，规律锻炼能培养意志力。} (Guīlǜ duànliàn néng péiyǎng yìzhìlì.) — \cle{规律}, \cle{培养}, \cle{意志力}.
    \end{itemize}
    \item \textbf{Précaution (Niveau B1+) :} \zh{当然，运动量要适度，避免受伤。} (Dāngrán, yùndòngliàng yào shìdù, bìmiǎn shòushāng.) — \cle{适度}, \cle{避免}.
    \item \textbf{Conclusion synthétique :} \zh{总之，让我们动起来吧！} (Zǒngzhī, ràng wǒmen dòng qǐlai ba!) — \cle{动起来}, \cle{吧}.
\end{enumerate}
\textbf{Schéma de structure (TikZ) :}
\begin{popfigure}
\begin{tikzpicture}[node distance=1.0cm, every node/.style={align=center, font=\small, inner sep=3pt}]
\node[draw, rounded corners, fill=popblueL, text width=3.8cm] (thesis) {\textbf{Thèse}\\\zh{坚持锻炼很重要}};
\node[draw, rounded corners, fill=popgreenL, text width=3.8cm, below of=thesis] (arg1) {\textbf{Arg 1 (Physio)}\\\zh{首先：增强体质、预防病}};
\node[draw, rounded corners, fill=popgreenL, text width=3.8cm, below of=arg1] (arg2) {\textbf{Arg 2 (Psycho/Social)}\\\zh{其次：缓解压力、交友}};
\node[draw, rounded corners, fill=popgreenL, text width=3.8cm, below of=arg2] (arg3) {\textbf{Arg 3 (Mental)}\\\zh{最后：培养意志}};
\node[draw, rounded corners, fill=poporangeL, text width=3.8cm, below of=arg3] (counter) {\textbf{Précaution}\\\zh{当然：适度、避免伤}};
\node[draw, rounded corners, fill=poppinkL, text width=3.8cm, below of=counter] (concl) {\textbf{Conclusion}\\\zh{总之：动起来！}};
\draw[-{Stealth[length=3mm]}, popblue, thick] (thesis) -- (arg1);
\draw[-{Stealth[length=3mm]}, popblue, thick] (arg1) -- (arg2);
\draw[-{Stealth[length=3mm]}, popblue, thick] (arg2) -- (arg3);
\draw[-{Stealth[length=3mm]}, popblue, thick] (arg3) -- (counter);
\draw[-{Stealth[length=3mm]}, popblue, thick] (counter) -- (concl);
\end{tikzpicture}
\legende{Structure type d'un paragraphe argumentatif écrit (HSK 3-4)}
\end{popfigure}
\end{methode}

\begin{exoresolu}[Rédaction écrite : « Les bienfaits du sport pour les lycéens »]
\textbf{Énoncé :} Rédigez en chinois (caractères + pinyin) un paragraphe (~90 caractères) sur : \cle{中学生应该坚持体育锻炼}. Utilisez \textbf{au minimum} : \cle{首先}, \cle{其次}, \cle{最后}, \cle{当然}, \cle{总之}, \cle{增强体质}, \cle{缓解压力}, \cle{适度}.
\tcblower
\textbf{Solution détaillée :}
\textbf{1. Plan (application méthode) :}
\begin{itemize}
    \item Thèse : \zh{中学生应该坚持体育锻炼。}
    \item Arg 1 : \zh{首先，能增强体质，预防疾病。}
    \item Arg 2 : \zh{其次，能缓解学习压力，心情变好。}
    \item Arg 3 : \zh{最后，能培养坚强的意志力。}
    \item Précaution : \zh{当然，运动量要适度，避免受伤。}
    \item Conclusion : \zh{总之，健康是革命的本钱，动起来吧！} (Réf. culturelle : Mao Zedong).
\end{itemize}
\textbf{2. Texte final (Caractères) :}
\begin{textebox}[Rédaction : Bienfaits du sport pour lycéens]
中学生应该坚持体育锻炼。首先，能增强体质，预防疾病。其次，能缓解学习压力，心情变好。最后，能培养坚强的意志力。当然，运动量要适度，避免受伤。总之，健康是革命的本钱，动起来吧！
\end{textebox}
\textbf{3. Pinyin (tons corrects, sandhi noté si oral) :}
Zhōngxuéshēng yīnggāi jiānchí tǐyù duànliàn. Shǒuxiān, néng zēngqiáng tǐzhì, yùfáng jíbìng. Qícì, néng huǎnjiě xuéxí yālì, xīnqíng biàn hǎo. Zuìhòu, néng péiyǎng jiānqiáng de yìzhìlì. Dāngrán, yùndòngliàng yào shìdù, bìmiǎn shòushāng. Zǒngzhī, jiànkāng shì gémìng de běnqián, dòng qǐlai ba!
\textbf{4. Traduction :} « Les lycéens devraient persévérer dans l'exercice physique. Premièrement, cela peut renforcer la constitution et prévenir les maladies. Deuxièmement, cela peut soulager la pression des études et améliorer l'humeur. Troisièmement, cela peut cultiver une volonté forte. Bien sûr, la dose d'exercice doit être modérée pour éviter les blessures. En bref, la santé est le capital de la révolution, bougeons-nous ! »
\textbf{5. Points de grammaire vérifiés :}
\begin{itemize}
    \item \zh{应该} (yīnggāi) + V : obligation/conseil.
    \item \zh{能} (néng) répété : parallélisme rhétorique.
    \item \zh{心情变好} (xīnqíng biàn hǎo) : Sujet + \cle{变} + Adj.
    \item \zh{要适度} (yào shìdù) : \cle{要} (devoir) + Adj prédicatif.
    \item \zh{避免受伤} (bìmiǎn shòushāng) : \cle{避免} + V.
\end{itemize}
\textbf{Conclusion :} Modèle de « version » type Bac : vocabulaire programme, connecteurs obligatoires, structures ciblées (\cle{能, 要, 避免, 变}), référence culturelle.
\end{exoresolu}

\courssub{Activité de classe : Jeux de rôle et débat guidé}

\begin{experience}[Débat guidé : « Le sport à l'école : tous les jours ou 3 fois par semaine ? »]
\textbf{Consigne :} En groupes de 4 (2 « Pour le quotidien », 2 « Pour 3 fois/semaine »).
\textbf{Étapes :}
\begin{enumerate}
    \item \textbf{Préparation (5 min) :} Chaque camp note 3 arguments avec \cle{因为…所以…} et 1 concession \cle{虽然…但是…} sur fiche.
    \item \textbf{Débat (10 min) :} Tour de table. Un élève lance : \zh{我认为…因为…} ; adversaire répond : \zh{我不同意，虽然…但是…} ; observateurs notent formules utilisées (\cle{你呢？}, \cle{真的吗？}, \cle{原来如此。}).
    \item \textbf{Synthèse orale (3 min) :} Un rapporteur par groupe résume : \zh{我们组认为…总之…}.
    \item \textbf{Retour prof :} Correction tons (neutres, sandhi), ordre mots (Lieu avant V), classificateurs (\cle{次}).
\end{enumerate}
\textbf{Critères d'évaluation (grille Bac) :} Prononciation (tons), Cohérence (connecteurs), Interaction (relances/réactifs), Vocabulaire (précision), Grammaire (structures ciblées).
\end{experience}

\courssub{Éléments socioculturels : Cameroun — Chine}

\begin{savaistu}[Sport et société : regards croisés]
\textbf{Chine :} La \cle{radio gymnastique} (广播体操 \zh{guǎngbō tǐcāo}) est une institution depuis les années 1950. Chaque matin, dans les écoles, usines et parcs, des millions de Chinois pratiquent ces exercices codifiés sur fond musical. C'est un héritage de l'hygiénisme d'État et du maoïsme (\zh{发展体育运动，增强人民体质} — « Développer le sport, renforcer la constitution du peuple », 1952). Aujourd'hui, le \cle{sport scolaire} est noté au \cle{Gaokao} (épreuve de 50-100 pts selon provinces : course, saut, lancer). Sports stars : Yao Ming (basket), Liu Xiang (haies), Sun Yang (natation).
\textbf{Cameroun :} Le \cle{football} est roi (Lions Indomptables, CAN 1984, 1988, 2000, 2002, 2017). Samuel Eto'o, Roger Milla icônes nationales. L'\cle{athlétisme} (triple saut : Françoise Mbango, double championne olympique) et la \cle{boxe} (Hassan N'Dam, champion du monde) brillent. L'\cle{EPS} est obligatoire au lycée (2h/semaine) mais peu évaluée au Bac. Infrastructures : Stade Ahmadou Ahidjo (Yaoundé), Stade de la Réunification (Douala).
\textbf{Comparaison :} Chine = sport de masse encadré par l'État, routine matinale, notation académique forte. Cameroun = sport passion populaire (football), talents individuels émergents, infrastructure inégale. Point commun : le sport comme fierté nationale et vecteur de santé publique.
\textbf{Vocabulaire culturel :} \zh{广播体操} (guǎngbō tǐcāo), \zh{高考体育} (gāokǎo tǐyù), \zh{狮子不可征服} (Shīzi bùkě zhēngfú — Lions Indomptables), \zh{体育强国} (tǐyù qiángguó — nation forte par le sport).
\end{savaistu}

\courssub{Écriture des caractères clés : Règles et pièges}

\begin{aretenir}[Caractères de la leçon : ordre des traits et radicaux]
\begin{itemize}
    \item \zh{运} (yùn) : Radicale \cle{辶} (chuò, marche). Ordre : haut→bas (云), puis \cle{辶} (point, horizontale, pliure, point). \textit{Sens :} mouvement, transporter.
    \item \zh{动} (dòng) : Radicale \cle{力} (lì, force). Haut : \cle{云} (nuage) simplifié. Ordre : haut→bas, \cle{力} en dernier. \textit{Sens :} bouger, action.
    \item \zh{锻} (duàn) : Radicale \cle{钅} (jīn, métal/or). Droite : \cle{段} (duàn, segment). Ordre : \cle{钅} (point, horizontale, verticale, point), puis \cle{段}. \textit{Sens :} forger, tremper (métal) $\rightarrow$ s'entraîner (corps).
    \item \zh{炼} (liàn) : Radicale \cle{火} (huǒ, feu). Droite : \cle{东} (dōng, est). Ordre : \cle{火} (point, trait, trait, point), puis \cle{东}. \textit{Sens :} raffiner (métal), affiner $\rightarrow$ s'entraîner (esprit/corps).
    \item \zh{身} (shēn) : Corps. 7 traits. Ordre : horizontale, verticale, horizontale, pliure, horizontale, verticale, horizontale. \textbf{Piège :} Ne pas confondre avec \zh{申} (shēn, étendre/demander, 5 traits, verticale initiale longue).
    \item \zh{体} (tǐ) : Radicale \cle{亻} (rén, homme). Droite : \cle{本} (běn, racine). Ordre : \cle{亻} (no, pié), puis \cle{本}.
\end{itemize}
\textbf{Distinction cruciale :} \zh{锻炼} (duànliàn, exercice physique standard) $\neq$ *\zh{锻练} (rare, forger/entraîner). \zh{身体} (shēntǐ, corps) $\neq$ *\zh{申体}.
\end{aretenir}

\begin{attention}[Les 5 erreurs qui coûtent des points au Bac (Leçon 11)]
\begin{enumerate}
    \item \textbf{Tons : Particules à ton neutre prononcées en ton plein.}
    \textit{Exemple :} *Wǒ tī \textbf{lè} qiú (\zh{了} $\rightarrow$ 4e ton « joie ») au lieu de \zh{le} (neutre) ; *Nǐ qù \textbf{mà} (\zh{吗} $\rightarrow$ 4e ton « gronder ») au lieu de \zh{ma} (neutre) ; *Wǒmen qù \textbf{bā} (\zh{吧} $\rightarrow$ 1er ton « bar ») au lieu de \zh{ba} (neutre).
    \textbf{Conséquence :} Changement de sens radical + perte de points « Prononciation » et « Grammaire ».
    \item \textbf{Ordre des mots : Complément de lieu APRÈS le verbe (calque français).}
    \textit{Erreur :} *\zh{我踢足球在体育场。} (Wǒ tī zúqiú zài tǐyùchǎng.) — « Je joue foot AU stade ».
    \textit{Correct :} \zh{我\textbf{在体育场}踢足球。} (Wǒ \textbf{zài tǐyùchǎng} tī zúqiú.)
    \textbf{Règle d'or :} \cle{Sujet + Temps + Lieu (+ Manière) + Verbe}.
    \item \textbf{Mots-outils : Confusion \cle{了} (accompli) / \cle{过} (expérience) / \cle{着} (duratif).}
    \textit{Contexte débat :} « J'ai déjà joué au foot » (expérience de vie) $\rightarrow$ \zh{我踢\textbf{过}足球。}
    *\zh{我踢了足球} = « J'ai joué au foot (action terminée, focus sur la fin) ».
    *\zh{我踢着足球} = « Je suis en train de jouer au foot » (rare, \zh{正在踢} préféré).
    \item \textbf{Classificateurs : Oubli ou erreur sur activités sportives.}
    \textit{Erreur :} *\zh{一场游泳} (yī chǎng yóuyǒng — \cle{场} pour match/film/spectacle), *\zh{一个跑步} (yī gè pǎobù — \cle{个} générique mal adapté).
    \textit{Correct :} \zh{一次游泳} (yī cì yóuyǒng — \cle{次} fois/séance), \zh{一场比赛} (yī chǎng bǐsài — match), \zh{一节课} (yī jié kè — séance de cours), \zh{一趟} (yī tàng — aller-retour/parcours).
    \textbf{Astuce :} \cle{次} (cì) = classificateur universel pour « fois / séance » d'activité.
    \item \textbf{Caractères proches / Homophones : \cle{练} (liàn, pratiquer) vs \cle{炼} (liàn, raffiner) ; \cle{锻} (duàn, forger) + \cle{炼} (liàn, affiner) = \zh{锻炼} (exercice physique).}
    \textit{Erreur fréquente :} Écrire *\zh{锻练} au lieu de \zh{锻炼}. Confondre \zh{身} (shēn, corps, 7 traits) et \zh{申} (shēn, étendre, 5 traits) dans \zh{身体} vs *\zh{申体}.
\end{enumerate}
\textbf{Conseil final :} À l'oral, si blocage lexical : stratégie de \cle{paraphrase} : \zh{就是…那个…运动…} (C'est… celui-là… le sport…). Ne restez pas silencieux ! La fluidité prime sur le mot exact.
\end{attention}

\newpage

\coursec{Exercices}

\courssub{Je vérifie mes connaissances}

\begin{exercice}[QCM : le vocabulaire du sport]
Pour chaque question, entoure la bonne réponse.

1. \zh{运动} (yùndòng) signifie :
\begin{itemize}
\item a) la santé
\item b) le sport
\item c) le repos
\end{itemize}

2. Pour dire « jouer au football », on dit :
\begin{itemize}
\item a) \zh{打篮球} (dǎ lánqiú)
\item b) \zh{踢足球} (tī zúqiú)
\item c) \zh{跑步} (pǎobù)
\end{itemize}

3. \zh{身体} (shēntǐ) signifie :
\begin{itemize}
\item a) le corps
\item b) l'esprit
\item c) la maladie
\end{itemize}

4. Quelle phrase signifie « Le sport est très important pour la santé » ?
\begin{itemize}
\item a) \zh{运动对健康很重要。}
\item b) \zh{健康对运动很重要。}
\item c) \zh{运动很健康重要。}
\end{itemize}
\end{exercice}

\begin{exercice}[Vrai ou faux justifié]
Dis si les phrases suivantes sont vraies ou fausses, puis justifie ta réponse en français.

1. \zh{我觉得运动很有意思。} (Wǒ juéde yùndòng hěn yǒu yìsi.) signifie « Je pense que le sport est ennuyeux ».
2. \zh{因为...所以...} (yīnwei... suǒyǐ...) sert à exprimer une cause et sa conséquence.
3. \zh{我不同意。} (Wǒ bù tóngyì.) est une formule pour donner son accord.
4. \zh{你呢？} (Nǐ ne ?) permet de relancer la parole vers son interlocuteur.
\end{exercice}

\begin{exercice}[Vocabulaire : complète en chinois]
Complète chaque phrase avec le mot chinois qui convient, choisi dans la liste : \zh{锻炼} (duànliàn), \zh{游泳} (yóuyǒng), \zh{健康} (jiànkāng), \zh{比赛} (bǐsài), \zh{每天} (měitiān).

1. Le week-end, j'aime nager : 周末我喜欢\_\_\_\_\_\_。
2. Le sport est bon pour la santé : 运动对\_\_\_\_\_\_有好处。
3. Je fais de l'exercice tous les jours : 我\_\_\_\_\_\_锻炼。
4. Samedi, il y a un match de football à Yaoundé : 星期六雅温得有足球\_\_\_\_\_\_。
5. Mon grand-père fait de la gymnastique le matin : 我爷爷早上\_\_\_\_\_\_。
\end{exercice}

\begin{exercice}[Associe le pinyin aux caractères]
Relie chaque mot en caractères à son pinyin.

\begin{center}
\begin{tabularx}{\linewidth}{|X|X|}
\hline
\rowcolor{popblueL} Caractères & Pinyin \\
\hline
1. 运动 & a) jiànkāng \\
\hline
2. 跑步 & b) tóngyì \\
\hline
3. 健康 & c) yùndòng \\
\hline
4. 同意 & d) pǎobù \\
\hline
5. 重要 & e) zhòngyào \\
\hline
\end{tabularx}
\end{center}
\end{exercice}

\courssub{Je m'entraîne}

\begin{exercice}[Traduis en chinois]
Traduis les phrases suivantes en chinois (caractères + pinyin).

1. Je pense que le sport est très important.
2. Parce que je veux être en bonne santé, je cours tous les matins.
3. Je ne suis pas d'accord : faire du sport tous les jours, c'est trop fatigant.
4. Tu as raison, alors nous pouvons faire du sport ensemble trois fois par semaine.
5. Pour moi, le football est le sport le plus intéressant.
\end{exercice}

\begin{exercice}[Dialogue à trous : les formules fonctionnelles]
Complète le dialogue avec les formules suivantes : \zh{你呢？} (Nǐ ne ?), \zh{我不同意} (Wǒ bù tóngyì), \zh{你说得对} (Nǐ shuō de duì), \zh{那我们可以} (Nà wǒmen kěyǐ). Puis joue le dialogue en binôme.

\begin{textebox}[Dialogue à compléter]
A：我觉得我们应该每天运动。你觉得怎么样？\\
Wǒ juéde wǒmen yīnggāi měitiān yùndòng. Nǐ juéde zěnmeyàng ?\\
B：\_\_\_\_\_\_！每天运动太累了。\\
A：可是运动对健康很重要。\\
B：嗯，\_\_\_\_\_\_。\_\_\_\_\_\_每个星期运动三次。\\
A：好主意！\_\_\_\_\_\_ 你最喜欢什么运动？\\
B：我最喜欢踢足球。
\end{textebox}
\end{exercice}

\begin{exercice}[Discrimination des tons]
Le professeur lit un mot de chaque paire. Entoure le mot entendu, puis lis correctement les deux mots à voix haute.

\begin{center}
\begin{tabularx}{\linewidth}{|X|X|}
\hline
\rowcolor{popblueL} Paire 1 & Paire 2 \\
\hline
pǎo 跑 (courir) / pào 炮 (canon) & liàn 练 (s'entraîner) / lián 连 (relier) \\
\hline
shēn 身 (corps) / shén 神 (esprit) & qiú 球 (ballon) / qiū 秋 (automne) \\
\hline
hǎo 好 (bon) / hào 号 (numéro) & mǎi 买 (acheter) / mài 卖 (vendre) \\
\hline
\end{tabularx}
\end{center}
\end{exercice}

\begin{exercice}[Remets le mini-débat dans l'ordre]
Remets dans l'ordre les répliques de ce débat sur le thème « Faut-il faire du sport tous les jours ? », puis joue-le sans support avec un camarade.

\begin{itemize}
\item a) \zh{那我们可以每个星期运动三四次，好吗？} (Nà wǒmen kěyǐ měi ge xīngqī yùndòng sān sì cì, hǎo ma ?)
\item b) \zh{我觉得我们应该每天运动，因为运动对健康很重要。} (Wǒ juéde wǒmen yīnggāi měitiān yùndòng, yīnwei yùndòng duì jiànkāng hěn zhòngyào.)
\item c) \zh{好！你说得对。从明天开始吧！} (Hǎo ! Nǐ shuō de duì. Cóng míngtiān kāishǐ ba !)
\item d) \zh{我不同意。每天运动太累了，我们也需要休息。} (Wǒ bù tóngyì. Měitiān yùndòng tài lèi le, wǒmen yě xūyào xiūxi.)
\end{itemize}
\end{exercice}

\courssub{Je me prépare au Bac}

\begin{exercice}[Type Bac — Compréhension de l'écrit et questions de langue]
Lis le texte suivant, puis réponds aux questions.

\begin{textebox}[Le sport au lycée]
我叫李雷，是雅温得一所高中的学生。我们学校的学生都很喜欢运动。每天下午四点以后，很多学生在操场上跑步、踢足球或者打篮球。\\
Wǒ jiào Lǐ Léi, shì Yǎwēndé yì suǒ gāozhōng de xuésheng. Wǒmen xuéxiào de xuésheng dōu hěn xǐhuan yùndòng. Měitiān xiàwǔ sì diǎn yǐhòu, hěn duō xuésheng zài cāochǎng shang pǎobù, tī zúqiú huòzhě dǎ lánqiú.\\
我觉得运动对我们非常重要。因为运动可以让我们的身体更健康，也可以让我们学习得更好。所以，我每个星期都锻炼三次。\\
Wǒ juéde yùndòng duì wǒmen fēicháng zhòngyào. Yīnwei yùndòng kěyǐ ràng wǒmen de shēntǐ gèng jiànkāng, yě kěyǐ ràng wǒmen xuéxí de gèng hǎo. Suǒyǐ, wǒ měi ge xīngqī dōu duànliàn sān cì.\\
你呢？你喜欢运动吗？\\
Nǐ ne ? Nǐ xǐhuan yùndòng ma ?
\end{textebox}

\textbf{Questions de compréhension} (réponds en français) :
\begin{enumerate}
\item Où étudie Li Lei ?
\item Que font les élèves après quatre heures de l'après-midi ?
\item Pourquoi Li Lei pense-t-il que le sport est important ? Donne deux raisons.
\item Combien de fois par semaine Li Lei s'entraîne-t-il ?
\end{enumerate}

\textbf{Questions de langue} :
\begin{enumerate}
\item Relève dans le texte une phrase avec la structure \zh{因为...所以...} et traduis-la.
\item Explique le rôle du mot-outil \zh{得} dans la phrase \zh{让我们学习得更好}.
\item Trouve dans le texte deux classificateurs et donne le nom qu'ils accompagnent.
\end{enumerate}
\end{exercice}

\begin{exercice}[Type Bac — Thème et version]
\textbf{Thème.} Traduis en chinois (caractères + pinyin) :
\begin{enumerate}
\item Le sport est très important pour notre santé.
\item Parce que je veux avoir un corps fort, je fais du sport tous les jours.
\item Je ne suis pas d'accord : nous avons aussi besoin de repos.
\end{enumerate}

\textbf{Version.} Traduis en français :

\begin{textebox}[Version]
对我来说，运动很有意思。我每个周末都和朋友一起踢足球。运动以后，我虽然很累，但是很高兴。\\
Duì wǒ láishuō, yùndòng hěn yǒu yìsi. Wǒ měi ge zhōumò dōu hé péngyou yìqǐ tī zúqiú. Yùndòng yǐhòu, wǒ suīrán hěn lèi, dànshì hěn gāoxìng.
\end{textebox}
\end{exercice}

\begin{exercice}[Type Bac — Expression écrite de préparation à l'oral]
\textbf{Sujet :} Tu dois présenter à la classe un exposé oral sur le thème : \zh{运动对健康很重要} (Yùndòng duì jiànkāng hěn zhòngyào — « Le sport est très important pour la santé »).

\textbf{Consignes :}
\begin{enumerate}
\item Rédige par écrit un exposé de \textbf{6 phrases minimum} (environ 60 à 80 caractères chinois).
\item Utilise obligatoirement les structures suivantes : \zh{我觉得...} (je pense que...), \zh{因为...所以...} (parce que... donc...), \zh{对我来说...} (pour moi...) et une conclusion avec \zh{所以，我认为...} (donc, je considère que...).
\item Soigne les tons : ton professeur écoutera ta prononciation.
\item Présente ensuite ton exposé à la classe, sans lire intégralement ton texte.
\end{enumerate}

\textbf{Bonus (jeu de rôle noté) :} En binôme, prépare une interview d'un sportif camerounais célèbre (par exemple Samuel Eto'o) avec \textbf{au moins 4 questions et 4 réponses} en chinois. Utilise les formules de relance : \zh{你呢？}, \zh{为什么？}, \zh{你觉得怎么样？}
\end{exercice}

\newpage

\coursec{Corrigés des exercices}

\begin{corrige}[Exercice 1 — QCM : le vocabulaire du sport]
1. \textbf{réponse b)} : \zh{运动} (yùndòng) signifie « le sport ». « La santé » se dit \zh{健康} (jiànkāng) et « le repos » \zh{休息} (xiūxi).

2. \textbf{réponse b)} : « jouer au football » se dit \zh{踢足球} (tī zúqiú), car le football se joue avec le pied, d'où le verbe \zh{踢} (tī, donner un coup de pied). \zh{打篮球} (dǎ lánqiú) signifie « jouer au basket-ball » (le ballon se frappe avec les mains, d'où \zh{打}) et \zh{跑步} (pǎobù) signifie « courir ».

3. \textbf{réponse a)} : \zh{身体} (shēntǐ) signifie « le corps ». « L'esprit » se dit \zh{精神} (jīngshén) et « la maladie » \zh{病} (bìng).

4. \textbf{réponse a)} : \zh{运动对健康很重要。} (Yùndòng duì jiànkāng hěn zhòngyào.) — « Le sport est très important pour la santé ». La structure est \cle{A 对 B 很 + adjectif} : c'est A qui agit sur B. La phrase b) inverse les rôles (« la santé est importante pour le sport ») et la phrase c) est agrammaticale : \zh{健康} ne peut pas modifier directement \zh{重要}.
\end{corrige}

\begin{corrige}[Exercice 2 — Vrai ou faux justifié]
1. \textbf{Faux.} \zh{我觉得运动很有意思。} (Wǒ juéde yùndòng hěn yǒu yìsi.) signifie « Je pense que le sport est très intéressant ». \zh{有意思} (yǒu yìsi) signifie « intéressant, amusant » ; « ennuyeux » se dirait \zh{没意思} (méi yìsi). Attention à la négation : \zh{没} et non \zh{不} devant \zh{有}.

2. \textbf{Vrai.} \zh{因为...所以...} (yīnwei... suǒyǐ...) exprime bien une cause et sa conséquence : « parce que... donc... ». Exemple : \zh{因为运动对健康有好处，所以我每天锻炼。} (Yīnwei yùndòng duì jiànkāng yǒu hǎochù, suǒyǐ wǒ měitiān duànliàn.) — « Parce que le sport est bon pour la santé, je m'entraîne tous les jours. »

3. \textbf{Faux.} \zh{我不同意。} (Wǒ bù tóngyì.) exprime le \emph{désaccord} : « je ne suis pas d'accord ». Pour donner son accord, on dit \zh{我同意。} (Wǒ tóngyì.) ou \zh{你说得对。} (Nǐ shuō de duì. — « Tu as raison »).

4. \textbf{Vrai.} \zh{你呢？} (Nǐ ne ?) est une formule de relance : « et toi ? ». La particule \zh{呢} (ne) reprend le sujet de la question précédente et la retourne vers l'interlocuteur. C'est un outil essentiel pour maintenir un dialogue oral.
\end{corrige}

\begin{corrige}[Exercice 3 — Vocabulaire : complète en chinois]
1. 周末我喜欢\zh{游泳}。 (Zhōumò wǒ xǐhuan yóuyǒng.) — « Le week-end, j'aime nager. »

2. 运动对\zh{健康}有好处。 (Yùndòng duì jiànkāng yǒu hǎochù.) — « Le sport est bon pour la santé. » Structure : \cle{A 对 B 有好处}.

3. 我\zh{每天}锻炼。 (Wǒ měitiān duànliàn.) — « Je fais de l'exercice tous les jours. » Le complément de temps \zh{每天} se place \emph{avant} le verbe.

4. 星期六雅温得有足球\zh{比赛}。 (Xīngqīliù Yǎwēndé yǒu zúqiú bǐsài.) — « Samedi, il y a un match de football à Yaoundé. »

5. 我爷爷早上\zh{锻炼}。 (Wǒ yéye zǎoshang duànliàn.) — « Mon grand-père fait de la gymnastique le matin. » \zh{锻炼} (duànliàn) signifie « s'entraîner, faire de l'exercice » ; c'est le seul verbe possible ici, \zh{游泳} exigeant un contexte de piscine.

\textbf{Point de vigilance :} en chinois, le complément de temps (\zh{每天}, \zh{早上}, \zh{周末}) se place toujours \emph{avant} le verbe, contrairement au français où il se met souvent à la fin.
\end{corrige}

\begin{corrige}[Exercice 4 — Associe le pinyin aux caractères]
1. \zh{运动} — \textbf{c)} yùndòng (le sport)

2. \zh{跑步} — \textbf{d)} pǎobù (courir)

3. \zh{健康} — \textbf{a)} jiànkāng (la santé)

4. \zh{同意} — \textbf{b)} tóngyì (être d'accord)

5. \zh{重要} — \textbf{e)} zhòngyào (important)

\textbf{Point de vigilance sur les tons :} \zh{重} se lit ici zhòng (4\ieme{} ton) dans \zh{重要} ; dans \zh{重新} (chóngxīn, « de nouveau »), il se lit chóng (2\ieme{} ton). De même, ne confonds pas tóngyì (être d'accord) avec tǒngyī (统一, unifier) : seuls les tons les distinguent.
\end{corrige}

\begin{corrige}[Exercice 5 — Traduis en chinois]
1. 我觉得运动很重要。 (Wǒ juéde yùndòng hěn zhòngyào.) — Structure : \cle{我觉得 + phrase}. Pas de \zh{是} devant l'adjectif \zh{重要} : en chinois, l'adjectif fait office de verbe.

2. 因为我想身体健康，所以我每天早上跑步。 (Yīnwei wǒ xiǎng shēntǐ jiànkāng, suǒyǐ wǒ měitiān zǎoshang pǎobù.) — Variante acceptée : \zh{因为我想保持健康，...} (Yīnwei wǒ xiǎng bǎochí jiànkāng, ...) — « parce que je veux rester en bonne santé ».

3. 我不同意：每天运动太累了。 (Wǒ bù tóngyì : měitiān yùndòng tài lèi le.) — Structure : \cle{太 + adjectif + 了} pour exprimer « trop... ». On peut ajouter \zh{我们也需要休息} (wǒmen yě xūyào xiūxi — « nous avons aussi besoin de repos »).

4. 你说得对，那我们可以每个星期一起运动三次。 (Nǐ shuō de duì, nà wǒmen kěyǐ měi ge xīngqī yìqǐ yùndòng sān cì.) — « Trois fois par semaine » : \zh{每个星期...三次} ; le complément de fréquence \zh{三次} se place après le verbe.

5. 对我来说，足球是最有意思的运动。 (Duì wǒ láishuō, zúqiú shì zuì yǒu yìsi de yùndòng.) — Superlatif : \cle{最 + adjectif + 的 + nom}. Ici \zh{是} est nécessaire car on identifie une chose à une autre (nom = nom).
\end{corrige}

\begin{corrige}[Exercice 6 — Dialogue à trous : les formules fonctionnelles]
Dialogue complété :

A：我觉得我们应该每天运动。你觉得怎么样？\\
B：\zh{我不同意}！每天运动太累了。\\
A：可是运动对健康很重要。\\
B：嗯，\zh{你说得对}。\zh{那我们可以}每个星期运动三次。\\
A：好主意！\zh{你呢？} 你最喜欢什么运动？\\
B：我最喜欢踢足球。

\textbf{Traduction :} « A : Je pense que nous devrions faire du sport tous les jours. Qu'en penses-tu ? — B : Je ne suis pas d'accord ! Faire du sport tous les jours, c'est trop fatigant. — A : Mais le sport est très important pour la santé. — B : Hmm, tu as raison. Alors nous pouvons faire du sport trois fois par semaine. — A : Bonne idée ! Et toi, quel sport préfères-tu ? — B : Je préfère jouer au football. »

\textbf{Justification :} la première réplique de B contient « 太累了 » (trop fatigant), ce qui suppose un désaccord : \zh{我不同意}. Ensuite B reconnaît le tort : \zh{你说得对}, puis propose une solution avec \zh{那我们可以}. Enfin A relance la parole avec \zh{你呢？} avant de préciser sa question.

\textbf{Point de vigilance :} dans \zh{你说得对}, le mot-outil \zh{得} introduit un complément de manière : « tu parles \emph{de manière} juste ». Ne pas écrire \zh{你说的对}.
\end{corrige}

\begin{corrige}[Exercice 7 — Discrimination des tons]
Cet exercice est oral : la correction porte sur la prononciation attendue.

\textbf{Paire 1 :} pǎo 跑 (3\ieme{} ton, le ton descend puis remonte : courir) / pào 炮 (4\ieme{} ton, le ton tombe franchement : canon). \textbf{Paire 2 :} liàn 练 (4\ieme{} ton : s'entraîner) / lián 连 (2\ieme{} ton, le ton monte : relier).

\textbf{Paire 3 :} shēn 身 (1\ier{} ton, haut et plat : corps) / shén 神 (2\ieme{} ton, montant : esprit). \textbf{Paire 4 :} qiú 球 (2\ieme{} ton : ballon) / qiū 秋 (1\ier{} ton : automne).

\textbf{Paire 5 :} hǎo 好 (3\ieme{} ton : bon) / hào 号 (4\ieme{} ton : numéro). \textbf{Paire 6 :} mǎi 买 (3\ieme{} ton : acheter) / mài 卖 (4\ieme{} ton : vendre).

\textbf{Conseil de prononciation :} pour le 3\ieme{} ton à l'intérieur d'une phrase, prononce seulement la partie descendante (demi-3\ieme{} ton) ; le ton complet ne s'entend qu'en fin de phrase ou à l'isolement. Devant un autre 3\ieme{} ton, le premier devient un 2\ieme{} ton : \zh{你好} se prononce en réalité ní hǎo. Pour mǎi/mài, retiens le moyen mnémotechnique : « j'achète en hésitant (3\ieme{} ton), je vends d'un ton décidé (4\ieme{} ton) ».
\end{corrige}

\begin{corrige}[Exercice 8 — Remets le mini-débat dans l'ordre]
Ordre correct : \textbf{b — d — a — c}.

1. \textbf{b)} \zh{我觉得我们应该每天运动，因为运动对健康很重要。} (Wǒ juéde wǒmen yīnggāi měitiān yùndòng, yīnwei yùndòng duì jiànkāng hěn zhòngyào.) — « Je pense que nous devrions faire du sport tous les jours, parce que le sport est très important pour la santé. » (ouverture du débat : opinion + justification)

2. \textbf{d)} \zh{我不同意。每天运动太累了，我们也需要休息。} (Wǒ bù tóngyì. Měitiān yùndòng tài lèi le, wǒmen yě xūyào xiūxi.) — « Je ne suis pas d'accord. Faire du sport tous les jours, c'est trop fatigant, nous avons aussi besoin de repos. » (désaccord argumenté)

3. \textbf{a)} \zh{那我们可以每个星期运动三四次，好吗？} (Nà wǒmen kěyǐ měi ge xīngqī yùndòng sān sì cì, hǎo ma ?) — « Alors nous pouvons faire du sport trois ou quatre fois par semaine, d'accord ? » (proposition de compromis ; \zh{三四次} = « trois ou quatre fois », sans \zh{或} entre deux chiffres voisins)

4. \textbf{c)} \zh{好！你说得对。从明天开始吧！} (Hǎo ! Nǐ shuō de duì. Cóng míngtiān kāishǐ ba !) — « D'accord ! Tu as raison. Commençons dès demain ! » (accord final ; la particule \zh{吧} adoucit la suggestion)

\textbf{Logique du débat :} opinion → désaccord → compromis → accord. C'est la structure type à réutiliser dans vos propres débats.
\end{corrige}

\begin{corrige}[Exercice 9 — Type Bac — Compréhension de l'écrit et questions de langue]
\textbf{Barème indicatif : 10 points} (compréhension 5 pts ; langue 5 pts).

\textbf{Questions de compréhension} (1 pt + 1 pt + 2 pts + 1 pt) :

1. Li Lei étudie dans un lycée (une école secondaire) à Yaoundé : \zh{雅温得一所高中}.

2. Après quatre heures de l'après-midi, beaucoup d'élèves courent, jouent au football ou au basket-ball sur le terrain de sport : \zh{很多学生在操场上跑步、踢足球或者打篮球}.

3. Il pense que le sport est important pour deux raisons : (a) le sport peut rendre notre corps plus sain (\zh{运动可以让我们的身体更健康}) ; (b) il peut aussi nous permettre de mieux étudier (\zh{也可以让我们学习得更好}).

4. Il s'entraîne trois fois par semaine : \zh{我每个星期都锻炼三次}.

\textbf{Questions de langue} (2 pts + 2 pts + 1 pt) :

1. Phrase avec \zh{因为...所以...} : \zh{因为运动可以让我们的身体更健康，也可以让我们学习得更好。所以，我每个星期都锻炼三次。} — Traduction : « Parce que le sport peut rendre notre corps plus sain et nous permettre de mieux étudier, je m'entraîne trois fois par semaine. » (On accepte aussi le découpage en deux phrases, fréquent en chinois : \zh{因为...} puis \zh{所以...} au début de la phrase suivante.)

2. Dans \zh{让我们学习得更好}, le mot-outil \zh{得} (de) introduit un \cle{complément de manière ou de degré} après le verbe \zh{学习} : il indique \emph{comment} on étudie, ici « mieux » (\zh{更好}). Structure : \cle{verbe + 得 + adjectif}. Attention à ne pas confondre avec \zh{的} (de, lien de possession ou de qualification avant un nom) ni avec \zh{地} (de, devant un verbe pour un adverbe).

3. Deux classificateurs : \zh{所} (suǒ) dans \zh{一所高中} (« un lycée » — classificateur des établissements) et \zh{次} (cì) dans \zh{锻炼三次} (« s'entraîner trois fois » — classificateur des occurrences, des fois). On accepte aussi \zh{个} (ge) dans \zh{每个星期}.

\textbf{Points de vigilance :} \zh{或者} (huòzhě, « ou ») s'emploie dans les phrases affirmatives ; dans une question, on utilise \zh{还是} (háishi). Ne traduis pas \zh{让} (ràng) par « laisser » ici : c'est « faire que, permettre ».
\end{corrige}

\begin{corrige}[Exercice 10 — Type Bac — Thème et version]
\textbf{Barème indicatif : 10 points} (thème 5 pts ; version 5 pts).

\textbf{Thème} (environ 1,5 pt par phrase) :

1. 运动对我们的健康很重要。 (Yùndòng duì wǒmen de jiànkāng hěn zhòngyào.) — Structure \cle{A 对 B 很重要} ; ne pas oublier \zh{的} entre \zh{我们} et \zh{健康}.

2. 因为我想有一个强壮的身体，所以我每天运动。 (Yīnwei wǒ xiǎng yǒu yí ge qiángzhuàng de shēntǐ, suǒyǐ wǒ měitiān yùndòng.) — Variante acceptée : \zh{因为我想让身体更强壮，...} (Yīnwei wǒ xiǎng ràng shēntǐ gèng qiángzhuàng, ...). « Un corps fort » : \zh{强壮的身体} ; le classificateur de \zh{身体} est \zh{个}.

3. 我不同意：我们也需要休息。 (Wǒ bù tóngyì : wǒmen yě xūyào xiūxi.) — « Avoir besoin de » : \zh{需要 + nom} ; « repos » : \zh{休息} (xiūxi).

\textbf{Version} (5 pts) :

« Pour moi, le sport est très intéressant. Chaque week-end, je joue au football avec mes amis. Après le sport, bien que je sois très fatigué, je suis très content. »

Détail grammatical : \zh{对我来说} (duì wǒ láishuō) = « pour moi, en ce qui me concerne » ; \zh{虽然...但是...} (suīrán... dànshì...) = « bien que... mais... » — en chinois, les deux termes se combinent dans la même phrase, contrairement au français où « bien que » suffit ; \zh{和...一起} (hé... yìqǐ) = « ensemble avec... ».

\textbf{Points de vigilance :} ne traduis pas \zh{有意思} par « avoir un sens » : c'est « intéressant, amusant ». \zh{运动以后} = « après le sport » (et non « après, le sport »).
\end{corrige}

\begin{corrige}[Exercice 11 — Type Bac — Expression écrite de préparation à l'oral]
\textbf{Barème indicatif : 10 points} — respect des 6 phrases minimum et des 4 structures imposées (4 pts) ; correction grammaticale et vocabulaire (3 pts) ; richesse et cohérence du propos (2 pts) ; prononciation et tons à l'oral (1 pt). Bonus jeu de rôle : 2 pts (4 questions et 4 réponses correctes, formules de relance utilisées).

\textbf{Proposition de rédaction modèle :}

\begin{textebox}[Exposé modèle : 运动对健康很重要]
大家好！今天我说一说运动和健康。\\
Dàjiā hǎo ! Jīntiān wǒ shuō yi shuō yùndòng hé jiànkāng.\\
对我来说，运动是生活中很重要的一部分。\\
Duì wǒ láishuō, yùndòng shì shēnghuó zhōng hěn zhòngyào de yí bùfen.\\
我觉得运动可以让我们的身体更健康，也可以让我们的心情更好。\\
Wǒ juéde yùndòng kěyǐ ràng wǒmen de shēntǐ gèng jiànkāng, yě kěyǐ ràng wǒmen de xīnqíng gèng hǎo.\\
因为运动对身体有好处，所以我每个星期锻炼三次。\\
Yīnwei yùndòng duì shēntǐ yǒu hǎochù, suǒyǐ wǒ měi ge xīngqī duànliàn sān cì.\\
有时候我跑步，有时候我和朋友一起踢足球。\\
Yǒu shíhou wǒ pǎobù, yǒu shíhou wǒ hé péngyou yìqǐ tī zúqiú.\\
所以，我认为每个人都应该多运动。谢谢大家！\\
Suǒyǐ, wǒ rènwéi měi ge rén dōu yīnggāi duō yùndòng. Xièxie dàjiā !
\end{textebox}

\textbf{Traduction :} « Bonjour à tous ! Aujourd'hui, je vais vous parler du sport et de la santé. Pour moi, le sport est une partie très importante de la vie. Je pense que le sport peut rendre notre corps plus sain et aussi améliorer notre humeur. Parce que le sport est bon pour le corps, je m'entraîne trois fois par semaine. Parfois je cours, parfois je joue au football avec mes amis. Donc, je considère que tout le monde devrait faire plus de sport. Merci à tous ! »

\textbf{Proposition de bonus (interview de Samuel Eto'o) :}

\begin{textebox}[Interview modèle]
A：埃托奥先生，您好！您最喜欢什么运动？\\
Àituō'ào xiānsheng, nín hǎo ! Nín zuì xǐhuan shénme yùndòng ?\\
B：你好！我当然最喜欢足球。\\
Nǐ hǎo ! Wǒ dāngrán zuì xǐhuan zúqiú.\\
A：您为什么喜欢足球？\\
Nín wèishénme xǐhuan zúqiú ?\\
B：因为足球很有意思，也可以让很多人快乐。\\
Yīnwei zúqiú hěn yǒu yìsi, yě kěyǐ ràng hěn duō rén kuàilè.\\
A：您觉得喀麦隆的年轻人应该每天运动吗？\\
Nín juéde Kāmàilóng de niánqīngrén yīnggāi měitiān yùndòng ma ?\\
B：我觉得他们应该经常锻炼，但是也要好好休息。你呢？你喜欢运动吗？\\
Wǒ juéde tāmen yīnggāi jīngcháng duànliàn, dànshì yě yào hǎohao xiūxi. Nǐ ne ? Nǐ xǐhuan yùndòng ma ?\\
A：我也喜欢！谢谢您！\\
Wǒ yě xǐhuan ! Xièxie nín !
\end{textebox}

\textbf{Points de vigilance :}
\begin{itemize}
\item Pas de \zh{是} devant un adjectif : écris \zh{运动很重要}, jamais \zh{运动是很重要}.
\item \zh{因为} et \zh{所以} peuvent se combiner dans la même phrase en chinois, contrairement au français.
\item Le complément de fréquence \zh{三次} se place après le verbe : \zh{锻炼三次}.
\item Utilise \zh{您} (nín, « vous » poli) pour interviewer une personnalité, et non \zh{你}.
\item À l'oral, soigne particulièrement les 4\ieme{} tons de \zh{运动} (yùndòng) et le 3\ieme{} ton de \zh{所以} (suǒyǐ).
\end{itemize}
\end{corrige}

\newpage

\coursec{Fiche bilan}

\courssub{Fiche bilan — Leçon 11 : Expression orale — Importance de la pratique du sport}

\begin{popfigure}
\begin{tikzpicture}[
  mindmap,
  grow cyclic,
  every node/.style={concept, execute at begin node=\hskip0pt},
  concept color=popblue,
  level 1/.style={level distance=4.5cm, sibling angle=360/7},
  level 2/.style={level distance=3cm, sibling angle=360/4},
  branch/.style={concept color=#1, edge from parent/.style={draw=#1, thick}}
]
\node[concept, text width=2.8cm, font=\bfseries\large, minimum size=1.8cm] {Leçon 11\\Sport \& Santé}
  child[branch=poporange] { node[concept, text width=2.2cm] {Vocabulaire\\Sports \& Santé} }
  child[branch=popgreen] { node[concept, text width=2.2cm] {Dialogue 1\\Mes sports} }
  child[branch=poppurple] { node[concept, text width=2.2cm] {Formules\\Fonctionnelles} }
  child[branch=poppink] { node[concept, text width=2.2cm] {Tons \&\\Prononciation} }
  child[branch=popgold] { node[concept, text width=2.2cm] {Dialogue 2\\Mini-débat} }
  child[branch=popblue] { node[concept, text width=2.2cm] {Jeux de Rôle\\Débat guidé} }
  child[branch=popdark] { node[concept, text width=2.2cm] {Socio-culturel\\Cameroun / Chine} };
\end{tikzpicture}
\legende{Carte mentale de la leçon 11 : Expression orale — Importance du sport}
\end{popfigure}

\begin{bilan}

\textbf{Lexique clé (chinois / français)} \\[4pt]
\begin{center}
\begin{tabularx}{\linewidth}{|X|X|X|X|}
\hline
\rowcolor{popblueL} \textbf{Caractères} & \textbf{Pinyin} & \textbf{Nature} & \textbf{Traduction} \\ \hline
运动 & yùndòng & nom/verbe & sport ; faire du sport \\ \hline
锻炼 & duànliàn & verbe & s'exercer, se renforcer (corps) \\ \hline
健康 & jiànkāng & adj/nom & santé ; en bonne santé \\ \hline
身体 & shēntǐ & nom & corps, condition physique \\ \hline
坚持 & jiānchí & verbe & persévérer, tenir bon \\ \hline
比赛 & bǐsài & nom/verbe & match, compétition ; concourir \\ \hline
队友 & duìyǒu & nom & coéquipier \\ \hline
规则 & guīzé & nom & règle (de jeu) \\ \hline
胜负 & shèngfù & nom & victoire et défaite, issue \\ \hline
习惯 & xíguàn & nom/verbe & habitude ; avoir l'habitude \\ \hline
每天 & měitiān & adv & chaque jour \\ \hline
至少 & zhìshǎo & adv & au moins \\ \hline
应该 & yīnggāi & verbe modal & devoir (conseil moral) \\ \hline
虽然…但是… & suīrán… dànshì… & conj & bien que… cependant… \\ \hline
不仅…而且… & bùjǐn… érqiě… & conj & non seulement… mais encore… \\ \hline
\end{tabularx}
\end{center}

\textbf{Structures grammaticales \& Formules fonctionnelles à connaître par cœur} \\[4pt]
\begin{center}
\begin{tabularx}{\linewidth}{|X|X|X|}
\hline
\rowcolor{poporangeL} \textbf{Fonction / Structure} & \textbf{Exemple (caractères + pinyin)} & \textbf{Traduction} \\ \hline
\multicolumn{3}{|l|}{\textbf{1. Présenter un exposé}} \\ \hline
首先…，其次…，最后… & 首先，运动有益健康。 & Premièrement, le sport est bon pour la santé. \\ \hline
我认为 / 我觉得 & 我认为坚持很重要。 & Je pense / je trouve que la persévérance est importante. \\ \hline
\multicolumn{3}{|l|}{\textbf{2. Donner son avis / Argumenter}} \\ \hline
一方面…，另一方面… & 一方面锻炼身体，另一方面交朋友。 & D'un côté renforcer le corps, de l'autre se faire des amis. \\ \hline
因为…所以… & 因为下雨，所以没去踢球。 & Parce qu'il pleut, je n'ai pas joué au foot. \\ \hline
\multicolumn{3}{|l|}{\textbf{3. Exprimer l'obligation / le conseil}} \\ \hline
应该 / 应当 + V & 我们应该每天运动半小时。 & Nous devrions faire du sport 30 min par jour. \\ \hline
最好 + V & 你最好先热身。 & Tu ferais mieux de t'échauffer d'abord. \\ \hline
\multicolumn{3}{|l|}{\textbf{4. Relancer / Interagir en débat}} \\ \hline
你觉得呢？ & 这个观点你觉得呢？ & Qu'en penses-tu ? \\ \hline
我赞同 / 我不同意 & 我赞同你的看法。 / 我不同意。 & Je suis d'accord / Je ne suis pas d'accord. \\ \hline
为什么呢？ & 你为什么这么认为？ & Pourquoi penses-tu cela ? \\ \hline
\multicolumn{3}{|l|}{\textbf{5. Concession / Opposition (niveau HSK 3-4)}} \\ \hline
虽然…，但是/可是… & 虽然很累，但是我很开心。 & Bien que fatigué, je suis content. \\ \hline
尽管…，还是… & 尽管下雨，他还是去跑步了。 & Malgré la pluie, il est quand même allé courir. \\ \hline
\end{tabularx}
\end{center}

\textbf{Méthodes express (oral)} \\
\begin{itemize}
  \item \textbf{Préparer un exposé (2 min) :} 1) Accroche (question/chiffre) $\rightarrow$ 2) Annonce du plan (3 points) $\rightarrow$ 3) Développement avec connecteurs (首先/而且) $\rightarrow$ 4) Conclusion + ouverture (j'espère que...).
  \item \textbf{Réussir un débat :} Écouter $\rightarrow$ Noter mots-clés $\rightarrow$ Valider/Contredire avec \zh{我赞同/我不同意} + \zh{因为…所以…} $\rightarrow$ Relancer (\zh{你觉得呢？}).
  \item \textbf{Tons à surveiller :} \zh{运动 yùn-dòng (4-4)}, \zh{健康 jiàn-kāng (4-1)}, \zh{坚持 jiān-chí (1-2)}, \zh{应该 yīng-gāi (1-1)}. Attention au \emph{sandhi} du 3e ton (\zh{锻炼 duàn-liàn $\rightarrow$ duàn-liàn}).
\end{itemize}

\end{bilan}

\begin{aretenir}[Checklist avant le Bac]
$\square$ Je connais le vocabulaire des sports courants (足球, 篮球, 乒乓球, 跑步, 游泳) et des termes santé (健康, 身体, 锻炼). \\
$\square$ Je sais présenter un mini-exposé structuré (3 arguments) avec \zh{首先、其次、最后}. \\
$\square$ J'utilise correctement \zh{应该 / 最好} pour donner un conseil. \\
$\square$ Je maîtrise la structure \zh{虽然…但是…} pour nuancer mon avis. \\
$\square$ Je sais exprimer l'accord/désaccord : \zh{我赞同 / 我不同意 / 我有不同看法}. \\
$\square$ Je relance mon interlocuteur avec \zh{你觉得呢？ / 为什么呢？}. \\
$\square$ Je prononce correctement les tons des mots-clés (notamment les 4e tons et le sandhi du 3e ton). \\
$\square$ Je connais 2 différences culturelles sportives Cameroun/Chine (ex: popularité football vs ping-pong, places publiques vs stades). \\
$\square$ Je sais écrire les caractères clés : \zh{运、动、健、康、坚、持、赛}. \\
$\square$ Je peux traduire un court paragraphe « Pourquoi le sport est important » du français vers le chinois. \\
$\square$ Je gère mon temps de parole (1-2 min exposé, tours de parole équilibrés en débat). \\
$\square$ Je reste courtois : \zh{请说、谢谢、没关系} en interaction orale.
\end{aretenir}

\findecours

\end{document}
